\volpart{第四编\quad 公开作业与中文解答}
\originalchapter{8.511作业1--4：晶格与能带}
\psource{Patrick A. Lee，8.511 Fall 2004，公开Problem Sets 1--4；保留原编号。}
\originalsection{作业1（2004年9月21日截止）}
\begin{exercise}[原题1]\label{p:511-1-1}
证明第一Brillouin区体积与实空间Bravais原胞体积的乘积为$(2\pi)^3$。
\end{exercise}
\begin{solution}
令$A=(a_1,a_2,a_3)$，倒格基矩阵$B=2\pi A^{-T}$，则$|\det A||\det B|=(2\pi)^3$。任意倒格原胞与第一Brillouin区各铺砌倒空间，且每个胞恰含一个格点，故体积相同。原胞的任意整数幺模换基不改行列式绝对值，结论与胞的具体形状无关。
\end{solution}
\begin{exercise}[原题2]
证明使Bravais晶格保持不变的绕任意轴旋转只能是一、二、三、四或六重。
\end{exercise}
\begin{solution}
旋转在晶格基中的矩阵为整数矩阵，在三维的本征值是$1,\e^{\ii\theta},\e^{-\ii\theta}$，所以整数迹$1+2\cos\theta$。$2\cos\theta\in\{-2,-1,0,1,2\}$，得到$\theta=\pi,2\pi/3,\pi/2,\pi/3,0$及相应方向的等价旋转。有限阶基本角给所列阶数；五重要求$2\cos(2\pi/5)$非整数，不能保持平移晶格。
\end{solution}
\begin{exercise}[原题3]
CuO$_2$层中Cu形成间距$a$的正方格，O位于每条Cu--Cu最近邻键中点；各层沿$z$简单堆叠，层距$c$，忽略其他原子。(a)给出Bravais晶格、原矢、原胞与基底。(b)原题第二图中O交替向上、向下移位。将该图的条件写清：在Cu坐标$(na,ma)$的右方O，$u_z=\delta(-1)^{n+m}$；在其上方O，$u_z=-\delta(-1)^{n+m}$。给出畸变后原胞、晶格间距和倒晶格；讨论$\delta\to0$时X射线峰如何变化。上下图样翻转不改变平移晶格。
\end{exercise}
\begin{solution}
(a)取$a_1=a\hat x,a_2=a\hat y,a_3=c\hat z$，基底为Cu$(0,0,0)$、O$(a/2,0,0)$、O$(0,a/2,0)$，体积$a^2c$。平面是方格，三维为简单四方Bravais晶格。

(b)平移$(a,0)$改变所有O的移位符号，不能保持畸变；$(a,a)$与$(-a,a)$保持$n+m$偶性。因此新原矢$A_1=a(\hat x+\hat y),A_2=a(-\hat x+\hat y),A_3=c\hat z$，平面原胞边长$\sqrt2a$，体积$2a^2c$，基底含2Cu、4O。倒矢$B_1=(\pi/a)(\hat x+\hat y),B_2=(\pi/a)(-\hat x+\hat y),B_3=(2\pi/c)\hat z$。新增倒格点对应原来禁戒的超结构峰；位移引起振幅一阶$\propto Q_z\delta$，强度$\propto\delta^2$，在畸变消失时逐渐消失。原有主峰位置不需随$\delta$连续漂移，区别在结构因子。
\end{solution}
\originalsection{作业2（2004年9月28日截止）}
原题1开头只有“Marder第38页2.5(b)”指针，该外书题面未取得，不能从后续小问倒造原题。以下保留公开附加(a)、(b)和原题2的完整数学条件。
\begin{exercise}[原题1(a)]
HCP具有平行于包含$a,c$轴的平面的滑移面。说明可把此平面放在哪里，使沿$c/2$平移再反射保持HCP结构。
\end{exercise}
\begin{solution}
用三角原矢$a_1=a(1,0,0),a_2=a(1/2,\sqrt3/2,0),a_3=c\hat z$，两层基底可取$0$和$r_B=(a/2,\sqrt3a/6,c/2)$。滑移面可取竖直平面经过两层点连线的中点且平面法向沿$r_{B\perp}$，即$\vct r_\perp\cdot r_{B\perp}=|r_{B\perp}|^2/2$。其线性反射$R$把三角晶格映回自身：此方向的法向相对$a_1$为$30^\circ$，属于三角格反射对称。仿射反射把原点映到$r_{B\perp}$，再加$c/2$恰为B层；B层映到下一A层。这样明确给出纯$c/2$滑移所需的平面，不能任取竖直晶面。
\end{solution}
\begin{exercise}[原题1(b)(i)--(iii)，公开文件标注Ashcroft--Mermin第109页题5]
NaCl采用边长$a$的FCC晶格，基底$+$离子在原点、$-$离子在$(a/2)\hat x$。倒格矢$K=(4\pi/a)(\nu_1,\nu_2,\nu_3)$，$\nu_i$全为整数或全为半整数，形状因子为$f_\pm$。(i)求结构因子。(ii)闪锌矿的第二基底在$(a/4)(1,1,1)$，求结构因子。(iii)若闭壳层离子的$f_\pm$只依赖$|K|$，如何根据Bragg峰强度区分两种结构？
\end{exercise}
\begin{solution}
(i)$S_K=f_++f_-\e^{-2\pi\ii\nu_1}$，全整数时$f_++f_-$，全半整数时$f_+-f_-$。若$f_+=f_-$，后者消光，是两类原子变得等价导致额外平移对称。

(ii)$S_K=f_++f_-\e^{-\ii\pi(\nu_1+\nu_2+\nu_3)}$。全整数时按和的偶奇为$f_+\pm f_-$；全半整数时相位为$\pm\ii$，得$f_+\pm\ii f_-$。符号依指数约定，强度不变。(iii)在相同$|K|$或已校正原子形状因子的反射族间比较：NaCl的全半整数类按$|f_+-f_-|^2$，闪锌矿对应$|f_+\pm\ii f_-|^2$，实因子时为$f_+^2+f_-^2$；全整数族还按和偶奇分组。须校正几何、多重性、吸收和热振动，不能凭一个峰强度唯一判定。
\end{solution}
\begin{exercise}[原题2(a)--(b)，公开文件标注Ashcroft--Mermin第171页题3]
FCC的$W=(2\pi/a)(1,1/2,0)$处，四个自由波动量为$k,k-(2\pi/a)(1,1,1),k-(2\pi/a)(1,1,-1),k-(2\pi/a)(2,0,0)$。势反演对称，使$U_G$实，令$U_1=U_{111}=U_{11\bar1},U_2=U_{200}$。(a)写四态能量矩阵并求$W$劈裂。(b)在$U=(2\pi/a)(1,1/4,1/4)$用同法求简并能级劈裂。
\end{exercise}
\begin{solution}
(a)四态对角为$\epsilon_i^0$，非对角依次为
$\left(\begin{smallmatrix}0&U_1&U_1&U_2\\U_1&0&U_2&U_1\\U_1&U_2&0&U_1\\U_2&U_1&U_1&0\end{smallmatrix}\right)$。
在$W$对角相同$\epsilon_W$。反对称组$(1,0,0,-1)$与$(0,1,-1,0)$能量$\epsilon_W-U_2$；剩余组$(1,1,1,1)$与$(1,-1,-1,1)$能量$\epsilon_W+U_2\pm2U_1$。

(b)在$U$只有三波简并：第一、第二和第四，矩阵非对角为$(U_1,U_2,U_1)$。第一与第三分量反对称给$\epsilon_U-U_2$；对称组与中间态的$2\times2$块为$\left(\begin{smallmatrix}U_2&\sqrt2U_1\\\sqrt2U_1&0\end{smallmatrix}\right)$，另两根$\epsilon_U+[U_2\pm\sqrt{U_2^2+8U_1^2}]/2$。其余非简并波的一阶作用不包含在该三态公式中。
\end{solution}
\originalsection{作业3（2004年10月5日截止）}
\begin{exercise}[原题1(a)--(e)]
Al为FCC，立方边长$a=4.05\,\text{\AA}$。(a)构造自由电子最低两带，画$\Gamma\to X\to W\to L\to\Gamma$并标简并，$X=(2\pi/a)(1,0,0),L=(2\pi/a)(1/2,1/2,1/2)$。(b)空芯赝势为$r>R_c$时$-Ze_*^2/r$、$r<R_c$时零，$Z=3,R_c=0.6\,\text{\AA}$；$V(q)=-4\pi e_*^2Z\cos(qR_c)/(q^2\Omega)$，$\Omega=a^3/4$，求$X,W,L$劈裂。(c)$W$四态中$S_1$同时交换1与4、2与3，$S_2$同时交换1与2、3与4，证明两者与$H$对易且态可按$\pm1$分类。(d)讨论沿$XW$、$WL$的这些对称与带连接，允许用计算核对。(e)估计费米能并与题面所指商业教材KKR图比较。
\end{exercise}
\begin{solution}
(a)对每个路径点将$\hbar^2|k+G|^2/(2m)$排序，$G=(2\pi/a)(h,k,l)$且$h,k,l$全偶或全奇，取最低两个（遇简并保留重数）。$\Gamma$最低为非简并0；$X$最低2个波简并；$W$最低4个波简并；$L$最低2个波简并。源码中的数值脚本给这一路径的采样图与能量表，$k$点处的简并先排序再合组。

(b)$X$、$L$分别由$G_{200}$与$G_{111}$散射，劈裂$2|V_{200}|$、$2|V_{111}|$；$W$采用上一题四态根，劈裂相对$\epsilon_W$为$-V_{200}$（二重）及$V_{200}\pm2V_{111}$。使用$\hbar^2/(2m)=3.809982\,\mathrm{eV}\,\text{\AA}^{2}$、$e_*^2=14.39965\,\mathrm{eV}\,\text{\AA}$，数值为$V_{111}=0.18768$ eV、$V_{200}=0.97376$ eV，$X$隙1.94752 eV、$L$隙0.37536 eV；$W$偏移$-0.97376$ eV（二重）、0.59840 eV、1.34912 eV。随卷计算证据也保留这些数，核对单位$q^{-2}/\Omega$为长度倒数。

(c)对称矩阵是置换矩阵，$S_i^2=I$且$S_1S_2=S_2S_1$。$W$处相同对角与非对角$V_1,V_2$均不因交换改变，故对易。四个共同本征矢可取$(1,s_2,s_1s_2,s_1)/2$，$s_i=\pm1$，相应偏移$V_2s_1+V_1s_2(1+s_1)$，重现上一题的二重根。

(d)$XW$上动量$(1,y,0)$，自由对角满足$\epsilon_1=\epsilon_4$、$\epsilon_2=\epsilon_3$，所以$S_1$保留，$S_2$除$W$外不保留。$WL$上$k=(1-s/2,1/2,s/2)$，$\epsilon_1=\epsilon_2$、$\epsilon_3=\epsilon_4$，所以$S_2$保留，$S_1$除$W$外不保留。保留不同$S$本征值的支可交叉，同一对称支一般避免交叉。用四态矩阵沿两段对角化即可把端点根按各自本征值连续连接，不能只按排序序号追踪交叉。

(e)每原胞3电子，$n=3/\Omega$给$k_F=(3\pi^2n)^{1/3}$和$E_F=\hbar^2k_F^2/(2m)$，自由模型估计约$11.65\,\mathrm{eV}$。加入弱赝势后需按整个区的占据数积分重新定$\mu$，单靠高对称线不足以精确求填充。所引用Ashcroft--Mermin第205页及Marder第250页KKR图未作为公开OCW文件取得，比较项保留缺口，不伪造书中曲线。
\end{solution}
\begin{center}\begin{tikzpicture}\begin{axis}[width=11cm,height=5.7cm,xlabel={路径},ylabel={能量/eV},title={Al自由电子折叠带},xtick={0,1,2,3,4},xticklabels={$\Gamma$,$X$,$W$,$L$,$\Gamma$},legend pos=north east]
\addplot[thick,color=structurecolor] coordinates {(0,0) (0.016393443,0.00246441) (0.032786885,0.0098576399) (0.049180328,0.02217969) (0.06557377,0.03943056) (0.081967213,0.061610249) (0.098360656,0.088718759) (0.1147541,0.12075609) (0.13114754,0.15772224) (0.14754098,0.19961721) (0.16393443,0.246441) (0.18032787,0.29819361) (0.19672131,0.35487504) (0.21311475,0.41648529) (0.2295082,0.48302435) (0.24590164,0.55449224) (0.26229508,0.63088895) (0.27868852,0.71221448) (0.29508197,0.79846883) (0.31147541,0.889652) (0.32786885,0.98576399) (0.3442623,1.0868048) (0.36065574,1.1927744) (0.37704918,1.3036729) (0.39344262,1.4195001) (0.40983607,1.5402562) (0.42622951,1.6659411) (0.44262295,1.7965549) (0.45901639,1.9320974) (0.47540984,2.0725688) (0.49180328,2.217969) (0.50819672,2.368298) (0.52459016,2.5235558) (0.54098361,2.6837425) (0.55737705,2.8488579) (0.57377049,3.0189022) (0.59016393,3.1938753) (0.60655738,3.3737773) (0.62295082,3.558608) (0.63934426,3.7483676) (0.6557377,3.943056) (0.67213115,4.1426732) (0.68852459,4.3472192) (0.70491803,4.556694) (0.72131148,4.7710977) (0.73770492,4.9904302) (0.75409836,5.2146915) (0.7704918,5.4438816) (0.78688525,5.6780006) (0.80327869,5.9170483) (0.81967213,6.1610249) (0.83606557,6.4099303) (0.85245902,6.6637646) (0.86885246,6.9225276) (0.8852459,7.1862195) (0.90163934,7.4548402) (0.91803279,7.7283897) (0.93442623,8.006868) (0.95081967,8.2902752) (0.96721311,8.5786111) (0.98360656,8.8718759) (1,9.1700695) (1.0163934,9.1706856) (1.0327869,9.1725339) (1.0491803,9.1756144) (1.0655738,9.1799272) (1.0819672,9.1854721) (1.0983607,9.1922492) (1.1147541,9.2002585) (1.1311475,9.2095001) (1.147541,9.2199738) (1.1639344,9.2316798) (1.1803279,9.2446179) (1.1967213,9.2587883) (1.2131148,9.2741908) (1.2295082,9.2908256) (1.2459016,9.3086926) (1.2622951,9.3277918) (1.2786885,9.3481231) (1.295082,9.3696867) (1.3114754,9.3924825) (1.3278689,9.4165105) (1.3442623,9.4417707) (1.3606557,9.4682631) (1.3770492,9.4959877) (1.3934426,9.5249445) (1.4098361,9.5551336) (1.4262295,9.5865548) (1.442623,9.6192082) (1.4590164,9.6530939) (1.4754098,9.6882117) (1.4918033,9.7245618) (1.5081967,9.762144) (1.5245902,9.8009585) (1.5409836,9.8410051) (1.557377,9.882284) (1.5737705,9.9247951) (1.5901639,9.9685383) (1.6065574,10.013514) (1.6229508,10.059722) (1.6393443,10.107161) (1.6557377,10.155834) (1.6721311,10.205738) (1.6885246,10.256874) (1.704918,10.309243) (1.7213115,10.362844) (1.7377049,10.417677) (1.7540984,10.473742) (1.7704918,10.53104) (1.7868852,10.58957) (1.8032787,10.649332) (1.8196721,10.710326) (1.8360656,10.772552) (1.852459,10.836011) (1.8688525,10.900701) (1.8852459,10.966624) (1.9016393,11.03378) (1.9180328,11.102167) (1.9344262,11.171787) (1.9508197,11.242638) (1.9672131,11.314722) (1.9836066,11.388038) (2,11.462587) (2.0163934,11.31349) (2.0327869,11.166858) (2.0491803,11.02269) (2.0655738,10.880986) (2.0819672,10.741747) (2.0983607,10.604972) (2.1147541,10.470662) (2.1311475,10.338816) (2.147541,10.209434) (2.1639344,10.082517) (2.1803279,9.9580646) (2.1967213,9.8360763) (2.2131148,9.7165524) (2.2295082,9.599493) (2.2459016,9.4848979) (2.2622951,9.3727672) (2.2786885,9.263101) (2.295082,9.1558992) (2.3114754,9.0511617) (2.3278689,8.9488887) (2.3442623,8.8490801) (2.3606557,8.7517359) (2.3770492,8.6568561) (2.3934426,8.5644408) (2.4098361,8.4744898) (2.4262295,8.3870032) (2.442623,8.3019811) (2.4590164,8.2194234) (2.4754098,8.13933) (2.4918033,8.0617011) (2.5081967,7.9865366) (2.5245902,7.9138365) (2.5409836,7.8436008) (2.557377,7.7758296) (2.5737705,7.7105227) (2.5901639,7.6476803) (2.6065574,7.5873022) (2.6229508,7.5293886) (2.6393443,7.4739393) (2.6557377,7.4209545) (2.6721311,7.3704341) (2.6885246,7.3223781) (2.704918,7.2767866) (2.7213115,7.2336594) (2.7377049,7.1929966) (2.7540984,7.1547983) (2.7704918,7.1190643) (2.7868852,7.0857948) (2.8032787,7.0549897) (2.8196721,7.0266489) (2.8360656,7.0007726) (2.852459,6.9773607) (2.8688525,6.9564133) (2.8852459,6.9379302) (2.9016393,6.9219115) (2.9180328,6.9083573) (2.9344262,6.8972674) (2.9508197,6.888642) (2.9672131,6.882481) (2.9836066,6.8787843) (3,6.8775521) (3.0166667,6.6502108) (3.0333333,6.4266904) (3.05,6.2069908) (3.0666667,5.9911121) (3.0833333,5.7790542) (3.1,5.5708172) (3.1166667,5.3664011) (3.1333333,5.1658058) (3.15,4.9690314) (3.1666667,4.7760779) (3.1833333,4.5869452) (3.2,4.4016334) (3.2166667,4.2201424) (3.2333333,4.0424723) (3.25,3.8686231) (3.2666667,3.6985947) (3.2833333,3.5323872) (3.3,3.3700005) (3.3166667,3.2114348) (3.3333333,3.0566898) (3.35,2.9057658) (3.3666667,2.7586626) (3.3833333,2.6153802) (3.4,2.4759188) (3.4166667,2.3402782) (3.4333333,2.2084584) (3.45,2.0804595) (3.4666667,1.9562815) (3.4833333,1.8359243) (3.5,1.719388) (3.5166667,1.6066726) (3.5333333,1.497778) (3.55,1.3927043) (3.5666667,1.2914515) (3.5833333,1.1940195) (3.6,1.1004083) (3.6166667,1.0106181) (3.6333333,0.92464868) (3.65,0.84250014) (3.6666667,0.76417246) (3.6833333,0.68966564) (3.7,0.61897969) (3.7166667,0.5521146) (3.7333333,0.48907037) (3.75,0.42984701) (3.7666667,0.37444451) (3.7833333,0.32286286) (3.8,0.27510209) (3.8166667,0.23116217) (3.8333333,0.19104311) (3.85,0.15474492) (3.8666667,0.12226759) (3.8833333,0.093611126) (3.9,0.068775521) (3.9166667,0.047760779) (3.9333333,0.030566898) (3.95,0.01719388) (3.9666667,0.0076417246) (3.9833333,0.0019104311) (4,0)};
\addlegendentry{第一带}
\addplot[thick,color=main] coordinates {(0,27.510209) (0.016393443,27.212015) (0.032786885,26.91875) (0.049180328,26.630414) (0.06557377,26.347007) (0.081967213,26.068529) (0.098360656,25.794979) (0.1147541,25.526359) (0.13114754,25.262667) (0.14754098,25.003904) (0.16393443,24.750069) (0.18032787,24.501164) (0.19672131,24.257187) (0.21311475,24.01814) (0.2295082,23.784021) (0.24590164,23.554831) (0.26229508,23.330569) (0.27868852,23.111237) (0.29508197,22.896833) (0.31147541,22.687358) (0.32786885,22.482812) (0.3442623,22.283195) (0.36065574,22.088507) (0.37704918,21.898747) (0.39344262,21.713916) (0.40983607,21.534014) (0.42622951,21.359041) (0.44262295,21.188997) (0.45901639,21.023881) (0.47540984,20.863695) (0.49180328,20.708437) (0.50819672,20.407779) (0.52459016,19.961721) (0.54098361,19.520591) (0.55737705,19.084391) (0.57377049,18.653119) (0.59016393,18.226776) (0.60655738,17.805362) (0.62295082,17.388877) (0.63934426,16.97732) (0.6557377,16.570693) (0.67213115,16.168994) (0.68852459,15.772224) (0.70491803,15.380383) (0.72131148,14.99347) (0.73770492,14.611487) (0.75409836,14.234432) (0.7704918,13.862306) (0.78688525,13.495109) (0.80327869,13.132841) (0.81967213,12.775501) (0.83606557,12.423091) (0.85245902,12.075609) (0.86885246,11.733056) (0.8852459,11.395432) (0.90163934,11.062736) (0.91803279,10.73497) (0.93442623,10.412132) (0.95081967,10.094223) (0.96721311,9.7812432) (0.98360656,9.4731919) (1,9.1700695) (1.0163934,9.1706856) (1.0327869,9.1725339) (1.0491803,9.1756144) (1.0655738,9.1799272) (1.0819672,9.1854721) (1.0983607,9.1922492) (1.1147541,9.2002585) (1.1311475,9.2095001) (1.147541,9.2199738) (1.1639344,9.2316798) (1.1803279,9.2446179) (1.1967213,9.2587883) (1.2131148,9.2741908) (1.2295082,9.2908256) (1.2459016,9.3086926) (1.2622951,9.3277918) (1.2786885,9.3481231) (1.295082,9.3696867) (1.3114754,9.3924825) (1.3278689,9.4165105) (1.3442623,9.4417707) (1.3606557,9.4682631) (1.3770492,9.4959877) (1.3934426,9.5249445) (1.4098361,9.5551336) (1.4262295,9.5865548) (1.442623,9.6192082) (1.4590164,9.6530939) (1.4754098,9.6882117) (1.4918033,9.7245618) (1.5081967,9.762144) (1.5245902,9.8009585) (1.5409836,9.8410051) (1.557377,9.882284) (1.5737705,9.9247951) (1.5901639,9.9685383) (1.6065574,10.013514) (1.6229508,10.059722) (1.6393443,10.107161) (1.6557377,10.155834) (1.6721311,10.205738) (1.6885246,10.256874) (1.704918,10.309243) (1.7213115,10.362844) (1.7377049,10.417677) (1.7540984,10.473742) (1.7704918,10.53104) (1.7868852,10.58957) (1.8032787,10.649332) (1.8196721,10.710326) (1.8360656,10.772552) (1.852459,10.836011) (1.8688525,10.900701) (1.8852459,10.966624) (1.9016393,11.03378) (1.9180328,11.102167) (1.9344262,11.171787) (1.9508197,11.242638) (1.9672131,11.314722) (1.9836066,11.388038) (2,11.462587) (2.0163934,11.31349) (2.0327869,11.166858) (2.0491803,11.02269) (2.0655738,10.880986) (2.0819672,10.741747) (2.0983607,10.604972) (2.1147541,10.470662) (2.1311475,10.338816) (2.147541,10.209434) (2.1639344,10.082517) (2.1803279,9.9580646) (2.1967213,9.8360763) (2.2131148,9.7165524) (2.2295082,9.599493) (2.2459016,9.4848979) (2.2622951,9.3727672) (2.2786885,9.263101) (2.295082,9.1558992) (2.3114754,9.0511617) (2.3278689,8.9488887) (2.3442623,8.8490801) (2.3606557,8.7517359) (2.3770492,8.6568561) (2.3934426,8.5644408) (2.4098361,8.4744898) (2.4262295,8.3870032) (2.442623,8.3019811) (2.4590164,8.2194234) (2.4754098,8.13933) (2.4918033,8.0617011) (2.5081967,7.9865366) (2.5245902,7.9138365) (2.5409836,7.8436008) (2.557377,7.7758296) (2.5737705,7.7105227) (2.5901639,7.6476803) (2.6065574,7.5873022) (2.6229508,7.5293886) (2.6393443,7.4739393) (2.6557377,7.4209545) (2.6721311,7.3704341) (2.6885246,7.3223781) (2.704918,7.2767866) (2.7213115,7.2336594) (2.7377049,7.1929966) (2.7540984,7.1547983) (2.7704918,7.1190643) (2.7868852,7.0857948) (2.8032787,7.0549897) (2.8196721,7.0266489) (2.8360656,7.0007726) (2.852459,6.9773607) (2.8688525,6.9564133) (2.8852459,6.9379302) (2.9016393,6.9219115) (2.9180328,6.9083573) (2.9344262,6.8972674) (2.9508197,6.888642) (2.9672131,6.882481) (2.9836066,6.8787843) (3,6.8775521) (3.0166667,7.1087143) (3.0333333,7.3436973) (3.05,7.5825012) (3.0666667,7.825126) (3.0833333,8.0715716) (3.1,8.3218381) (3.1166667,8.5759254) (3.1333333,8.8338336) (3.15,9.0955627) (3.1666667,9.3611126) (3.1833333,9.6304834) (3.2,9.9036751) (3.2166667,10.180688) (3.2333333,10.461521) (3.25,10.746175) (3.2666667,11.03465) (3.2833333,11.326946) (3.3,11.623063) (3.3166667,11.923001) (3.3333333,12.226759) (3.35,12.534339) (3.3666667,12.845739) (3.3833333,13.16096) (3.4,13.480002) (3.4166667,13.802865) (3.4333333,14.129549) (3.45,14.460053) (3.4666667,14.794379) (3.4833333,15.132525) (3.5,15.474492) (3.5166667,15.82028) (3.5333333,16.169889) (3.55,16.523319) (3.5666667,16.88057) (3.5833333,17.241641) (3.6,17.606533) (3.6166667,17.975247) (3.6333333,18.347781) (3.65,18.724136) (3.6666667,19.104311) (3.6833333,19.488308) (3.7,19.876126) (3.7166667,20.267764) (3.7333333,20.663223) (3.75,21.062503) (3.7666667,21.465604) (3.7833333,21.872526) (3.8,22.283269) (3.8166667,22.697832) (3.8333333,23.116217) (3.85,23.538422) (3.8666667,23.964448) (3.8833333,24.394295) (3.9,24.827963) (3.9166667,25.265452) (3.9333333,25.706762) (3.95,26.151892) (3.9666667,26.600843) (3.9833333,27.053615) (4,27.510209)};
\addlegendentry{第二带}
\end{axis}\end{tikzpicture}\end{center}

\originalsection{作业4（2004年10月12日截止）}
\begin{exercise}[原题1(a)--(e)]
La$_2$CuO$_4$的平面模型中Cu在方格点、O在键中点，只保留Cu的$d_{x^2-y^2}$轨道（能量$E_d$）与每个O的沿键$p$轨道（能量$E_p$）。忽略轨道重叠及O--O跃迁，只保留Cu--最近O的幅度$t_{pd}$。(a)求三带结构。(b)取La为$3+$、非平面O为$2-$、Cu的$4s$空，证明保留轨道中每胞一空穴并找费米面。(c)找鞍点与态密度对数奇异。(d)若15\%的La位由Sr$^{2+}$替换，画空穴掺杂后的费米面；原题给Yoshida等PRB 63,220501(2001)为实验参考。(e)Mattheiss PRL58,1028(1987)称$E_p=E_d,t_{pd}=-1.85\,\mathrm{eV}$可拟合其APW；检验模型预言，原文三维$X,S$对应二维$(\pi/a,\pi/a)$，$\Gamma,Z$对应$(0,0)$。
\end{exercise}
\begin{solution}
(a)按轨道奇偶性选规范，$d,p_x,p_y$矩阵为
\begin{equation}
H(k)=\begin{pmatrix}E_d&2\ii t\sin(k_xa/2)&-2\ii t\sin(k_ya/2)\\-2\ii t\sin(k_xa/2)&E_p&0\\2\ii t\sin(k_ya/2)&0&E_p\end{pmatrix}.
\end{equation}
一条非键带$E_p$，另外$E_\pm=(E_d+E_p)/2\pm\sqrt{(E_d-E_p)^2/4+4t^2[s_x^2+s_y^2]}$。保留矩阵元相位保证Hermitian，改轨道规范不改能量。

(b)离子平衡给Cu$2+$即$d^9$，相对于满$d^{10}$一个空穴；两平面O的相关$p$轨道近满。因此三轨道6个自旋态有5电子，下两带满，上带半满。对仅最近Cu--O的模型，$s_x^2+s_y^2=1$是半填充等能线，即$\cos k_xa+\cos k_ya=0$，费米面为菱形。真实材料未掺杂时是Mott绝缘体，题目中的费米面是该单粒子模型的预言。

(c)$(\pi/a,0)$与$(0,\pi/a)$处，$s_x^2+s_y^2=1+[(\delta k_y)^2-(\delta k_x)^2]a^2/4+\cdots$，故$E_+-E_F=\alpha[(\delta k_y)^2-(\delta k_x)^2]$，$\alpha>0$。第\ref{ch:bands}章的$uv$积分给对数态密度。

(d)“15\%的La位”按字面每化学式有2个La，新增$0.30$空穴；若实验用La$_{2-x}$Sr$_x$且$x=0.15$，则是$0.15$空穴，两种口径必须分开。电子填充下降，$E_F$向上带较低能量移动，等能条件$s_x^2+s_y^2<1$，电子占据区环绕$\Gamma$收缩，空穴区扩张。精确等能值由面积满足上带电子数$1-p$确定，不能任意挪一条线冒称15\%面积。

(e)$E_p=E_d=E_0$时$E_\pm=E_0\pm2|t|\sqrt{s_x^2+s_y^2}$，$\Gamma$三带同能，$(\pi,0)$分裂$\pm3.70$ eV，$(\pi,\pi)$分裂$\pm5.233$ eV。这是可检验的二维拟合预言；对题面引用论文的完整三维APW曲线未取得可再分发原件，不能据这几个点声称已核验论文拟合。原题第三方晶体结构照片未复用，以上条件已足够重建本模型。
\end{solution}

\begin{center}\begin{tikzpicture}\begin{axis}[width=11cm,height=5.7cm,xlabel={$k_xa$},ylabel={$k_ya$},title={CuO$_2$模型的费米面},axis equal image,xmin=-3.2,xmax=3.2,ymin=-3.2,ymax=3.2,legend pos=outer north east]
\addplot[thick,color=structurecolor] coordinates {(-3.1415927,0.0039220882) (-3.1363567,0.0065420497) (-3.1311207,0.011182367) (-3.1258847,0.016190229) (-3.1206487,0.02130805) (-3.1154127,0.026472131) (-3.1101767,0.031659844) (-3.1049407,0.036861213) (-3.0997048,0.042071173) (-3.0944688,0.047286885) (-3.0892328,0.052506634) (-3.0839968,0.057729325) (-3.0787608,0.062954227) (-3.0735248,0.06818083) (-3.0682888,0.073408772) (-3.0630528,0.078637786) (-3.0578168,0.08386767) (-3.0525809,0.089098273) (-3.0473449,0.094329473) (-3.0421089,0.099561178) (-3.0368729,0.10479331) (-3.0316369,0.11002581) (-3.0264009,0.11525863) (-3.0211649,0.12049172) (-3.0159289,0.12572506) (-3.010693,0.13095861) (-3.005457,0.13619234) (-3.000221,0.14142625) (-2.994985,0.1466603) (-2.989749,0.15189448) (-2.984513,0.15712879) (-2.979277,0.16236321) (-2.974041,0.16759772) (-2.9688051,0.17283233) (-2.9635691,0.17806701) (-2.9583331,0.18330177) (-2.9530971,0.1885366) (-2.9478611,0.19377149) (-2.9426251,0.19900644) (-2.9373891,0.20424145) (-2.9321531,0.2094765) (-2.9269172,0.2147116) (-2.9216812,0.21994674) (-2.9164452,0.22518192) (-2.9112092,0.23041714) (-2.9059732,0.23565239) (-2.9007372,0.24088768) (-2.8955012,0.24612299) (-2.8902652,0.25135834) (-2.8850293,0.25659371) (-2.8797933,0.2618291) (-2.8745573,0.26706452) (-2.8693213,0.27229996) (-2.8640853,0.27753542) (-2.8588493,0.28277091) (-2.8536133,0.28800641) (-2.8483773,0.29324192) (-2.8431414,0.29847746) (-2.8379054,0.30371301) (-2.8326694,0.30894857) (-2.8274334,0.31418415) (-2.8221974,0.31941975) (-2.8169614,0.32465535) (-2.8117254,0.32989097) (-2.8064894,0.3351266) (-2.8012534,0.34036224) (-2.7960175,0.3455979) (-2.7907815,0.35083356) (-2.7855455,0.35606923) (-2.7803095,0.36130491) (-2.7750735,0.3665406) (-2.7698375,0.3717763) (-2.7646015,0.37701201) (-2.7593655,0.38224773) (-2.7541296,0.38748345) (-2.7488936,0.39271918) (-2.7436576,0.39795492) (-2.7384216,0.40319066) (-2.7331856,0.40842641) (-2.7279496,0.41366217) (-2.7227136,0.41889793) (-2.7174776,0.4241337) (-2.7122417,0.42936947) (-2.7070057,0.43460525) (-2.7017697,0.43984104) (-2.6965337,0.44507682) (-2.6912977,0.45031262) (-2.6860617,0.45554842) (-2.6808257,0.46078422) (-2.6755897,0.46602003) (-2.6703538,0.47125584) (-2.6651178,0.47649166) (-2.6598818,0.48172747) (-2.6546458,0.4869633) (-2.6494098,0.49219913) (-2.6441738,0.49743496) (-2.6389378,0.50267079) (-2.6337018,0.50790663) (-2.6284659,0.51314247) (-2.6232299,0.51837831) (-2.6179939,0.52361416) (-2.6127579,0.52885001) (-2.6075219,0.53408586) (-2.6022859,0.53932172) (-2.5970499,0.54455757) (-2.5918139,0.54979343) (-2.586578,0.5550293) (-2.581342,0.56026516) (-2.576106,0.56550103) (-2.57087,0.5707369) (-2.565634,0.57597277) (-2.560398,0.58120865) (-2.555162,0.58644453) (-2.549926,0.59168041) (-2.54469,0.59691629) (-2.5394541,0.60215217) (-2.5342181,0.60738806) (-2.5289821,0.61262394) (-2.5237461,0.61785983) (-2.5185101,0.62309572) (-2.5132741,0.62833162) (-2.5080381,0.63356751) (-2.5028021,0.63880341) (-2.4975662,0.6440393) (-2.4923302,0.6492752) (-2.4870942,0.6545111) (-2.4818582,0.65974701) (-2.4766222,0.66498291) (-2.4713862,0.67021882) (-2.4661502,0.67545472) (-2.4609142,0.68069063) (-2.4556783,0.68592654) (-2.4504423,0.69116245) (-2.4452063,0.69639836) (-2.4399703,0.70163428) (-2.4347343,0.70687019) (-2.4294983,0.71210611) (-2.4242623,0.71734202) (-2.4190263,0.72257794) (-2.4137904,0.72781386) (-2.4085544,0.73304978) (-2.4033184,0.7382857) (-2.3980824,0.74352162) (-2.3928464,0.74875755) (-2.3876104,0.75399347) (-2.3823744,0.7592294) (-2.3771384,0.76446532) (-2.3719025,0.76970125) (-2.3666665,0.77493718) (-2.3614305,0.78017311) (-2.3561945,0.78540904) (-2.3509585,0.79064497) (-2.3457225,0.7958809) (-2.3404865,0.80111684) (-2.3352505,0.80635277) (-2.3300146,0.81158871) (-2.3247786,0.81682464) (-2.3195426,0.82206058) (-2.3143066,0.82729651) (-2.3090706,0.83253245) (-2.3038346,0.83776839) (-2.2985986,0.84300433) (-2.2933626,0.84824027) (-2.2881266,0.85347621) (-2.2828907,0.85871215) (-2.2776547,0.86394809) (-2.2724187,0.86918404) (-2.2671827,0.87441998) (-2.2619467,0.87965593) (-2.2567107,0.88489187) (-2.2514747,0.89012782) (-2.2462387,0.89536376) (-2.2410028,0.90059971) (-2.2357668,0.90583566) (-2.2305308,0.9110716) (-2.2252948,0.91630755) (-2.2200588,0.9215435) (-2.2148228,0.92677945) (-2.2095868,0.9320154) (-2.2043508,0.93725135) (-2.1991149,0.9424873) (-2.1938789,0.94772325) (-2.1886429,0.95295921) (-2.1834069,0.95819516) (-2.1781709,0.96343111) (-2.1729349,0.96866707) (-2.1676989,0.97390302) (-2.1624629,0.97913898) (-2.157227,0.98437493) (-2.151991,0.98961089) (-2.146755,0.99484684) (-2.141519,1.0000828) (-2.136283,1.0053188) (-2.131047,1.0105547) (-2.125811,1.0157907) (-2.120575,1.0210266) (-2.1153391,1.0262626) (-2.1101031,1.0314986) (-2.1048671,1.0367345) (-2.0996311,1.0419705) (-2.0943951,1.0472064) (-2.0891591,1.0524424) (-2.0839231,1.0576784) (-2.0786871,1.0629143) (-2.0734512,1.0681503) (-2.0682152,1.0733862) (-2.0629792,1.0786222) (-2.0577432,1.0838582) (-2.0525072,1.0890941) (-2.0472712,1.0943301) (-2.0420352,1.0995661) (-2.0367992,1.104802) (-2.0315632,1.110038) (-2.0263273,1.115274) (-2.0210913,1.1205099) (-2.0158553,1.1257459) (-2.0106193,1.1309819) (-2.0053833,1.1362178) (-2.0001473,1.1414538) (-1.9949113,1.1466898) (-1.9896753,1.1519257) (-1.9844394,1.1571617) (-1.9792034,1.1623977) (-1.9739674,1.1676336) (-1.9687314,1.1728696) (-1.9634954,1.1781056) (-1.9582594,1.1833415) (-1.9530234,1.1885775) (-1.9477874,1.1938135) (-1.9425515,1.1990495) (-1.9373155,1.2042854) (-1.9320795,1.2095214) (-1.9268435,1.2147574) (-1.9216075,1.2199933) (-1.9163715,1.2252293) (-1.9111355,1.2304653) (-1.9058995,1.2357013) (-1.9006636,1.2409372) (-1.8954276,1.2461732) (-1.8901916,1.2514092) (-1.8849556,1.2566451) (-1.8797196,1.2618811) (-1.8744836,1.2671171) (-1.8692476,1.2723531) (-1.8640116,1.277589) (-1.8587757,1.282825) (-1.8535397,1.288061) (-1.8483037,1.293297) (-1.8430677,1.2985329) (-1.8378317,1.3037689) (-1.8325957,1.3090049) (-1.8273597,1.3142409) (-1.8221237,1.3194769) (-1.8168878,1.3247128) (-1.8116518,1.3299488) (-1.8064158,1.3351848) (-1.8011798,1.3404208) (-1.7959438,1.3456567) (-1.7907078,1.3508927) (-1.7854718,1.3561287) (-1.7802358,1.3613647) (-1.7749998,1.3666007) (-1.7697639,1.3718366) (-1.7645279,1.3770726) (-1.7592919,1.3823086) (-1.7540559,1.3875446) (-1.7488199,1.3927806) (-1.7435839,1.3980165) (-1.7383479,1.4032525) (-1.7331119,1.4084885) (-1.727876,1.4137245) (-1.72264,1.4189605) (-1.717404,1.4241964) (-1.712168,1.4294324) (-1.706932,1.4346684) (-1.701696,1.4399044) (-1.69646,1.4451404) (-1.691224,1.4503764) (-1.6859881,1.4556123) (-1.6807521,1.4608483) (-1.6755161,1.4660843) (-1.6702801,1.4713203) (-1.6650441,1.4765563) (-1.6598081,1.4817923) (-1.6545721,1.4870282) (-1.6493361,1.4922642) (-1.6441002,1.4975002) (-1.6388642,1.5027362) (-1.6336282,1.5079722) (-1.6283922,1.5132082) (-1.6231562,1.5184442) (-1.6179202,1.5236801) (-1.6126842,1.5289161) (-1.6074482,1.5341521) (-1.6022123,1.5393881) (-1.5969763,1.5446241) (-1.5917403,1.5498601) (-1.5865043,1.5550961) (-1.5812683,1.560332) (-1.5760323,1.565568) (-1.5707963,1.570804) (-1.5655603,1.57604) (-1.5603244,1.581276) (-1.5550884,1.586512) (-1.5498524,1.591748) (-1.5446164,1.596984) (-1.5393804,1.6022199) (-1.5341444,1.6074559) (-1.5289084,1.6126919) (-1.5236724,1.6179279) (-1.5184364,1.6231639) (-1.5132005,1.6283999) (-1.5079645,1.6336359) (-1.5027285,1.6388719) (-1.4974925,1.6441079) (-1.4922565,1.6493439) (-1.4870205,1.6545798) (-1.4817845,1.6598158) (-1.4765485,1.6650518) (-1.4713126,1.6702878) (-1.4660766,1.6755238) (-1.4608406,1.6807598) (-1.4556046,1.6859958) (-1.4503686,1.6912318) (-1.4451326,1.6964678) (-1.4398966,1.7017038) (-1.4346606,1.7069398) (-1.4294247,1.7121758) (-1.4241887,1.7174118) (-1.4189527,1.7226478) (-1.4137167,1.7278837) (-1.4084807,1.7331197) (-1.4032447,1.7383557) (-1.3980087,1.7435917) (-1.3927727,1.7488277) (-1.3875368,1.7540637) (-1.3823008,1.7592997) (-1.3770648,1.7645357) (-1.3718288,1.7697717) (-1.3665928,1.7750077) (-1.3613568,1.7802437) (-1.3561208,1.7854797) (-1.3508848,1.7907157) (-1.3456489,1.7959517) (-1.3404129,1.8011877) (-1.3351769,1.8064237) (-1.3299409,1.8116597) (-1.3247049,1.8168957) (-1.3194689,1.8221317) (-1.3142329,1.8273677) (-1.3089969,1.8326037) (-1.303761,1.8378397) (-1.298525,1.8430757) (-1.293289,1.8483117) (-1.288053,1.8535477) (-1.282817,1.8587837) (-1.277581,1.8640197) (-1.272345,1.8692557) (-1.267109,1.8744917) (-1.261873,1.8797277) (-1.2566371,1.8849637) (-1.2514011,1.8901997) (-1.2461651,1.8954357) (-1.2409291,1.9006717) (-1.2356931,1.9059077) (-1.2304571,1.9111437) (-1.2252211,1.9163797) (-1.2199851,1.9216157) (-1.2147492,1.9268517) (-1.2095132,1.9320877) (-1.2042772,1.9373237) (-1.1990412,1.9425597) (-1.1938052,1.9477957) (-1.1885692,1.9530317) (-1.1833332,1.9582677) (-1.1780972,1.9635037) (-1.1728613,1.9687397) (-1.1676253,1.9739757) (-1.1623893,1.9792118) (-1.1571533,1.9844478) (-1.1519173,1.9896838) (-1.1466813,1.9949198) (-1.1414453,2.0001558) (-1.1362093,2.0053918) (-1.1309734,2.0106278) (-1.1257374,2.0158638) (-1.1205014,2.0210998) (-1.1152654,2.0263358) (-1.1100294,2.0315718) (-1.1047934,2.0368078) (-1.0995574,2.0420439) (-1.0943214,2.0472799) (-1.0890855,2.0525159) (-1.0838495,2.0577519) (-1.0786135,2.0629879) (-1.0733775,2.0682239) (-1.0681415,2.0734599) (-1.0629055,2.0786959) (-1.0576695,2.083932) (-1.0524335,2.089168) (-1.0471976,2.094404) (-1.0419616,2.09964) (-1.0367256,2.104876) (-1.0314896,2.110112) (-1.0262536,2.115348) (-1.0210176,2.1205841) (-1.0157816,2.1258201) (-1.0105456,2.1310561) (-1.0053096,2.1362921) (-1.0000737,2.1415281) (-0.99483767,2.1467642) (-0.98960169,2.1520002) (-0.9843657,2.1572362) (-0.97912971,2.1624722) (-0.97389372,2.1677082) (-0.96865773,2.1729443) (-0.96342175,2.1781803) (-0.95818576,2.1834163) (-0.95294977,2.1886523) (-0.94771378,2.1938883) (-0.9424778,2.1991244) (-0.93724181,2.2043604) (-0.93200582,2.2095964) (-0.92676983,2.2148324) (-0.92153385,2.2200685) (-0.91629786,2.2253045) (-0.91106187,2.2305405) (-0.90582588,2.2357765) (-0.90058989,2.2410126) (-0.89535391,2.2462486) (-0.89011792,2.2514846) (-0.88488193,2.2567207) (-0.87964594,2.2619567) (-0.87440996,2.2671927) (-0.86917397,2.2724288) (-0.86393798,2.2776648) (-0.85870199,2.2829008) (-0.853466,2.2881369) (-0.84823002,2.2933729) (-0.84299403,2.2986089) (-0.83775804,2.303845) (-0.83252205,2.309081) (-0.82728607,2.314317) (-0.82205008,2.3195531) (-0.81681409,2.3247891) (-0.8115781,2.3300252) (-0.80634211,2.3352612) (-0.80110613,2.3404972) (-0.79587014,2.3457333) (-0.79063415,2.3509693) (-0.78539816,2.3562054) (-0.78016218,2.3614414) (-0.77492619,2.3666775) (-0.7696902,2.3719135) (-0.76445421,2.3771496) (-0.75921822,2.3823856) (-0.75398224,2.3876217) (-0.74874625,2.3928577) (-0.74351026,2.3980938) (-0.73827427,2.4033298) (-0.73303829,2.4085659) (-0.7278023,2.4138019) (-0.72256631,2.419038) (-0.71733032,2.424274) (-0.71209433,2.4295101) (-0.70685835,2.4347461) (-0.70162236,2.4399822) (-0.69638637,2.4452183) (-0.69115038,2.4504543) (-0.6859144,2.4556904) (-0.68067841,2.4609265) (-0.67544242,2.4661625) (-0.67020643,2.4713986) (-0.66497045,2.4766347) (-0.65973446,2.4818707) (-0.65449847,2.4871068) (-0.64926248,2.4923429) (-0.64402649,2.497579) (-0.63879051,2.502815) (-0.63355452,2.5080511) (-0.62831853,2.5132872) (-0.62308254,2.5185233) (-0.61784656,2.5237594) (-0.61261057,2.5289955) (-0.60737458,2.5342316) (-0.60213859,2.5394676) (-0.5969026,2.5447037) (-0.59166662,2.5499398) (-0.58643063,2.5551759) (-0.58119464,2.560412) (-0.57595865,2.5656481) (-0.57072267,2.5708842) (-0.56548668,2.5761203) (-0.56025069,2.5813564) (-0.5550147,2.5865925) (-0.54977871,2.5918287) (-0.54454273,2.5970648) (-0.53930674,2.6023009) (-0.53407075,2.607537) (-0.52883476,2.6127731) (-0.52359878,2.6180093) (-0.51836279,2.6232454) (-0.5131268,2.6284815) (-0.50789081,2.6337177) (-0.50265482,2.6389538) (-0.49741884,2.6441899) (-0.49218285,2.6494261) (-0.48694686,2.6546622) (-0.48171087,2.6598984) (-0.47647489,2.6651345) (-0.4712389,2.6703707) (-0.46600291,2.6756069) (-0.46076692,2.680843) (-0.45553093,2.6860792) (-0.45029495,2.6913154) (-0.44505896,2.6965516) (-0.43982297,2.7017877) (-0.43458698,2.7070239) (-0.429351,2.7122601) (-0.42411501,2.7174963) (-0.41887902,2.7227325) (-0.41364303,2.7279688) (-0.40840704,2.733205) (-0.40317106,2.7384412) (-0.39793507,2.7436774) (-0.39269908,2.7489137) (-0.38746309,2.7541499) (-0.38222711,2.7593862) (-0.37699112,2.7646224) (-0.37175513,2.7698587) (-0.36651914,2.775095) (-0.36128316,2.7803313) (-0.35604717,2.7855676) (-0.35081118,2.7908039) (-0.34557519,2.7960402) (-0.3403392,2.8012765) (-0.33510322,2.8065128) (-0.32986723,2.8117492) (-0.32463124,2.8169855) (-0.31939525,2.8222219) (-0.31415927,2.8274583) (-0.30892328,2.8326947) (-0.30368729,2.8379311) (-0.2984513,2.8431675) (-0.29321531,2.848404) (-0.28797933,2.8536404) (-0.28274334,2.8588769) (-0.27750735,2.8641134) (-0.27227136,2.8693499) (-0.26703538,2.8745864) (-0.26179939,2.879823) (-0.2565634,2.8850596) (-0.25132741,2.8902962) (-0.24609142,2.8955328) (-0.24085544,2.9007695) (-0.23561945,2.9060062) (-0.23038346,2.9112429) (-0.22514747,2.9164796) (-0.21991149,2.9217164) (-0.2146755,2.9269533) (-0.20943951,2.9321901) (-0.20420352,2.9374271) (-0.19896753,2.942664) (-0.19373155,2.9479011) (-0.18849556,2.9531381) (-0.18325957,2.9583753) (-0.17802358,2.9636125) (-0.1727876,2.9688498) (-0.16755161,2.9740872) (-0.16231562,2.9793246) (-0.15707963,2.9845622) (-0.15184364,2.9897999) (-0.14660766,2.9950377) (-0.14137167,3.0002756) (-0.13613568,3.0055137) (-0.13089969,3.0107519) (-0.12566371,3.0159903) (-0.12042772,3.021229) (-0.11519173,3.0264679) (-0.10995574,3.031707) (-0.10471976,3.0369465) (-0.099483767,3.0421864) (-0.09424778,3.0474266) (-0.089011792,3.0526674) (-0.083775804,3.0579088) (-0.078539816,3.0631509) (-0.073303829,3.0683939) (-0.068067841,3.073638) (-0.062831853,3.0788834) (-0.057595865,3.0841306) (-0.052359878,3.0893799) (-0.04712389,3.0946323) (-0.041887902,3.0998888) (-0.036651914,3.1051512) (-0.031415927,3.1104226) (-0.026179939,3.1157082) (-0.020943951,3.1210192) (-0.015707963,3.1263822) (-0.010471976,3.1318829) (-0.0052359878,3.1381238) (0.0052359878,3.1381238) (0.010471976,3.1318829) (0.015707963,3.1263822) (0.020943951,3.1210192) (0.026179939,3.1157082) (0.031415927,3.1104226) (0.036651914,3.1051512) (0.041887902,3.0998888) (0.04712389,3.0946323) (0.052359878,3.0893799) (0.057595865,3.0841306) (0.062831853,3.0788834) (0.068067841,3.073638) (0.073303829,3.0683939) (0.078539816,3.0631509) (0.083775804,3.0579088) (0.089011792,3.0526674) (0.09424778,3.0474266) (0.099483767,3.0421864) (0.10471976,3.0369465) (0.10995574,3.031707) (0.11519173,3.0264679) (0.12042772,3.021229) (0.12566371,3.0159903) (0.13089969,3.0107519) (0.13613568,3.0055137) (0.14137167,3.0002756) (0.14660766,2.9950377) (0.15184364,2.9897999) (0.15707963,2.9845622) (0.16231562,2.9793246) (0.16755161,2.9740872) (0.1727876,2.9688498) (0.17802358,2.9636125) (0.18325957,2.9583753) (0.18849556,2.9531381) (0.19373155,2.9479011) (0.19896753,2.942664) (0.20420352,2.9374271) (0.20943951,2.9321901) (0.2146755,2.9269533) (0.21991149,2.9217164) (0.22514747,2.9164796) (0.23038346,2.9112429) (0.23561945,2.9060062) (0.24085544,2.9007695) (0.24609142,2.8955328) (0.25132741,2.8902962) (0.2565634,2.8850596) (0.26179939,2.879823) (0.26703538,2.8745864) (0.27227136,2.8693499) (0.27750735,2.8641134) (0.28274334,2.8588769) (0.28797933,2.8536404) (0.29321531,2.848404) (0.2984513,2.8431675) (0.30368729,2.8379311) (0.30892328,2.8326947) (0.31415927,2.8274583) (0.31939525,2.8222219) (0.32463124,2.8169855) (0.32986723,2.8117492) (0.33510322,2.8065128) (0.3403392,2.8012765) (0.34557519,2.7960402) (0.35081118,2.7908039) (0.35604717,2.7855676) (0.36128316,2.7803313) (0.36651914,2.775095) (0.37175513,2.7698587) (0.37699112,2.7646224) (0.38222711,2.7593862) (0.38746309,2.7541499) (0.39269908,2.7489137) (0.39793507,2.7436774) (0.40317106,2.7384412) (0.40840704,2.733205) (0.41364303,2.7279688) (0.41887902,2.7227325) (0.42411501,2.7174963) (0.429351,2.7122601) (0.43458698,2.7070239) (0.43982297,2.7017877) (0.44505896,2.6965516) (0.45029495,2.6913154) (0.45553093,2.6860792) (0.46076692,2.680843) (0.46600291,2.6756069) (0.4712389,2.6703707) (0.47647489,2.6651345) (0.48171087,2.6598984) (0.48694686,2.6546622) (0.49218285,2.6494261) (0.49741884,2.6441899) (0.50265482,2.6389538) (0.50789081,2.6337177) (0.5131268,2.6284815) (0.51836279,2.6232454) (0.52359878,2.6180093) (0.52883476,2.6127731) (0.53407075,2.607537) (0.53930674,2.6023009) (0.54454273,2.5970648) (0.54977871,2.5918287) (0.5550147,2.5865925) (0.56025069,2.5813564) (0.56548668,2.5761203) (0.57072267,2.5708842) (0.57595865,2.5656481) (0.58119464,2.560412) (0.58643063,2.5551759) (0.59166662,2.5499398) (0.5969026,2.5447037) (0.60213859,2.5394676) (0.60737458,2.5342316) (0.61261057,2.5289955) (0.61784656,2.5237594) (0.62308254,2.5185233) (0.62831853,2.5132872) (0.63355452,2.5080511) (0.63879051,2.502815) (0.64402649,2.497579) (0.64926248,2.4923429) (0.65449847,2.4871068) (0.65973446,2.4818707) (0.66497045,2.4766347) (0.67020643,2.4713986) (0.67544242,2.4661625) (0.68067841,2.4609265) (0.6859144,2.4556904) (0.69115038,2.4504543) (0.69638637,2.4452183) (0.70162236,2.4399822) (0.70685835,2.4347461) (0.71209433,2.4295101) (0.71733032,2.424274) (0.72256631,2.419038) (0.7278023,2.4138019) (0.73303829,2.4085659) (0.73827427,2.4033298) (0.74351026,2.3980938) (0.74874625,2.3928577) (0.75398224,2.3876217) (0.75921822,2.3823856) (0.76445421,2.3771496) (0.7696902,2.3719135) (0.77492619,2.3666775) (0.78016218,2.3614414) (0.78539816,2.3562054) (0.79063415,2.3509693) (0.79587014,2.3457333) (0.80110613,2.3404972) (0.80634211,2.3352612) (0.8115781,2.3300252) (0.81681409,2.3247891) (0.82205008,2.3195531) (0.82728607,2.314317) (0.83252205,2.309081) (0.83775804,2.303845) (0.84299403,2.2986089) (0.84823002,2.2933729) (0.853466,2.2881369) (0.85870199,2.2829008) (0.86393798,2.2776648) (0.86917397,2.2724288) (0.87440996,2.2671927) (0.87964594,2.2619567) (0.88488193,2.2567207) (0.89011792,2.2514846) (0.89535391,2.2462486) (0.90058989,2.2410126) (0.90582588,2.2357765) (0.91106187,2.2305405) (0.91629786,2.2253045) (0.92153385,2.2200685) (0.92676983,2.2148324) (0.93200582,2.2095964) (0.93724181,2.2043604) (0.9424778,2.1991244) (0.94771378,2.1938883) (0.95294977,2.1886523) (0.95818576,2.1834163) (0.96342175,2.1781803) (0.96865773,2.1729443) (0.97389372,2.1677082) (0.97912971,2.1624722) (0.9843657,2.1572362) (0.98960169,2.1520002) (0.99483767,2.1467642) (1.0000737,2.1415281) (1.0053096,2.1362921) (1.0105456,2.1310561) (1.0157816,2.1258201) (1.0210176,2.1205841) (1.0262536,2.115348) (1.0314896,2.110112) (1.0367256,2.104876) (1.0419616,2.09964) (1.0471976,2.094404) (1.0524335,2.089168) (1.0576695,2.083932) (1.0629055,2.0786959) (1.0681415,2.0734599) (1.0733775,2.0682239) (1.0786135,2.0629879) (1.0838495,2.0577519) (1.0890855,2.0525159) (1.0943214,2.0472799) (1.0995574,2.0420439) (1.1047934,2.0368078) (1.1100294,2.0315718) (1.1152654,2.0263358) (1.1205014,2.0210998) (1.1257374,2.0158638) (1.1309734,2.0106278) (1.1362093,2.0053918) (1.1414453,2.0001558) (1.1466813,1.9949198) (1.1519173,1.9896838) (1.1571533,1.9844478) (1.1623893,1.9792118) (1.1676253,1.9739757) (1.1728613,1.9687397) (1.1780972,1.9635037) (1.1833332,1.9582677) (1.1885692,1.9530317) (1.1938052,1.9477957) (1.1990412,1.9425597) (1.2042772,1.9373237) (1.2095132,1.9320877) (1.2147492,1.9268517) (1.2199851,1.9216157) (1.2252211,1.9163797) (1.2304571,1.9111437) (1.2356931,1.9059077) (1.2409291,1.9006717) (1.2461651,1.8954357) (1.2514011,1.8901997) (1.2566371,1.8849637) (1.261873,1.8797277) (1.267109,1.8744917) (1.272345,1.8692557) (1.277581,1.8640197) (1.282817,1.8587837) (1.288053,1.8535477) (1.293289,1.8483117) (1.298525,1.8430757) (1.303761,1.8378397) (1.3089969,1.8326037) (1.3142329,1.8273677) (1.3194689,1.8221317) (1.3247049,1.8168957) (1.3299409,1.8116597) (1.3351769,1.8064237) (1.3404129,1.8011877) (1.3456489,1.7959517) (1.3508848,1.7907157) (1.3561208,1.7854797) (1.3613568,1.7802437) (1.3665928,1.7750077) (1.3718288,1.7697717) (1.3770648,1.7645357) (1.3823008,1.7592997) (1.3875368,1.7540637) (1.3927727,1.7488277) (1.3980087,1.7435917) (1.4032447,1.7383557) (1.4084807,1.7331197) (1.4137167,1.7278837) (1.4189527,1.7226478) (1.4241887,1.7174118) (1.4294247,1.7121758) (1.4346606,1.7069398) (1.4398966,1.7017038) (1.4451326,1.6964678) (1.4503686,1.6912318) (1.4556046,1.6859958) (1.4608406,1.6807598) (1.4660766,1.6755238) (1.4713126,1.6702878) (1.4765485,1.6650518) (1.4817845,1.6598158) (1.4870205,1.6545798) (1.4922565,1.6493439) (1.4974925,1.6441079) (1.5027285,1.6388719) (1.5079645,1.6336359) (1.5132005,1.6283999) (1.5184364,1.6231639) (1.5236724,1.6179279) (1.5289084,1.6126919) (1.5341444,1.6074559) (1.5393804,1.6022199) (1.5446164,1.596984) (1.5498524,1.591748) (1.5550884,1.586512) (1.5603244,1.581276) (1.5655603,1.57604) (1.5707963,1.570804) (1.5760323,1.565568) (1.5812683,1.560332) (1.5865043,1.5550961) (1.5917403,1.5498601) (1.5969763,1.5446241) (1.6022123,1.5393881) (1.6074482,1.5341521) (1.6126842,1.5289161) (1.6179202,1.5236801) (1.6231562,1.5184442) (1.6283922,1.5132082) (1.6336282,1.5079722) (1.6388642,1.5027362) (1.6441002,1.4975002) (1.6493361,1.4922642) (1.6545721,1.4870282) (1.6598081,1.4817923) (1.6650441,1.4765563) (1.6702801,1.4713203) (1.6755161,1.4660843) (1.6807521,1.4608483) (1.6859881,1.4556123) (1.691224,1.4503764) (1.69646,1.4451404) (1.701696,1.4399044) (1.706932,1.4346684) (1.712168,1.4294324) (1.717404,1.4241964) (1.72264,1.4189605) (1.727876,1.4137245) (1.7331119,1.4084885) (1.7383479,1.4032525) (1.7435839,1.3980165) (1.7488199,1.3927806) (1.7540559,1.3875446) (1.7592919,1.3823086) (1.7645279,1.3770726) (1.7697639,1.3718366) (1.7749998,1.3666007) (1.7802358,1.3613647) (1.7854718,1.3561287) (1.7907078,1.3508927) (1.7959438,1.3456567) (1.8011798,1.3404208) (1.8064158,1.3351848) (1.8116518,1.3299488) (1.8168878,1.3247128) (1.8221237,1.3194769) (1.8273597,1.3142409) (1.8325957,1.3090049) (1.8378317,1.3037689) (1.8430677,1.2985329) (1.8483037,1.293297) (1.8535397,1.288061) (1.8587757,1.282825) (1.8640116,1.277589) (1.8692476,1.2723531) (1.8744836,1.2671171) (1.8797196,1.2618811) (1.8849556,1.2566451) (1.8901916,1.2514092) (1.8954276,1.2461732) (1.9006636,1.2409372) (1.9058995,1.2357013) (1.9111355,1.2304653) (1.9163715,1.2252293) (1.9216075,1.2199933) (1.9268435,1.2147574) (1.9320795,1.2095214) (1.9373155,1.2042854) (1.9425515,1.1990495) (1.9477874,1.1938135) (1.9530234,1.1885775) (1.9582594,1.1833415) (1.9634954,1.1781056) (1.9687314,1.1728696) (1.9739674,1.1676336) (1.9792034,1.1623977) (1.9844394,1.1571617) (1.9896753,1.1519257) (1.9949113,1.1466898) (2.0001473,1.1414538) (2.0053833,1.1362178) (2.0106193,1.1309819) (2.0158553,1.1257459) (2.0210913,1.1205099) (2.0263273,1.115274) (2.0315632,1.110038) (2.0367992,1.104802) (2.0420352,1.0995661) (2.0472712,1.0943301) (2.0525072,1.0890941) (2.0577432,1.0838582) (2.0629792,1.0786222) (2.0682152,1.0733862) (2.0734512,1.0681503) (2.0786871,1.0629143) (2.0839231,1.0576784) (2.0891591,1.0524424) (2.0943951,1.0472064) (2.0996311,1.0419705) (2.1048671,1.0367345) (2.1101031,1.0314986) (2.1153391,1.0262626) (2.120575,1.0210266) (2.125811,1.0157907) (2.131047,1.0105547) (2.136283,1.0053188) (2.141519,1.0000828) (2.146755,0.99484684) (2.151991,0.98961089) (2.157227,0.98437493) (2.1624629,0.97913898) (2.1676989,0.97390302) (2.1729349,0.96866707) (2.1781709,0.96343111) (2.1834069,0.95819516) (2.1886429,0.95295921) (2.1938789,0.94772325) (2.1991149,0.9424873) (2.2043508,0.93725135) (2.2095868,0.9320154) (2.2148228,0.92677945) (2.2200588,0.9215435) (2.2252948,0.91630755) (2.2305308,0.9110716) (2.2357668,0.90583566) (2.2410028,0.90059971) (2.2462387,0.89536376) (2.2514747,0.89012782) (2.2567107,0.88489187) (2.2619467,0.87965593) (2.2671827,0.87441998) (2.2724187,0.86918404) (2.2776547,0.86394809) (2.2828907,0.85871215) (2.2881266,0.85347621) (2.2933626,0.84824027) (2.2985986,0.84300433) (2.3038346,0.83776839) (2.3090706,0.83253245) (2.3143066,0.82729651) (2.3195426,0.82206058) (2.3247786,0.81682464) (2.3300146,0.81158871) (2.3352505,0.80635277) (2.3404865,0.80111684) (2.3457225,0.7958809) (2.3509585,0.79064497) (2.3561945,0.78540904) (2.3614305,0.78017311) (2.3666665,0.77493718) (2.3719025,0.76970125) (2.3771384,0.76446532) (2.3823744,0.7592294) (2.3876104,0.75399347) (2.3928464,0.74875755) (2.3980824,0.74352162) (2.4033184,0.7382857) (2.4085544,0.73304978) (2.4137904,0.72781386) (2.4190263,0.72257794) (2.4242623,0.71734202) (2.4294983,0.71210611) (2.4347343,0.70687019) (2.4399703,0.70163428) (2.4452063,0.69639836) (2.4504423,0.69116245) (2.4556783,0.68592654) (2.4609142,0.68069063) (2.4661502,0.67545472) (2.4713862,0.67021882) (2.4766222,0.66498291) (2.4818582,0.65974701) (2.4870942,0.6545111) (2.4923302,0.6492752) (2.4975662,0.6440393) (2.5028021,0.63880341) (2.5080381,0.63356751) (2.5132741,0.62833162) (2.5185101,0.62309572) (2.5237461,0.61785983) (2.5289821,0.61262394) (2.5342181,0.60738806) (2.5394541,0.60215217) (2.54469,0.59691629) (2.549926,0.59168041) (2.555162,0.58644453) (2.560398,0.58120865) (2.565634,0.57597277) (2.57087,0.5707369) (2.576106,0.56550103) (2.581342,0.56026516) (2.586578,0.5550293) (2.5918139,0.54979343) (2.5970499,0.54455757) (2.6022859,0.53932172) (2.6075219,0.53408586) (2.6127579,0.52885001) (2.6179939,0.52361416) (2.6232299,0.51837831) (2.6284659,0.51314247) (2.6337018,0.50790663) (2.6389378,0.50267079) (2.6441738,0.49743496) (2.6494098,0.49219913) (2.6546458,0.4869633) (2.6598818,0.48172747) (2.6651178,0.47649166) (2.6703538,0.47125584) (2.6755897,0.46602003) (2.6808257,0.46078422) (2.6860617,0.45554842) (2.6912977,0.45031262) (2.6965337,0.44507682) (2.7017697,0.43984104) (2.7070057,0.43460525) (2.7122417,0.42936947) (2.7174776,0.4241337) (2.7227136,0.41889793) (2.7279496,0.41366217) (2.7331856,0.40842641) (2.7384216,0.40319066) (2.7436576,0.39795492) (2.7488936,0.39271918) (2.7541296,0.38748345) (2.7593655,0.38224773) (2.7646015,0.37701201) (2.7698375,0.3717763) (2.7750735,0.3665406) (2.7803095,0.36130491) (2.7855455,0.35606923) (2.7907815,0.35083356) (2.7960175,0.3455979) (2.8012534,0.34036224) (2.8064894,0.3351266) (2.8117254,0.32989097) (2.8169614,0.32465535) (2.8221974,0.31941975) (2.8274334,0.31418415) (2.8326694,0.30894857) (2.8379054,0.30371301) (2.8431414,0.29847746) (2.8483773,0.29324192) (2.8536133,0.28800641) (2.8588493,0.28277091) (2.8640853,0.27753542) (2.8693213,0.27229996) (2.8745573,0.26706452) (2.8797933,0.2618291) (2.8850293,0.25659371) (2.8902652,0.25135834) (2.8955012,0.24612299) (2.9007372,0.24088768) (2.9059732,0.23565239) (2.9112092,0.23041714) (2.9164452,0.22518192) (2.9216812,0.21994674) (2.9269172,0.2147116) (2.9321531,0.2094765) (2.9373891,0.20424145) (2.9426251,0.19900644) (2.9478611,0.19377149) (2.9530971,0.1885366) (2.9583331,0.18330177) (2.9635691,0.17806701) (2.9688051,0.17283233) (2.974041,0.16759772) (2.979277,0.16236321) (2.984513,0.15712879) (2.989749,0.15189448) (2.994985,0.1466603) (3.000221,0.14142625) (3.005457,0.13619234) (3.010693,0.13095861) (3.0159289,0.12572506) (3.0211649,0.12049172) (3.0264009,0.11525863) (3.0316369,0.11002581) (3.0368729,0.10479331) (3.0421089,0.099561178) (3.0473449,0.094329473) (3.0525809,0.089098273) (3.0578168,0.08386767) (3.0630528,0.078637786) (3.0682888,0.073408772) (3.0735248,0.06818083) (3.0787608,0.062954227) (3.0839968,0.057729325) (3.0892328,0.052506634) (3.0944688,0.047286885) (3.0997048,0.042071173) (3.1049407,0.036861213) (3.1101767,0.031659844) (3.1154127,0.026472131) (3.1206487,0.02130805) (3.1258847,0.016190229) (3.1311207,0.011182367) (3.1363567,0.0065420497) (3.1415927,0.0039220882) (3.1415927,-0.0039220882) (3.1363567,-0.0065420497) (3.1311207,-0.011182367) (3.1258847,-0.016190229) (3.1206487,-0.02130805) (3.1154127,-0.026472131) (3.1101767,-0.031659844) (3.1049407,-0.036861213) (3.0997048,-0.042071173) (3.0944688,-0.047286885) (3.0892328,-0.052506634) (3.0839968,-0.057729325) (3.0787608,-0.062954227) (3.0735248,-0.06818083) (3.0682888,-0.073408772) (3.0630528,-0.078637786) (3.0578168,-0.08386767) (3.0525809,-0.089098273) (3.0473449,-0.094329473) (3.0421089,-0.099561178) (3.0368729,-0.10479331) (3.0316369,-0.11002581) (3.0264009,-0.11525863) (3.0211649,-0.12049172) (3.0159289,-0.12572506) (3.010693,-0.13095861) (3.005457,-0.13619234) (3.000221,-0.14142625) (2.994985,-0.1466603) (2.989749,-0.15189448) (2.984513,-0.15712879) (2.979277,-0.16236321) (2.974041,-0.16759772) (2.9688051,-0.17283233) (2.9635691,-0.17806701) (2.9583331,-0.18330177) (2.9530971,-0.1885366) (2.9478611,-0.19377149) (2.9426251,-0.19900644) (2.9373891,-0.20424145) (2.9321531,-0.2094765) (2.9269172,-0.2147116) (2.9216812,-0.21994674) (2.9164452,-0.22518192) (2.9112092,-0.23041714) (2.9059732,-0.23565239) (2.9007372,-0.24088768) (2.8955012,-0.24612299) (2.8902652,-0.25135834) (2.8850293,-0.25659371) (2.8797933,-0.2618291) (2.8745573,-0.26706452) (2.8693213,-0.27229996) (2.8640853,-0.27753542) (2.8588493,-0.28277091) (2.8536133,-0.28800641) (2.8483773,-0.29324192) (2.8431414,-0.29847746) (2.8379054,-0.30371301) (2.8326694,-0.30894857) (2.8274334,-0.31418415) (2.8221974,-0.31941975) (2.8169614,-0.32465535) (2.8117254,-0.32989097) (2.8064894,-0.3351266) (2.8012534,-0.34036224) (2.7960175,-0.3455979) (2.7907815,-0.35083356) (2.7855455,-0.35606923) (2.7803095,-0.36130491) (2.7750735,-0.3665406) (2.7698375,-0.3717763) (2.7646015,-0.37701201) (2.7593655,-0.38224773) (2.7541296,-0.38748345) (2.7488936,-0.39271918) (2.7436576,-0.39795492) (2.7384216,-0.40319066) (2.7331856,-0.40842641) (2.7279496,-0.41366217) (2.7227136,-0.41889793) (2.7174776,-0.4241337) (2.7122417,-0.42936947) (2.7070057,-0.43460525) (2.7017697,-0.43984104) (2.6965337,-0.44507682) (2.6912977,-0.45031262) (2.6860617,-0.45554842) (2.6808257,-0.46078422) (2.6755897,-0.46602003) (2.6703538,-0.47125584) (2.6651178,-0.47649166) (2.6598818,-0.48172747) (2.6546458,-0.4869633) (2.6494098,-0.49219913) (2.6441738,-0.49743496) (2.6389378,-0.50267079) (2.6337018,-0.50790663) (2.6284659,-0.51314247) (2.6232299,-0.51837831) (2.6179939,-0.52361416) (2.6127579,-0.52885001) (2.6075219,-0.53408586) (2.6022859,-0.53932172) (2.5970499,-0.54455757) (2.5918139,-0.54979343) (2.586578,-0.5550293) (2.581342,-0.56026516) (2.576106,-0.56550103) (2.57087,-0.5707369) (2.565634,-0.57597277) (2.560398,-0.58120865) (2.555162,-0.58644453) (2.549926,-0.59168041) (2.54469,-0.59691629) (2.5394541,-0.60215217) (2.5342181,-0.60738806) (2.5289821,-0.61262394) (2.5237461,-0.61785983) (2.5185101,-0.62309572) (2.5132741,-0.62833162) (2.5080381,-0.63356751) (2.5028021,-0.63880341) (2.4975662,-0.6440393) (2.4923302,-0.6492752) (2.4870942,-0.6545111) (2.4818582,-0.65974701) (2.4766222,-0.66498291) (2.4713862,-0.67021882) (2.4661502,-0.67545472) (2.4609142,-0.68069063) (2.4556783,-0.68592654) (2.4504423,-0.69116245) (2.4452063,-0.69639836) (2.4399703,-0.70163428) (2.4347343,-0.70687019) (2.4294983,-0.71210611) (2.4242623,-0.71734202) (2.4190263,-0.72257794) (2.4137904,-0.72781386) (2.4085544,-0.73304978) (2.4033184,-0.7382857) (2.3980824,-0.74352162) (2.3928464,-0.74875755) (2.3876104,-0.75399347) (2.3823744,-0.7592294) (2.3771384,-0.76446532) (2.3719025,-0.76970125) (2.3666665,-0.77493718) (2.3614305,-0.78017311) (2.3561945,-0.78540904) (2.3509585,-0.79064497) (2.3457225,-0.7958809) (2.3404865,-0.80111684) (2.3352505,-0.80635277) (2.3300146,-0.81158871) (2.3247786,-0.81682464) (2.3195426,-0.82206058) (2.3143066,-0.82729651) (2.3090706,-0.83253245) (2.3038346,-0.83776839) (2.2985986,-0.84300433) (2.2933626,-0.84824027) (2.2881266,-0.85347621) (2.2828907,-0.85871215) (2.2776547,-0.86394809) (2.2724187,-0.86918404) (2.2671827,-0.87441998) (2.2619467,-0.87965593) (2.2567107,-0.88489187) (2.2514747,-0.89012782) (2.2462387,-0.89536376) (2.2410028,-0.90059971) (2.2357668,-0.90583566) (2.2305308,-0.9110716) (2.2252948,-0.91630755) (2.2200588,-0.9215435) (2.2148228,-0.92677945) (2.2095868,-0.9320154) (2.2043508,-0.93725135) (2.1991149,-0.9424873) (2.1938789,-0.94772325) (2.1886429,-0.95295921) (2.1834069,-0.95819516) (2.1781709,-0.96343111) (2.1729349,-0.96866707) (2.1676989,-0.97390302) (2.1624629,-0.97913898) (2.157227,-0.98437493) (2.151991,-0.98961089) (2.146755,-0.99484684) (2.141519,-1.0000828) (2.136283,-1.0053188) (2.131047,-1.0105547) (2.125811,-1.0157907) (2.120575,-1.0210266) (2.1153391,-1.0262626) (2.1101031,-1.0314986) (2.1048671,-1.0367345) (2.0996311,-1.0419705) (2.0943951,-1.0472064) (2.0891591,-1.0524424) (2.0839231,-1.0576784) (2.0786871,-1.0629143) (2.0734512,-1.0681503) (2.0682152,-1.0733862) (2.0629792,-1.0786222) (2.0577432,-1.0838582) (2.0525072,-1.0890941) (2.0472712,-1.0943301) (2.0420352,-1.0995661) (2.0367992,-1.104802) (2.0315632,-1.110038) (2.0263273,-1.115274) (2.0210913,-1.1205099) (2.0158553,-1.1257459) (2.0106193,-1.1309819) (2.0053833,-1.1362178) (2.0001473,-1.1414538) (1.9949113,-1.1466898) (1.9896753,-1.1519257) (1.9844394,-1.1571617) (1.9792034,-1.1623977) (1.9739674,-1.1676336) (1.9687314,-1.1728696) (1.9634954,-1.1781056) (1.9582594,-1.1833415) (1.9530234,-1.1885775) (1.9477874,-1.1938135) (1.9425515,-1.1990495) (1.9373155,-1.2042854) (1.9320795,-1.2095214) (1.9268435,-1.2147574) (1.9216075,-1.2199933) (1.9163715,-1.2252293) (1.9111355,-1.2304653) (1.9058995,-1.2357013) (1.9006636,-1.2409372) (1.8954276,-1.2461732) (1.8901916,-1.2514092) (1.8849556,-1.2566451) (1.8797196,-1.2618811) (1.8744836,-1.2671171) (1.8692476,-1.2723531) (1.8640116,-1.277589) (1.8587757,-1.282825) (1.8535397,-1.288061) (1.8483037,-1.293297) (1.8430677,-1.2985329) (1.8378317,-1.3037689) (1.8325957,-1.3090049) (1.8273597,-1.3142409) (1.8221237,-1.3194769) (1.8168878,-1.3247128) (1.8116518,-1.3299488) (1.8064158,-1.3351848) (1.8011798,-1.3404208) (1.7959438,-1.3456567) (1.7907078,-1.3508927) (1.7854718,-1.3561287) (1.7802358,-1.3613647) (1.7749998,-1.3666007) (1.7697639,-1.3718366) (1.7645279,-1.3770726) (1.7592919,-1.3823086) (1.7540559,-1.3875446) (1.7488199,-1.3927806) (1.7435839,-1.3980165) (1.7383479,-1.4032525) (1.7331119,-1.4084885) (1.727876,-1.4137245) (1.72264,-1.4189605) (1.717404,-1.4241964) (1.712168,-1.4294324) (1.706932,-1.4346684) (1.701696,-1.4399044) (1.69646,-1.4451404) (1.691224,-1.4503764) (1.6859881,-1.4556123) (1.6807521,-1.4608483) (1.6755161,-1.4660843) (1.6702801,-1.4713203) (1.6650441,-1.4765563) (1.6598081,-1.4817923) (1.6545721,-1.4870282) (1.6493361,-1.4922642) (1.6441002,-1.4975002) (1.6388642,-1.5027362) (1.6336282,-1.5079722) (1.6283922,-1.5132082) (1.6231562,-1.5184442) (1.6179202,-1.5236801) (1.6126842,-1.5289161) (1.6074482,-1.5341521) (1.6022123,-1.5393881) (1.5969763,-1.5446241) (1.5917403,-1.5498601) (1.5865043,-1.5550961) (1.5812683,-1.560332) (1.5760323,-1.565568) (1.5707963,-1.570804) (1.5655603,-1.57604) (1.5603244,-1.581276) (1.5550884,-1.586512) (1.5498524,-1.591748) (1.5446164,-1.596984) (1.5393804,-1.6022199) (1.5341444,-1.6074559) (1.5289084,-1.6126919) (1.5236724,-1.6179279) (1.5184364,-1.6231639) (1.5132005,-1.6283999) (1.5079645,-1.6336359) (1.5027285,-1.6388719) (1.4974925,-1.6441079) (1.4922565,-1.6493439) (1.4870205,-1.6545798) (1.4817845,-1.6598158) (1.4765485,-1.6650518) (1.4713126,-1.6702878) (1.4660766,-1.6755238) (1.4608406,-1.6807598) (1.4556046,-1.6859958) (1.4503686,-1.6912318) (1.4451326,-1.6964678) (1.4398966,-1.7017038) (1.4346606,-1.7069398) (1.4294247,-1.7121758) (1.4241887,-1.7174118) (1.4189527,-1.7226478) (1.4137167,-1.7278837) (1.4084807,-1.7331197) (1.4032447,-1.7383557) (1.3980087,-1.7435917) (1.3927727,-1.7488277) (1.3875368,-1.7540637) (1.3823008,-1.7592997) (1.3770648,-1.7645357) (1.3718288,-1.7697717) (1.3665928,-1.7750077) (1.3613568,-1.7802437) (1.3561208,-1.7854797) (1.3508848,-1.7907157) (1.3456489,-1.7959517) (1.3404129,-1.8011877) (1.3351769,-1.8064237) (1.3299409,-1.8116597) (1.3247049,-1.8168957) (1.3194689,-1.8221317) (1.3142329,-1.8273677) (1.3089969,-1.8326037) (1.303761,-1.8378397) (1.298525,-1.8430757) (1.293289,-1.8483117) (1.288053,-1.8535477) (1.282817,-1.8587837) (1.277581,-1.8640197) (1.272345,-1.8692557) (1.267109,-1.8744917) (1.261873,-1.8797277) (1.2566371,-1.8849637) (1.2514011,-1.8901997) (1.2461651,-1.8954357) (1.2409291,-1.9006717) (1.2356931,-1.9059077) (1.2304571,-1.9111437) (1.2252211,-1.9163797) (1.2199851,-1.9216157) (1.2147492,-1.9268517) (1.2095132,-1.9320877) (1.2042772,-1.9373237) (1.1990412,-1.9425597) (1.1938052,-1.9477957) (1.1885692,-1.9530317) (1.1833332,-1.9582677) (1.1780972,-1.9635037) (1.1728613,-1.9687397) (1.1676253,-1.9739757) (1.1623893,-1.9792118) (1.1571533,-1.9844478) (1.1519173,-1.9896838) (1.1466813,-1.9949198) (1.1414453,-2.0001558) (1.1362093,-2.0053918) (1.1309734,-2.0106278) (1.1257374,-2.0158638) (1.1205014,-2.0210998) (1.1152654,-2.0263358) (1.1100294,-2.0315718) (1.1047934,-2.0368078) (1.0995574,-2.0420439) (1.0943214,-2.0472799) (1.0890855,-2.0525159) (1.0838495,-2.0577519) (1.0786135,-2.0629879) (1.0733775,-2.0682239) (1.0681415,-2.0734599) (1.0629055,-2.0786959) (1.0576695,-2.083932) (1.0524335,-2.089168) (1.0471976,-2.094404) (1.0419616,-2.09964) (1.0367256,-2.104876) (1.0314896,-2.110112) (1.0262536,-2.115348) (1.0210176,-2.1205841) (1.0157816,-2.1258201) (1.0105456,-2.1310561) (1.0053096,-2.1362921) (1.0000737,-2.1415281) (0.99483767,-2.1467642) (0.98960169,-2.1520002) (0.9843657,-2.1572362) (0.97912971,-2.1624722) (0.97389372,-2.1677082) (0.96865773,-2.1729443) (0.96342175,-2.1781803) (0.95818576,-2.1834163) (0.95294977,-2.1886523) (0.94771378,-2.1938883) (0.9424778,-2.1991244) (0.93724181,-2.2043604) (0.93200582,-2.2095964) (0.92676983,-2.2148324) (0.92153385,-2.2200685) (0.91629786,-2.2253045) (0.91106187,-2.2305405) (0.90582588,-2.2357765) (0.90058989,-2.2410126) (0.89535391,-2.2462486) (0.89011792,-2.2514846) (0.88488193,-2.2567207) (0.87964594,-2.2619567) (0.87440996,-2.2671927) (0.86917397,-2.2724288) (0.86393798,-2.2776648) (0.85870199,-2.2829008) (0.853466,-2.2881369) (0.84823002,-2.2933729) (0.84299403,-2.2986089) (0.83775804,-2.303845) (0.83252205,-2.309081) (0.82728607,-2.314317) (0.82205008,-2.3195531) (0.81681409,-2.3247891) (0.8115781,-2.3300252) (0.80634211,-2.3352612) (0.80110613,-2.3404972) (0.79587014,-2.3457333) (0.79063415,-2.3509693) (0.78539816,-2.3562054) (0.78016218,-2.3614414) (0.77492619,-2.3666775) (0.7696902,-2.3719135) (0.76445421,-2.3771496) (0.75921822,-2.3823856) (0.75398224,-2.3876217) (0.74874625,-2.3928577) (0.74351026,-2.3980938) (0.73827427,-2.4033298) (0.73303829,-2.4085659) (0.7278023,-2.4138019) (0.72256631,-2.419038) (0.71733032,-2.424274) (0.71209433,-2.4295101) (0.70685835,-2.4347461) (0.70162236,-2.4399822) (0.69638637,-2.4452183) (0.69115038,-2.4504543) (0.6859144,-2.4556904) (0.68067841,-2.4609265) (0.67544242,-2.4661625) (0.67020643,-2.4713986) (0.66497045,-2.4766347) (0.65973446,-2.4818707) (0.65449847,-2.4871068) (0.64926248,-2.4923429) (0.64402649,-2.497579) (0.63879051,-2.502815) (0.63355452,-2.5080511) (0.62831853,-2.5132872) (0.62308254,-2.5185233) (0.61784656,-2.5237594) (0.61261057,-2.5289955) (0.60737458,-2.5342316) (0.60213859,-2.5394676) (0.5969026,-2.5447037) (0.59166662,-2.5499398) (0.58643063,-2.5551759) (0.58119464,-2.560412) (0.57595865,-2.5656481) (0.57072267,-2.5708842) (0.56548668,-2.5761203) (0.56025069,-2.5813564) (0.5550147,-2.5865925) (0.54977871,-2.5918287) (0.54454273,-2.5970648) (0.53930674,-2.6023009) (0.53407075,-2.607537) (0.52883476,-2.6127731) (0.52359878,-2.6180093) (0.51836279,-2.6232454) (0.5131268,-2.6284815) (0.50789081,-2.6337177) (0.50265482,-2.6389538) (0.49741884,-2.6441899) (0.49218285,-2.6494261) (0.48694686,-2.6546622) (0.48171087,-2.6598984) (0.47647489,-2.6651345) (0.4712389,-2.6703707) (0.46600291,-2.6756069) (0.46076692,-2.680843) (0.45553093,-2.6860792) (0.45029495,-2.6913154) (0.44505896,-2.6965516) (0.43982297,-2.7017877) (0.43458698,-2.7070239) (0.429351,-2.7122601) (0.42411501,-2.7174963) (0.41887902,-2.7227325) (0.41364303,-2.7279688) (0.40840704,-2.733205) (0.40317106,-2.7384412) (0.39793507,-2.7436774) (0.39269908,-2.7489137) (0.38746309,-2.7541499) (0.38222711,-2.7593862) (0.37699112,-2.7646224) (0.37175513,-2.7698587) (0.36651914,-2.775095) (0.36128316,-2.7803313) (0.35604717,-2.7855676) (0.35081118,-2.7908039) (0.34557519,-2.7960402) (0.3403392,-2.8012765) (0.33510322,-2.8065128) (0.32986723,-2.8117492) (0.32463124,-2.8169855) (0.31939525,-2.8222219) (0.31415927,-2.8274583) (0.30892328,-2.8326947) (0.30368729,-2.8379311) (0.2984513,-2.8431675) (0.29321531,-2.848404) (0.28797933,-2.8536404) (0.28274334,-2.8588769) (0.27750735,-2.8641134) (0.27227136,-2.8693499) (0.26703538,-2.8745864) (0.26179939,-2.879823) (0.2565634,-2.8850596) (0.25132741,-2.8902962) (0.24609142,-2.8955328) (0.24085544,-2.9007695) (0.23561945,-2.9060062) (0.23038346,-2.9112429) (0.22514747,-2.9164796) (0.21991149,-2.9217164) (0.2146755,-2.9269533) (0.20943951,-2.9321901) (0.20420352,-2.9374271) (0.19896753,-2.942664) (0.19373155,-2.9479011) (0.18849556,-2.9531381) (0.18325957,-2.9583753) (0.17802358,-2.9636125) (0.1727876,-2.9688498) (0.16755161,-2.9740872) (0.16231562,-2.9793246) (0.15707963,-2.9845622) (0.15184364,-2.9897999) (0.14660766,-2.9950377) (0.14137167,-3.0002756) (0.13613568,-3.0055137) (0.13089969,-3.0107519) (0.12566371,-3.0159903) (0.12042772,-3.021229) (0.11519173,-3.0264679) (0.10995574,-3.031707) (0.10471976,-3.0369465) (0.099483767,-3.0421864) (0.09424778,-3.0474266) (0.089011792,-3.0526674) (0.083775804,-3.0579088) (0.078539816,-3.0631509) (0.073303829,-3.0683939) (0.068067841,-3.073638) (0.062831853,-3.0788834) (0.057595865,-3.0841306) (0.052359878,-3.0893799) (0.04712389,-3.0946323) (0.041887902,-3.0998888) (0.036651914,-3.1051512) (0.031415927,-3.1104226) (0.026179939,-3.1157082) (0.020943951,-3.1210192) (0.015707963,-3.1263822) (0.010471976,-3.1318829) (0.0052359878,-3.1381238) (-0.0052359878,-3.1381238) (-0.010471976,-3.1318829) (-0.015707963,-3.1263822) (-0.020943951,-3.1210192) (-0.026179939,-3.1157082) (-0.031415927,-3.1104226) (-0.036651914,-3.1051512) (-0.041887902,-3.0998888) (-0.04712389,-3.0946323) (-0.052359878,-3.0893799) (-0.057595865,-3.0841306) (-0.062831853,-3.0788834) (-0.068067841,-3.073638) (-0.073303829,-3.0683939) (-0.078539816,-3.0631509) (-0.083775804,-3.0579088) (-0.089011792,-3.0526674) (-0.09424778,-3.0474266) (-0.099483767,-3.0421864) (-0.10471976,-3.0369465) (-0.10995574,-3.031707) (-0.11519173,-3.0264679) (-0.12042772,-3.021229) (-0.12566371,-3.0159903) (-0.13089969,-3.0107519) (-0.13613568,-3.0055137) (-0.14137167,-3.0002756) (-0.14660766,-2.9950377) (-0.15184364,-2.9897999) (-0.15707963,-2.9845622) (-0.16231562,-2.9793246) (-0.16755161,-2.9740872) (-0.1727876,-2.9688498) (-0.17802358,-2.9636125) (-0.18325957,-2.9583753) (-0.18849556,-2.9531381) (-0.19373155,-2.9479011) (-0.19896753,-2.942664) (-0.20420352,-2.9374271) (-0.20943951,-2.9321901) (-0.2146755,-2.9269533) (-0.21991149,-2.9217164) (-0.22514747,-2.9164796) (-0.23038346,-2.9112429) (-0.23561945,-2.9060062) (-0.24085544,-2.9007695) (-0.24609142,-2.8955328) (-0.25132741,-2.8902962) (-0.2565634,-2.8850596) (-0.26179939,-2.879823) (-0.26703538,-2.8745864) (-0.27227136,-2.8693499) (-0.27750735,-2.8641134) (-0.28274334,-2.8588769) (-0.28797933,-2.8536404) (-0.29321531,-2.848404) (-0.2984513,-2.8431675) (-0.30368729,-2.8379311) (-0.30892328,-2.8326947) (-0.31415927,-2.8274583) (-0.31939525,-2.8222219) (-0.32463124,-2.8169855) (-0.32986723,-2.8117492) (-0.33510322,-2.8065128) (-0.3403392,-2.8012765) (-0.34557519,-2.7960402) (-0.35081118,-2.7908039) (-0.35604717,-2.7855676) (-0.36128316,-2.7803313) (-0.36651914,-2.775095) (-0.37175513,-2.7698587) (-0.37699112,-2.7646224) (-0.38222711,-2.7593862) (-0.38746309,-2.7541499) (-0.39269908,-2.7489137) (-0.39793507,-2.7436774) (-0.40317106,-2.7384412) (-0.40840704,-2.733205) (-0.41364303,-2.7279688) (-0.41887902,-2.7227325) (-0.42411501,-2.7174963) (-0.429351,-2.7122601) (-0.43458698,-2.7070239) (-0.43982297,-2.7017877) (-0.44505896,-2.6965516) (-0.45029495,-2.6913154) (-0.45553093,-2.6860792) (-0.46076692,-2.680843) (-0.46600291,-2.6756069) (-0.4712389,-2.6703707) (-0.47647489,-2.6651345) (-0.48171087,-2.6598984) (-0.48694686,-2.6546622) (-0.49218285,-2.6494261) (-0.49741884,-2.6441899) (-0.50265482,-2.6389538) (-0.50789081,-2.6337177) (-0.5131268,-2.6284815) (-0.51836279,-2.6232454) (-0.52359878,-2.6180093) (-0.52883476,-2.6127731) (-0.53407075,-2.607537) (-0.53930674,-2.6023009) (-0.54454273,-2.5970648) (-0.54977871,-2.5918287) (-0.5550147,-2.5865925) (-0.56025069,-2.5813564) (-0.56548668,-2.5761203) (-0.57072267,-2.5708842) (-0.57595865,-2.5656481) (-0.58119464,-2.560412) (-0.58643063,-2.5551759) (-0.59166662,-2.5499398) (-0.5969026,-2.5447037) (-0.60213859,-2.5394676) (-0.60737458,-2.5342316) (-0.61261057,-2.5289955) (-0.61784656,-2.5237594) (-0.62308254,-2.5185233) (-0.62831853,-2.5132872) (-0.63355452,-2.5080511) (-0.63879051,-2.502815) (-0.64402649,-2.497579) (-0.64926248,-2.4923429) (-0.65449847,-2.4871068) (-0.65973446,-2.4818707) (-0.66497045,-2.4766347) (-0.67020643,-2.4713986) (-0.67544242,-2.4661625) (-0.68067841,-2.4609265) (-0.6859144,-2.4556904) (-0.69115038,-2.4504543) (-0.69638637,-2.4452183) (-0.70162236,-2.4399822) (-0.70685835,-2.4347461) (-0.71209433,-2.4295101) (-0.71733032,-2.424274) (-0.72256631,-2.419038) (-0.7278023,-2.4138019) (-0.73303829,-2.4085659) (-0.73827427,-2.4033298) (-0.74351026,-2.3980938) (-0.74874625,-2.3928577) (-0.75398224,-2.3876217) (-0.75921822,-2.3823856) (-0.76445421,-2.3771496) (-0.7696902,-2.3719135) (-0.77492619,-2.3666775) (-0.78016218,-2.3614414) (-0.78539816,-2.3562054) (-0.79063415,-2.3509693) (-0.79587014,-2.3457333) (-0.80110613,-2.3404972) (-0.80634211,-2.3352612) (-0.8115781,-2.3300252) (-0.81681409,-2.3247891) (-0.82205008,-2.3195531) (-0.82728607,-2.314317) (-0.83252205,-2.309081) (-0.83775804,-2.303845) (-0.84299403,-2.2986089) (-0.84823002,-2.2933729) (-0.853466,-2.2881369) (-0.85870199,-2.2829008) (-0.86393798,-2.2776648) (-0.86917397,-2.2724288) (-0.87440996,-2.2671927) (-0.87964594,-2.2619567) (-0.88488193,-2.2567207) (-0.89011792,-2.2514846) (-0.89535391,-2.2462486) (-0.90058989,-2.2410126) (-0.90582588,-2.2357765) (-0.91106187,-2.2305405) (-0.91629786,-2.2253045) (-0.92153385,-2.2200685) (-0.92676983,-2.2148324) (-0.93200582,-2.2095964) (-0.93724181,-2.2043604) (-0.9424778,-2.1991244) (-0.94771378,-2.1938883) (-0.95294977,-2.1886523) (-0.95818576,-2.1834163) (-0.96342175,-2.1781803) (-0.96865773,-2.1729443) (-0.97389372,-2.1677082) (-0.97912971,-2.1624722) (-0.9843657,-2.1572362) (-0.98960169,-2.1520002) (-0.99483767,-2.1467642) (-1.0000737,-2.1415281) (-1.0053096,-2.1362921) (-1.0105456,-2.1310561) (-1.0157816,-2.1258201) (-1.0210176,-2.1205841) (-1.0262536,-2.115348) (-1.0314896,-2.110112) (-1.0367256,-2.104876) (-1.0419616,-2.09964) (-1.0471976,-2.094404) (-1.0524335,-2.089168) (-1.0576695,-2.083932) (-1.0629055,-2.0786959) (-1.0681415,-2.0734599) (-1.0733775,-2.0682239) (-1.0786135,-2.0629879) (-1.0838495,-2.0577519) (-1.0890855,-2.0525159) (-1.0943214,-2.0472799) (-1.0995574,-2.0420439) (-1.1047934,-2.0368078) (-1.1100294,-2.0315718) (-1.1152654,-2.0263358) (-1.1205014,-2.0210998) (-1.1257374,-2.0158638) (-1.1309734,-2.0106278) (-1.1362093,-2.0053918) (-1.1414453,-2.0001558) (-1.1466813,-1.9949198) (-1.1519173,-1.9896838) (-1.1571533,-1.9844478) (-1.1623893,-1.9792118) (-1.1676253,-1.9739757) (-1.1728613,-1.9687397) (-1.1780972,-1.9635037) (-1.1833332,-1.9582677) (-1.1885692,-1.9530317) (-1.1938052,-1.9477957) (-1.1990412,-1.9425597) (-1.2042772,-1.9373237) (-1.2095132,-1.9320877) (-1.2147492,-1.9268517) (-1.2199851,-1.9216157) (-1.2252211,-1.9163797) (-1.2304571,-1.9111437) (-1.2356931,-1.9059077) (-1.2409291,-1.9006717) (-1.2461651,-1.8954357) (-1.2514011,-1.8901997) (-1.2566371,-1.8849637) (-1.261873,-1.8797277) (-1.267109,-1.8744917) (-1.272345,-1.8692557) (-1.277581,-1.8640197) (-1.282817,-1.8587837) (-1.288053,-1.8535477) (-1.293289,-1.8483117) (-1.298525,-1.8430757) (-1.303761,-1.8378397) (-1.3089969,-1.8326037) (-1.3142329,-1.8273677) (-1.3194689,-1.8221317) (-1.3247049,-1.8168957) (-1.3299409,-1.8116597) (-1.3351769,-1.8064237) (-1.3404129,-1.8011877) (-1.3456489,-1.7959517) (-1.3508848,-1.7907157) (-1.3561208,-1.7854797) (-1.3613568,-1.7802437) (-1.3665928,-1.7750077) (-1.3718288,-1.7697717) (-1.3770648,-1.7645357) (-1.3823008,-1.7592997) (-1.3875368,-1.7540637) (-1.3927727,-1.7488277) (-1.3980087,-1.7435917) (-1.4032447,-1.7383557) (-1.4084807,-1.7331197) (-1.4137167,-1.7278837) (-1.4189527,-1.7226478) (-1.4241887,-1.7174118) (-1.4294247,-1.7121758) (-1.4346606,-1.7069398) (-1.4398966,-1.7017038) (-1.4451326,-1.6964678) (-1.4503686,-1.6912318) (-1.4556046,-1.6859958) (-1.4608406,-1.6807598) (-1.4660766,-1.6755238) (-1.4713126,-1.6702878) (-1.4765485,-1.6650518) (-1.4817845,-1.6598158) (-1.4870205,-1.6545798) (-1.4922565,-1.6493439) (-1.4974925,-1.6441079) (-1.5027285,-1.6388719) (-1.5079645,-1.6336359) (-1.5132005,-1.6283999) (-1.5184364,-1.6231639) (-1.5236724,-1.6179279) (-1.5289084,-1.6126919) (-1.5341444,-1.6074559) (-1.5393804,-1.6022199) (-1.5446164,-1.596984) (-1.5498524,-1.591748) (-1.5550884,-1.586512) (-1.5603244,-1.581276) (-1.5655603,-1.57604) (-1.5707963,-1.570804) (-1.5760323,-1.565568) (-1.5812683,-1.560332) (-1.5865043,-1.5550961) (-1.5917403,-1.5498601) (-1.5969763,-1.5446241) (-1.6022123,-1.5393881) (-1.6074482,-1.5341521) (-1.6126842,-1.5289161) (-1.6179202,-1.5236801) (-1.6231562,-1.5184442) (-1.6283922,-1.5132082) (-1.6336282,-1.5079722) (-1.6388642,-1.5027362) (-1.6441002,-1.4975002) (-1.6493361,-1.4922642) (-1.6545721,-1.4870282) (-1.6598081,-1.4817923) (-1.6650441,-1.4765563) (-1.6702801,-1.4713203) (-1.6755161,-1.4660843) (-1.6807521,-1.4608483) (-1.6859881,-1.4556123) (-1.691224,-1.4503764) (-1.69646,-1.4451404) (-1.701696,-1.4399044) (-1.706932,-1.4346684) (-1.712168,-1.4294324) (-1.717404,-1.4241964) (-1.72264,-1.4189605) (-1.727876,-1.4137245) (-1.7331119,-1.4084885) (-1.7383479,-1.4032525) (-1.7435839,-1.3980165) (-1.7488199,-1.3927806) (-1.7540559,-1.3875446) (-1.7592919,-1.3823086) (-1.7645279,-1.3770726) (-1.7697639,-1.3718366) (-1.7749998,-1.3666007) (-1.7802358,-1.3613647) (-1.7854718,-1.3561287) (-1.7907078,-1.3508927) (-1.7959438,-1.3456567) (-1.8011798,-1.3404208) (-1.8064158,-1.3351848) (-1.8116518,-1.3299488) (-1.8168878,-1.3247128) (-1.8221237,-1.3194769) (-1.8273597,-1.3142409) (-1.8325957,-1.3090049) (-1.8378317,-1.3037689) (-1.8430677,-1.2985329) (-1.8483037,-1.293297) (-1.8535397,-1.288061) (-1.8587757,-1.282825) (-1.8640116,-1.277589) (-1.8692476,-1.2723531) (-1.8744836,-1.2671171) (-1.8797196,-1.2618811) (-1.8849556,-1.2566451) (-1.8901916,-1.2514092) (-1.8954276,-1.2461732) (-1.9006636,-1.2409372) (-1.9058995,-1.2357013) (-1.9111355,-1.2304653) (-1.9163715,-1.2252293) (-1.9216075,-1.2199933) (-1.9268435,-1.2147574) (-1.9320795,-1.2095214) (-1.9373155,-1.2042854) (-1.9425515,-1.1990495) (-1.9477874,-1.1938135) (-1.9530234,-1.1885775) (-1.9582594,-1.1833415) (-1.9634954,-1.1781056) (-1.9687314,-1.1728696) (-1.9739674,-1.1676336) (-1.9792034,-1.1623977) (-1.9844394,-1.1571617) (-1.9896753,-1.1519257) (-1.9949113,-1.1466898) (-2.0001473,-1.1414538) (-2.0053833,-1.1362178) (-2.0106193,-1.1309819) (-2.0158553,-1.1257459) (-2.0210913,-1.1205099) (-2.0263273,-1.115274) (-2.0315632,-1.110038) (-2.0367992,-1.104802) (-2.0420352,-1.0995661) (-2.0472712,-1.0943301) (-2.0525072,-1.0890941) (-2.0577432,-1.0838582) (-2.0629792,-1.0786222) (-2.0682152,-1.0733862) (-2.0734512,-1.0681503) (-2.0786871,-1.0629143) (-2.0839231,-1.0576784) (-2.0891591,-1.0524424) (-2.0943951,-1.0472064) (-2.0996311,-1.0419705) (-2.1048671,-1.0367345) (-2.1101031,-1.0314986) (-2.1153391,-1.0262626) (-2.120575,-1.0210266) (-2.125811,-1.0157907) (-2.131047,-1.0105547) (-2.136283,-1.0053188) (-2.141519,-1.0000828) (-2.146755,-0.99484684) (-2.151991,-0.98961089) (-2.157227,-0.98437493) (-2.1624629,-0.97913898) (-2.1676989,-0.97390302) (-2.1729349,-0.96866707) (-2.1781709,-0.96343111) (-2.1834069,-0.95819516) (-2.1886429,-0.95295921) (-2.1938789,-0.94772325) (-2.1991149,-0.9424873) (-2.2043508,-0.93725135) (-2.2095868,-0.9320154) (-2.2148228,-0.92677945) (-2.2200588,-0.9215435) (-2.2252948,-0.91630755) (-2.2305308,-0.9110716) (-2.2357668,-0.90583566) (-2.2410028,-0.90059971) (-2.2462387,-0.89536376) (-2.2514747,-0.89012782) (-2.2567107,-0.88489187) (-2.2619467,-0.87965593) (-2.2671827,-0.87441998) (-2.2724187,-0.86918404) (-2.2776547,-0.86394809) (-2.2828907,-0.85871215) (-2.2881266,-0.85347621) (-2.2933626,-0.84824027) (-2.2985986,-0.84300433) (-2.3038346,-0.83776839) (-2.3090706,-0.83253245) (-2.3143066,-0.82729651) (-2.3195426,-0.82206058) (-2.3247786,-0.81682464) (-2.3300146,-0.81158871) (-2.3352505,-0.80635277) (-2.3404865,-0.80111684) (-2.3457225,-0.7958809) (-2.3509585,-0.79064497) (-2.3561945,-0.78540904) (-2.3614305,-0.78017311) (-2.3666665,-0.77493718) (-2.3719025,-0.76970125) (-2.3771384,-0.76446532) (-2.3823744,-0.7592294) (-2.3876104,-0.75399347) (-2.3928464,-0.74875755) (-2.3980824,-0.74352162) (-2.4033184,-0.7382857) (-2.4085544,-0.73304978) (-2.4137904,-0.72781386) (-2.4190263,-0.72257794) (-2.4242623,-0.71734202) (-2.4294983,-0.71210611) (-2.4347343,-0.70687019) (-2.4399703,-0.70163428) (-2.4452063,-0.69639836) (-2.4504423,-0.69116245) (-2.4556783,-0.68592654) (-2.4609142,-0.68069063) (-2.4661502,-0.67545472) (-2.4713862,-0.67021882) (-2.4766222,-0.66498291) (-2.4818582,-0.65974701) (-2.4870942,-0.6545111) (-2.4923302,-0.6492752) (-2.4975662,-0.6440393) (-2.5028021,-0.63880341) (-2.5080381,-0.63356751) (-2.5132741,-0.62833162) (-2.5185101,-0.62309572) (-2.5237461,-0.61785983) (-2.5289821,-0.61262394) (-2.5342181,-0.60738806) (-2.5394541,-0.60215217) (-2.54469,-0.59691629) (-2.549926,-0.59168041) (-2.555162,-0.58644453) (-2.560398,-0.58120865) (-2.565634,-0.57597277) (-2.57087,-0.5707369) (-2.576106,-0.56550103) (-2.581342,-0.56026516) (-2.586578,-0.5550293) (-2.5918139,-0.54979343) (-2.5970499,-0.54455757) (-2.6022859,-0.53932172) (-2.6075219,-0.53408586) (-2.6127579,-0.52885001) (-2.6179939,-0.52361416) (-2.6232299,-0.51837831) (-2.6284659,-0.51314247) (-2.6337018,-0.50790663) (-2.6389378,-0.50267079) (-2.6441738,-0.49743496) (-2.6494098,-0.49219913) (-2.6546458,-0.4869633) (-2.6598818,-0.48172747) (-2.6651178,-0.47649166) (-2.6703538,-0.47125584) (-2.6755897,-0.46602003) (-2.6808257,-0.46078422) (-2.6860617,-0.45554842) (-2.6912977,-0.45031262) (-2.6965337,-0.44507682) (-2.7017697,-0.43984104) (-2.7070057,-0.43460525) (-2.7122417,-0.42936947) (-2.7174776,-0.4241337) (-2.7227136,-0.41889793) (-2.7279496,-0.41366217) (-2.7331856,-0.40842641) (-2.7384216,-0.40319066) (-2.7436576,-0.39795492) (-2.7488936,-0.39271918) (-2.7541296,-0.38748345) (-2.7593655,-0.38224773) (-2.7646015,-0.37701201) (-2.7698375,-0.3717763) (-2.7750735,-0.3665406) (-2.7803095,-0.36130491) (-2.7855455,-0.35606923) (-2.7907815,-0.35083356) (-2.7960175,-0.3455979) (-2.8012534,-0.34036224) (-2.8064894,-0.3351266) (-2.8117254,-0.32989097) (-2.8169614,-0.32465535) (-2.8221974,-0.31941975) (-2.8274334,-0.31418415) (-2.8326694,-0.30894857) (-2.8379054,-0.30371301) (-2.8431414,-0.29847746) (-2.8483773,-0.29324192) (-2.8536133,-0.28800641) (-2.8588493,-0.28277091) (-2.8640853,-0.27753542) (-2.8693213,-0.27229996) (-2.8745573,-0.26706452) (-2.8797933,-0.2618291) (-2.8850293,-0.25659371) (-2.8902652,-0.25135834) (-2.8955012,-0.24612299) (-2.9007372,-0.24088768) (-2.9059732,-0.23565239) (-2.9112092,-0.23041714) (-2.9164452,-0.22518192) (-2.9216812,-0.21994674) (-2.9269172,-0.2147116) (-2.9321531,-0.2094765) (-2.9373891,-0.20424145) (-2.9426251,-0.19900644) (-2.9478611,-0.19377149) (-2.9530971,-0.1885366) (-2.9583331,-0.18330177) (-2.9635691,-0.17806701) (-2.9688051,-0.17283233) (-2.974041,-0.16759772) (-2.979277,-0.16236321) (-2.984513,-0.15712879) (-2.989749,-0.15189448) (-2.994985,-0.1466603) (-3.000221,-0.14142625) (-3.005457,-0.13619234) (-3.010693,-0.13095861) (-3.0159289,-0.12572506) (-3.0211649,-0.12049172) (-3.0264009,-0.11525863) (-3.0316369,-0.11002581) (-3.0368729,-0.10479331) (-3.0421089,-0.099561178) (-3.0473449,-0.094329473) (-3.0525809,-0.089098273) (-3.0578168,-0.08386767) (-3.0630528,-0.078637786) (-3.0682888,-0.073408772) (-3.0735248,-0.06818083) (-3.0787608,-0.062954227) (-3.0839968,-0.057729325) (-3.0892328,-0.052506634) (-3.0944688,-0.047286885) (-3.0997048,-0.042071173) (-3.1049407,-0.036861213) (-3.1101767,-0.031659844) (-3.1154127,-0.026472131) (-3.1206487,-0.02130805) (-3.1258847,-0.016190229) (-3.1311207,-0.011182367) (-3.1363567,-0.0065420497) (-3.1415927,-0.0039220882)};
\addlegendentry{$p=0$}
\addplot[thick,color=main] coordinates {(-2.586578,0.057843896) (-2.581342,0.094303316) (-2.576106,0.12038277) (-2.57087,0.1419246) (-2.565634,0.16076333) (-2.560398,0.17776042) (-2.555162,0.193402) (-2.549926,0.20799408) (-2.54469,0.22174392) (-2.5394541,0.23479951) (-2.5342181,0.24727081) (-2.5289821,0.25924214) (-2.5237461,0.27077978) (-2.5185101,0.28193697) (-2.5132741,0.29275719) (-2.5080381,0.30327648) (-2.5028021,0.31352511) (-2.4975662,0.32352878) (-2.4923302,0.33330953) (-2.4870942,0.34288642) (-2.4818582,0.35227605) (-2.4766222,0.36149299) (-2.4713862,0.37055012) (-2.4661502,0.37945885) (-2.4609142,0.38822939) (-2.4556783,0.39687087) (-2.4504423,0.40539153) (-2.4452063,0.41379881) (-2.4399703,0.42209948) (-2.4347343,0.43029968) (-2.4294983,0.43840503) (-2.4242623,0.44642069) (-2.4190263,0.45435138) (-2.4137904,0.46220145) (-2.4085544,0.46997493) (-2.4033184,0.47767555) (-2.3980824,0.48530674) (-2.3928464,0.49287173) (-2.3876104,0.50037348) (-2.3823744,0.5078148) (-2.3771384,0.51519828) (-2.3719025,0.52252636) (-2.3666665,0.52980132) (-2.3614305,0.53702529) (-2.3561945,0.5442003) (-2.3509585,0.55132824) (-2.3457225,0.55841089) (-2.3404865,0.56544994) (-2.3352505,0.57244698) (-2.3300146,0.57940351) (-2.3247786,0.58632097) (-2.3195426,0.59320069) (-2.3143066,0.60004397) (-2.3090706,0.60685202) (-2.3038346,0.61362599) (-2.2985986,0.62036699) (-2.2933626,0.62707605) (-2.2881266,0.63375418) (-2.2828907,0.64040234) (-2.2776547,0.64702142) (-2.2724187,0.65361229) (-2.2671827,0.66017578) (-2.2619467,0.66671268) (-2.2567107,0.67322376) (-2.2514747,0.67970973) (-2.2462387,0.68617129) (-2.2410028,0.69260911) (-2.2357668,0.69902383) (-2.2305308,0.70541605) (-2.2252948,0.71178638) (-2.2200588,0.71813537) (-2.2148228,0.72446357) (-2.2095868,0.7307715) (-2.2043508,0.73705967) (-2.1991149,0.74332856) (-2.1938789,0.74957863) (-2.1886429,0.75581034) (-2.1834069,0.76202412) (-2.1781709,0.76822039) (-2.1729349,0.77439955) (-2.1676989,0.78056199) (-2.1624629,0.78670808) (-2.157227,0.79283819) (-2.151991,0.79895267) (-2.146755,0.80505186) (-2.141519,0.81113609) (-2.136283,0.81720566) (-2.131047,0.82326089) (-2.125811,0.82930208) (-2.120575,0.83532951) (-2.1153391,0.84134346) (-2.1101031,0.8473442) (-2.1048671,0.85333199) (-2.0996311,0.85930708) (-2.0943951,0.86526971) (-2.0891591,0.87122013) (-2.0839231,0.87715857) (-2.0786871,0.88308525) (-2.0734512,0.88900038) (-2.0682152,0.89490418) (-2.0629792,0.90079685) (-2.0577432,0.90667858) (-2.0525072,0.91254958) (-2.0472712,0.91841003) (-2.0420352,0.92426012) (-2.0367992,0.9301) (-2.0315632,0.93592987) (-2.0263273,0.94174989) (-2.0210913,0.94756022) (-2.0158553,0.95336101) (-2.0106193,0.95915242) (-2.0053833,0.96493461) (-2.0001473,0.97070771) (-1.9949113,0.97647187) (-1.9896753,0.98222723) (-1.9844394,0.98797392) (-1.9792034,0.99371206) (-1.9739674,0.9994418) (-1.9687314,1.0051632) (-1.9634954,1.0108765) (-1.9582594,1.0165818) (-1.9530234,1.022279) (-1.9477874,1.0279685) (-1.9425515,1.0336503) (-1.9373155,1.0393244) (-1.9320795,1.0449911) (-1.9268435,1.0506503) (-1.9216075,1.0563022) (-1.9163715,1.0619469) (-1.9111355,1.0675845) (-1.9058995,1.0732151) (-1.9006636,1.0788387) (-1.8954276,1.0844555) (-1.8901916,1.0900655) (-1.8849556,1.0956688) (-1.8797196,1.1012656) (-1.8744836,1.1068558) (-1.8692476,1.1124395) (-1.8640116,1.1180169) (-1.8587757,1.1235879) (-1.8535397,1.1291527) (-1.8483037,1.1347114) (-1.8430677,1.140264) (-1.8378317,1.1458105) (-1.8325957,1.1513511) (-1.8273597,1.1568857) (-1.8221237,1.1624145) (-1.8168878,1.1679376) (-1.8116518,1.1734549) (-1.8064158,1.1789665) (-1.8011798,1.1844725) (-1.7959438,1.189973) (-1.7907078,1.195468) (-1.7854718,1.2009575) (-1.7802358,1.2064415) (-1.7749998,1.2119203) (-1.7697639,1.2173937) (-1.7645279,1.2228618) (-1.7592919,1.2283248) (-1.7540559,1.2337825) (-1.7488199,1.2392352) (-1.7435839,1.2446827) (-1.7383479,1.2501252) (-1.7331119,1.2555627) (-1.727876,1.2609952) (-1.72264,1.2664227) (-1.717404,1.2718454) (-1.712168,1.2772632) (-1.706932,1.2826762) (-1.701696,1.2880843) (-1.69646,1.2934877) (-1.691224,1.2988864) (-1.6859881,1.3042804) (-1.6807521,1.3096697) (-1.6755161,1.3150544) (-1.6702801,1.3204345) (-1.6650441,1.3258099) (-1.6598081,1.3311808) (-1.6545721,1.3365472) (-1.6493361,1.3419091) (-1.6441002,1.3472665) (-1.6388642,1.3526194) (-1.6336282,1.3579679) (-1.6283922,1.363312) (-1.6231562,1.3686517) (-1.6179202,1.373987) (-1.6126842,1.3793179) (-1.6074482,1.3846446) (-1.6022123,1.3899669) (-1.5969763,1.3952849) (-1.5917403,1.4005986) (-1.5865043,1.405908) (-1.5812683,1.4112132) (-1.5760323,1.4165142) (-1.5707963,1.421811) (-1.5655603,1.4271035) (-1.5603244,1.4323918) (-1.5550884,1.437676) (-1.5498524,1.442956) (-1.5446164,1.4482318) (-1.5393804,1.4535035) (-1.5341444,1.4587711) (-1.5289084,1.4640345) (-1.5236724,1.4692938) (-1.5184364,1.4745489) (-1.5132005,1.4798) (-1.5079645,1.485047) (-1.5027285,1.4902899) (-1.4974925,1.4955287) (-1.4922565,1.5007634) (-1.4870205,1.505994) (-1.4817845,1.5112206) (-1.4765485,1.5164431) (-1.4713126,1.5216615) (-1.4660766,1.5268759) (-1.4608406,1.5320862) (-1.4556046,1.5372924) (-1.4503686,1.5424946) (-1.4451326,1.5476927) (-1.4398966,1.5528867) (-1.4346606,1.5580767) (-1.4294247,1.5632627) (-1.4241887,1.5684445) (-1.4189527,1.5736223) (-1.4137167,1.5787961) (-1.4084807,1.5839657) (-1.4032447,1.5891313) (-1.3980087,1.5942928) (-1.3927727,1.5994502) (-1.3875368,1.6046035) (-1.3823008,1.6097527) (-1.3770648,1.6148978) (-1.3718288,1.6200388) (-1.3665928,1.6251756) (-1.3613568,1.6303083) (-1.3561208,1.6354369) (-1.3508848,1.6405613) (-1.3456489,1.6456816) (-1.3404129,1.6507977) (-1.3351769,1.6559096) (-1.3299409,1.6610173) (-1.3247049,1.6661208) (-1.3194689,1.6712201) (-1.3142329,1.6763152) (-1.3089969,1.681406) (-1.303761,1.6864925) (-1.298525,1.6915748) (-1.293289,1.6966527) (-1.288053,1.7017264) (-1.282817,1.7067957) (-1.277581,1.7118607) (-1.272345,1.7169213) (-1.267109,1.7219775) (-1.261873,1.7270294) (-1.2566371,1.7320768) (-1.2514011,1.7371198) (-1.2461651,1.7421583) (-1.2409291,1.7471923) (-1.2356931,1.7522218) (-1.2304571,1.7572468) (-1.2252211,1.7622672) (-1.2199851,1.7672831) (-1.2147492,1.7722943) (-1.2095132,1.7773009) (-1.2042772,1.7823029) (-1.1990412,1.7873002) (-1.1938052,1.7922928) (-1.1885692,1.7972806) (-1.1833332,1.8022637) (-1.1780972,1.8072419) (-1.1728613,1.8122154) (-1.1676253,1.817184) (-1.1623893,1.8221477) (-1.1571533,1.8271065) (-1.1519173,1.8320603) (-1.1466813,1.8370091) (-1.1414453,1.8419529) (-1.1362093,1.8468917) (-1.1309734,1.8518254) (-1.1257374,1.8567539) (-1.1205014,1.8616773) (-1.1152654,1.8665954) (-1.1100294,1.8715084) (-1.1047934,1.876416) (-1.0995574,1.8813183) (-1.0943214,1.8862152) (-1.0890855,1.8911068) (-1.0838495,1.8959928) (-1.0786135,1.9008734) (-1.0733775,1.9057484) (-1.0681415,1.9106179) (-1.0629055,1.9154817) (-1.0576695,1.9203398) (-1.0524335,1.9251922) (-1.0471976,1.9300388) (-1.0419616,1.9348796) (-1.0367256,1.9397145) (-1.0314896,1.9445435) (-1.0262536,1.9493664) (-1.0210176,1.9541834) (-1.0157816,1.9589942) (-1.0105456,1.9637989) (-1.0053096,1.9685973) (-1.0000737,1.9733895) (-0.99483767,1.9781754) (-0.98960169,1.9829548) (-0.9843657,1.9877278) (-0.97912971,1.9924943) (-0.97389372,1.9972542) (-0.96865773,2.0020075) (-0.96342175,2.0067541) (-0.95818576,2.0114938) (-0.95294977,2.0162268) (-0.94771378,2.0209528) (-0.9424778,2.0256718) (-0.93724181,2.0303837) (-0.93200582,2.0350885) (-0.92676983,2.0397861) (-0.92153385,2.0444765) (-0.91629786,2.0491594) (-0.91106187,2.0538349) (-0.90582588,2.0585029) (-0.90058989,2.0631632) (-0.89535391,2.0678159) (-0.89011792,2.0724608) (-0.88488193,2.0770978) (-0.87964594,2.0817269) (-0.87440996,2.0863479) (-0.86917397,2.0909608) (-0.86393798,2.0955655) (-0.85870199,2.1001618) (-0.853466,2.1047498) (-0.84823002,2.1093292) (-0.84299403,2.1139) (-0.83775804,2.1184621) (-0.83252205,2.1230153) (-0.82728607,2.1275597) (-0.82205008,2.132095) (-0.81681409,2.1366212) (-0.8115781,2.1411381) (-0.80634211,2.1456457) (-0.80110613,2.1501438) (-0.79587014,2.1546323) (-0.79063415,2.1591111) (-0.78539816,2.16358) (-0.78016218,2.1680391) (-0.77492619,2.172488) (-0.7696902,2.1769267) (-0.76445421,2.1813552) (-0.75921822,2.1857731) (-0.75398224,2.1901805) (-0.74874625,2.1945771) (-0.74351026,2.1989629) (-0.73827427,2.2033376) (-0.73303829,2.2077012) (-0.7278023,2.2120535) (-0.72256631,2.2163944) (-0.71733032,2.2207237) (-0.71209433,2.2250412) (-0.70685835,2.2293469) (-0.70162236,2.2336405) (-0.69638637,2.2379218) (-0.69115038,2.2421908) (-0.6859144,2.2464473) (-0.68067841,2.250691) (-0.67544242,2.2549219) (-0.67020643,2.2591397) (-0.66497045,2.2633443) (-0.65973446,2.2675354) (-0.65449847,2.271713) (-0.64926248,2.2758768) (-0.64402649,2.2800266) (-0.63879051,2.2841622) (-0.63355452,2.2882835) (-0.62831853,2.2923903) (-0.62308254,2.2964823) (-0.61784656,2.3005593) (-0.61261057,2.3046211) (-0.60737458,2.3086676) (-0.60213859,2.3126985) (-0.5969026,2.3167135) (-0.59166662,2.3207126) (-0.58643063,2.3246953) (-0.58119464,2.3286616) (-0.57595865,2.3326111) (-0.57072267,2.3365437) (-0.56548668,2.3404591) (-0.56025069,2.3443571) (-0.5550147,2.3482373) (-0.54977871,2.3520996) (-0.54454273,2.3559437) (-0.53930674,2.3597694) (-0.53407075,2.3635763) (-0.52883476,2.3673643) (-0.52359878,2.371133) (-0.51836279,2.3748821) (-0.5131268,2.3786115) (-0.50789081,2.3823207) (-0.50265482,2.3860096) (-0.49741884,2.3896778) (-0.49218285,2.393325) (-0.48694686,2.396951) (-0.48171087,2.4005554) (-0.47647489,2.4041379) (-0.4712389,2.4076983) (-0.46600291,2.4112361) (-0.46076692,2.4147511) (-0.45553093,2.4182429) (-0.45029495,2.4217113) (-0.44505896,2.4251559) (-0.43982297,2.4285763) (-0.43458698,2.4319722) (-0.429351,2.4353432) (-0.42411501,2.4386891) (-0.41887902,2.4420095) (-0.41364303,2.4453039) (-0.40840704,2.4485721) (-0.40317106,2.4518137) (-0.39793507,2.4550282) (-0.39269908,2.4582154) (-0.38746309,2.4613749) (-0.38222711,2.4645062) (-0.37699112,2.467609) (-0.37175513,2.4706829) (-0.36651914,2.4737275) (-0.36128316,2.4767425) (-0.35604717,2.4797273) (-0.35081118,2.4826816) (-0.34557519,2.4856051) (-0.3403392,2.4884972) (-0.33510322,2.4913576) (-0.32986723,2.4941859) (-0.32463124,2.4969817) (-0.31939525,2.4997445) (-0.31415927,2.5024739) (-0.30892328,2.5051694) (-0.30368729,2.5078308) (-0.2984513,2.5104575) (-0.29321531,2.5130491) (-0.28797933,2.5156051) (-0.28274334,2.5181252) (-0.27750735,2.5206089) (-0.27227136,2.5230558) (-0.26703538,2.5254655) (-0.26179939,2.5278375) (-0.2565634,2.5301713) (-0.25132741,2.5324666) (-0.24609142,2.5347229) (-0.24085544,2.5369398) (-0.23561945,2.5391168) (-0.23038346,2.5412535) (-0.22514747,2.5433496) (-0.21991149,2.5454045) (-0.2146755,2.5474179) (-0.20943951,2.5493893) (-0.20420352,2.5513183) (-0.19896753,2.5532045) (-0.19373155,2.5550475) (-0.18849556,2.5568469) (-0.18325957,2.5586022) (-0.17802358,2.5603132) (-0.1727876,2.5619794) (-0.16755161,2.5636003) (-0.16231562,2.5651758) (-0.15707963,2.5667052) (-0.15184364,2.5681884) (-0.14660766,2.569625) (-0.14137167,2.5710145) (-0.13613568,2.5723567) (-0.13089969,2.5736512) (-0.12566371,2.5748977) (-0.12042772,2.5760959) (-0.11519173,2.5772455) (-0.10995574,2.5783461) (-0.10471976,2.5793976) (-0.099483767,2.5803996) (-0.09424778,2.5813518) (-0.089011792,2.582254) (-0.083775804,2.5831061) (-0.078539816,2.5839076) (-0.073303829,2.5846585) (-0.068067841,2.5853585) (-0.062831853,2.5860074) (-0.057595865,2.5866051) (-0.052359878,2.5871514) (-0.04712389,2.5876461) (-0.041887902,2.5880891) (-0.036651914,2.5884802) (-0.031415927,2.5888195) (-0.026179939,2.5891067) (-0.020943951,2.5893418) (-0.015707963,2.5895247) (-0.010471976,2.5896554) (-0.0052359878,2.5897339) (-4.4408921e-16,2.58976) (0.0052359878,2.5897339) (0.010471976,2.5896554) (0.015707963,2.5895247) (0.020943951,2.5893418) (0.026179939,2.5891067) (0.031415927,2.5888195) (0.036651914,2.5884802) (0.041887902,2.5880891) (0.04712389,2.5876461) (0.052359878,2.5871514) (0.057595865,2.5866051) (0.062831853,2.5860074) (0.068067841,2.5853585) (0.073303829,2.5846585) (0.078539816,2.5839076) (0.083775804,2.5831061) (0.089011792,2.582254) (0.09424778,2.5813518) (0.099483767,2.5803996) (0.10471976,2.5793976) (0.10995574,2.5783461) (0.11519173,2.5772455) (0.12042772,2.5760959) (0.12566371,2.5748977) (0.13089969,2.5736512) (0.13613568,2.5723567) (0.14137167,2.5710145) (0.14660766,2.569625) (0.15184364,2.5681884) (0.15707963,2.5667052) (0.16231562,2.5651758) (0.16755161,2.5636003) (0.1727876,2.5619794) (0.17802358,2.5603132) (0.18325957,2.5586022) (0.18849556,2.5568469) (0.19373155,2.5550475) (0.19896753,2.5532045) (0.20420352,2.5513183) (0.20943951,2.5493893) (0.2146755,2.5474179) (0.21991149,2.5454045) (0.22514747,2.5433496) (0.23038346,2.5412535) (0.23561945,2.5391168) (0.24085544,2.5369398) (0.24609142,2.5347229) (0.25132741,2.5324666) (0.2565634,2.5301713) (0.26179939,2.5278375) (0.26703538,2.5254655) (0.27227136,2.5230558) (0.27750735,2.5206089) (0.28274334,2.5181252) (0.28797933,2.5156051) (0.29321531,2.5130491) (0.2984513,2.5104575) (0.30368729,2.5078308) (0.30892328,2.5051694) (0.31415927,2.5024739) (0.31939525,2.4997445) (0.32463124,2.4969817) (0.32986723,2.4941859) (0.33510322,2.4913576) (0.3403392,2.4884972) (0.34557519,2.4856051) (0.35081118,2.4826816) (0.35604717,2.4797273) (0.36128316,2.4767425) (0.36651914,2.4737275) (0.37175513,2.4706829) (0.37699112,2.467609) (0.38222711,2.4645062) (0.38746309,2.4613749) (0.39269908,2.4582154) (0.39793507,2.4550282) (0.40317106,2.4518137) (0.40840704,2.4485721) (0.41364303,2.4453039) (0.41887902,2.4420095) (0.42411501,2.4386891) (0.429351,2.4353432) (0.43458698,2.4319722) (0.43982297,2.4285763) (0.44505896,2.4251559) (0.45029495,2.4217113) (0.45553093,2.4182429) (0.46076692,2.4147511) (0.46600291,2.4112361) (0.4712389,2.4076983) (0.47647489,2.4041379) (0.48171087,2.4005554) (0.48694686,2.396951) (0.49218285,2.393325) (0.49741884,2.3896778) (0.50265482,2.3860096) (0.50789081,2.3823207) (0.5131268,2.3786115) (0.51836279,2.3748821) (0.52359878,2.371133) (0.52883476,2.3673643) (0.53407075,2.3635763) (0.53930674,2.3597694) (0.54454273,2.3559437) (0.54977871,2.3520996) (0.5550147,2.3482373) (0.56025069,2.3443571) (0.56548668,2.3404591) (0.57072267,2.3365437) (0.57595865,2.3326111) (0.58119464,2.3286616) (0.58643063,2.3246953) (0.59166662,2.3207126) (0.5969026,2.3167135) (0.60213859,2.3126985) (0.60737458,2.3086676) (0.61261057,2.3046211) (0.61784656,2.3005593) (0.62308254,2.2964823) (0.62831853,2.2923903) (0.63355452,2.2882835) (0.63879051,2.2841622) (0.64402649,2.2800266) (0.64926248,2.2758768) (0.65449847,2.271713) (0.65973446,2.2675354) (0.66497045,2.2633443) (0.67020643,2.2591397) (0.67544242,2.2549219) (0.68067841,2.250691) (0.6859144,2.2464473) (0.69115038,2.2421908) (0.69638637,2.2379218) (0.70162236,2.2336405) (0.70685835,2.2293469) (0.71209433,2.2250412) (0.71733032,2.2207237) (0.72256631,2.2163944) (0.7278023,2.2120535) (0.73303829,2.2077012) (0.73827427,2.2033376) (0.74351026,2.1989629) (0.74874625,2.1945771) (0.75398224,2.1901805) (0.75921822,2.1857731) (0.76445421,2.1813552) (0.7696902,2.1769267) (0.77492619,2.172488) (0.78016218,2.1680391) (0.78539816,2.16358) (0.79063415,2.1591111) (0.79587014,2.1546323) (0.80110613,2.1501438) (0.80634211,2.1456457) (0.8115781,2.1411381) (0.81681409,2.1366212) (0.82205008,2.132095) (0.82728607,2.1275597) (0.83252205,2.1230153) (0.83775804,2.1184621) (0.84299403,2.1139) (0.84823002,2.1093292) (0.853466,2.1047498) (0.85870199,2.1001618) (0.86393798,2.0955655) (0.86917397,2.0909608) (0.87440996,2.0863479) (0.87964594,2.0817269) (0.88488193,2.0770978) (0.89011792,2.0724608) (0.89535391,2.0678159) (0.90058989,2.0631632) (0.90582588,2.0585029) (0.91106187,2.0538349) (0.91629786,2.0491594) (0.92153385,2.0444765) (0.92676983,2.0397861) (0.93200582,2.0350885) (0.93724181,2.0303837) (0.9424778,2.0256718) (0.94771378,2.0209528) (0.95294977,2.0162268) (0.95818576,2.0114938) (0.96342175,2.0067541) (0.96865773,2.0020075) (0.97389372,1.9972542) (0.97912971,1.9924943) (0.9843657,1.9877278) (0.98960169,1.9829548) (0.99483767,1.9781754) (1.0000737,1.9733895) (1.0053096,1.9685973) (1.0105456,1.9637989) (1.0157816,1.9589942) (1.0210176,1.9541834) (1.0262536,1.9493664) (1.0314896,1.9445435) (1.0367256,1.9397145) (1.0419616,1.9348796) (1.0471976,1.9300388) (1.0524335,1.9251922) (1.0576695,1.9203398) (1.0629055,1.9154817) (1.0681415,1.9106179) (1.0733775,1.9057484) (1.0786135,1.9008734) (1.0838495,1.8959928) (1.0890855,1.8911068) (1.0943214,1.8862152) (1.0995574,1.8813183) (1.1047934,1.876416) (1.1100294,1.8715084) (1.1152654,1.8665954) (1.1205014,1.8616773) (1.1257374,1.8567539) (1.1309734,1.8518254) (1.1362093,1.8468917) (1.1414453,1.8419529) (1.1466813,1.8370091) (1.1519173,1.8320603) (1.1571533,1.8271065) (1.1623893,1.8221477) (1.1676253,1.817184) (1.1728613,1.8122154) (1.1780972,1.8072419) (1.1833332,1.8022637) (1.1885692,1.7972806) (1.1938052,1.7922928) (1.1990412,1.7873002) (1.2042772,1.7823029) (1.2095132,1.7773009) (1.2147492,1.7722943) (1.2199851,1.7672831) (1.2252211,1.7622672) (1.2304571,1.7572468) (1.2356931,1.7522218) (1.2409291,1.7471923) (1.2461651,1.7421583) (1.2514011,1.7371198) (1.2566371,1.7320768) (1.261873,1.7270294) (1.267109,1.7219775) (1.272345,1.7169213) (1.277581,1.7118607) (1.282817,1.7067957) (1.288053,1.7017264) (1.293289,1.6966527) (1.298525,1.6915748) (1.303761,1.6864925) (1.3089969,1.681406) (1.3142329,1.6763152) (1.3194689,1.6712201) (1.3247049,1.6661208) (1.3299409,1.6610173) (1.3351769,1.6559096) (1.3404129,1.6507977) (1.3456489,1.6456816) (1.3508848,1.6405613) (1.3561208,1.6354369) (1.3613568,1.6303083) (1.3665928,1.6251756) (1.3718288,1.6200388) (1.3770648,1.6148978) (1.3823008,1.6097527) (1.3875368,1.6046035) (1.3927727,1.5994502) (1.3980087,1.5942928) (1.4032447,1.5891313) (1.4084807,1.5839657) (1.4137167,1.5787961) (1.4189527,1.5736223) (1.4241887,1.5684445) (1.4294247,1.5632627) (1.4346606,1.5580767) (1.4398966,1.5528867) (1.4451326,1.5476927) (1.4503686,1.5424946) (1.4556046,1.5372924) (1.4608406,1.5320862) (1.4660766,1.5268759) (1.4713126,1.5216615) (1.4765485,1.5164431) (1.4817845,1.5112206) (1.4870205,1.505994) (1.4922565,1.5007634) (1.4974925,1.4955287) (1.5027285,1.4902899) (1.5079645,1.485047) (1.5132005,1.4798) (1.5184364,1.4745489) (1.5236724,1.4692938) (1.5289084,1.4640345) (1.5341444,1.4587711) (1.5393804,1.4535035) (1.5446164,1.4482318) (1.5498524,1.442956) (1.5550884,1.437676) (1.5603244,1.4323918) (1.5655603,1.4271035) (1.5707963,1.421811) (1.5760323,1.4165142) (1.5812683,1.4112132) (1.5865043,1.405908) (1.5917403,1.4005986) (1.5969763,1.3952849) (1.6022123,1.3899669) (1.6074482,1.3846446) (1.6126842,1.3793179) (1.6179202,1.373987) (1.6231562,1.3686517) (1.6283922,1.363312) (1.6336282,1.3579679) (1.6388642,1.3526194) (1.6441002,1.3472665) (1.6493361,1.3419091) (1.6545721,1.3365472) (1.6598081,1.3311808) (1.6650441,1.3258099) (1.6702801,1.3204345) (1.6755161,1.3150544) (1.6807521,1.3096697) (1.6859881,1.3042804) (1.691224,1.2988864) (1.69646,1.2934877) (1.701696,1.2880843) (1.706932,1.2826762) (1.712168,1.2772632) (1.717404,1.2718454) (1.72264,1.2664227) (1.727876,1.2609952) (1.7331119,1.2555627) (1.7383479,1.2501252) (1.7435839,1.2446827) (1.7488199,1.2392352) (1.7540559,1.2337825) (1.7592919,1.2283248) (1.7645279,1.2228618) (1.7697639,1.2173937) (1.7749998,1.2119203) (1.7802358,1.2064415) (1.7854718,1.2009575) (1.7907078,1.195468) (1.7959438,1.189973) (1.8011798,1.1844725) (1.8064158,1.1789665) (1.8116518,1.1734549) (1.8168878,1.1679376) (1.8221237,1.1624145) (1.8273597,1.1568857) (1.8325957,1.1513511) (1.8378317,1.1458105) (1.8430677,1.140264) (1.8483037,1.1347114) (1.8535397,1.1291527) (1.8587757,1.1235879) (1.8640116,1.1180169) (1.8692476,1.1124395) (1.8744836,1.1068558) (1.8797196,1.1012656) (1.8849556,1.0956688) (1.8901916,1.0900655) (1.8954276,1.0844555) (1.9006636,1.0788387) (1.9058995,1.0732151) (1.9111355,1.0675845) (1.9163715,1.0619469) (1.9216075,1.0563022) (1.9268435,1.0506503) (1.9320795,1.0449911) (1.9373155,1.0393244) (1.9425515,1.0336503) (1.9477874,1.0279685) (1.9530234,1.022279) (1.9582594,1.0165818) (1.9634954,1.0108765) (1.9687314,1.0051632) (1.9739674,0.9994418) (1.9792034,0.99371206) (1.9844394,0.98797392) (1.9896753,0.98222723) (1.9949113,0.97647187) (2.0001473,0.97070771) (2.0053833,0.96493461) (2.0106193,0.95915242) (2.0158553,0.95336101) (2.0210913,0.94756022) (2.0263273,0.94174989) (2.0315632,0.93592987) (2.0367992,0.9301) (2.0420352,0.92426012) (2.0472712,0.91841003) (2.0525072,0.91254958) (2.0577432,0.90667858) (2.0629792,0.90079685) (2.0682152,0.89490418) (2.0734512,0.88900038) (2.0786871,0.88308525) (2.0839231,0.87715857) (2.0891591,0.87122013) (2.0943951,0.86526971) (2.0996311,0.85930708) (2.1048671,0.85333199) (2.1101031,0.8473442) (2.1153391,0.84134346) (2.120575,0.83532951) (2.125811,0.82930208) (2.131047,0.82326089) (2.136283,0.81720566) (2.141519,0.81113609) (2.146755,0.80505186) (2.151991,0.79895267) (2.157227,0.79283819) (2.1624629,0.78670808) (2.1676989,0.78056199) (2.1729349,0.77439955) (2.1781709,0.76822039) (2.1834069,0.76202412) (2.1886429,0.75581034) (2.1938789,0.74957863) (2.1991149,0.74332856) (2.2043508,0.73705967) (2.2095868,0.7307715) (2.2148228,0.72446357) (2.2200588,0.71813537) (2.2252948,0.71178638) (2.2305308,0.70541605) (2.2357668,0.69902383) (2.2410028,0.69260911) (2.2462387,0.68617129) (2.2514747,0.67970973) (2.2567107,0.67322376) (2.2619467,0.66671268) (2.2671827,0.66017578) (2.2724187,0.65361229) (2.2776547,0.64702142) (2.2828907,0.64040234) (2.2881266,0.63375418) (2.2933626,0.62707605) (2.2985986,0.62036699) (2.3038346,0.61362599) (2.3090706,0.60685202) (2.3143066,0.60004397) (2.3195426,0.59320069) (2.3247786,0.58632097) (2.3300146,0.57940351) (2.3352505,0.57244698) (2.3404865,0.56544994) (2.3457225,0.55841089) (2.3509585,0.55132824) (2.3561945,0.5442003) (2.3614305,0.53702529) (2.3666665,0.52980132) (2.3719025,0.52252636) (2.3771384,0.51519828) (2.3823744,0.5078148) (2.3876104,0.50037348) (2.3928464,0.49287173) (2.3980824,0.48530674) (2.4033184,0.47767555) (2.4085544,0.46997493) (2.4137904,0.46220145) (2.4190263,0.45435138) (2.4242623,0.44642069) (2.4294983,0.43840503) (2.4347343,0.43029968) (2.4399703,0.42209948) (2.4452063,0.41379881) (2.4504423,0.40539153) (2.4556783,0.39687087) (2.4609142,0.38822939) (2.4661502,0.37945885) (2.4713862,0.37055012) (2.4766222,0.36149299) (2.4818582,0.35227605) (2.4870942,0.34288642) (2.4923302,0.33330953) (2.4975662,0.32352878) (2.5028021,0.31352511) (2.5080381,0.30327648) (2.5132741,0.29275719) (2.5185101,0.28193697) (2.5237461,0.27077978) (2.5289821,0.25924214) (2.5342181,0.24727081) (2.5394541,0.23479951) (2.54469,0.22174392) (2.549926,0.20799408) (2.555162,0.193402) (2.560398,0.17776042) (2.565634,0.16076333) (2.57087,0.1419246) (2.576106,0.12038277) (2.581342,0.094303316) (2.586578,0.057843896) (2.586578,-0.057843896) (2.581342,-0.094303316) (2.576106,-0.12038277) (2.57087,-0.1419246) (2.565634,-0.16076333) (2.560398,-0.17776042) (2.555162,-0.193402) (2.549926,-0.20799408) (2.54469,-0.22174392) (2.5394541,-0.23479951) (2.5342181,-0.24727081) (2.5289821,-0.25924214) (2.5237461,-0.27077978) (2.5185101,-0.28193697) (2.5132741,-0.29275719) (2.5080381,-0.30327648) (2.5028021,-0.31352511) (2.4975662,-0.32352878) (2.4923302,-0.33330953) (2.4870942,-0.34288642) (2.4818582,-0.35227605) (2.4766222,-0.36149299) (2.4713862,-0.37055012) (2.4661502,-0.37945885) (2.4609142,-0.38822939) (2.4556783,-0.39687087) (2.4504423,-0.40539153) (2.4452063,-0.41379881) (2.4399703,-0.42209948) (2.4347343,-0.43029968) (2.4294983,-0.43840503) (2.4242623,-0.44642069) (2.4190263,-0.45435138) (2.4137904,-0.46220145) (2.4085544,-0.46997493) (2.4033184,-0.47767555) (2.3980824,-0.48530674) (2.3928464,-0.49287173) (2.3876104,-0.50037348) (2.3823744,-0.5078148) (2.3771384,-0.51519828) (2.3719025,-0.52252636) (2.3666665,-0.52980132) (2.3614305,-0.53702529) (2.3561945,-0.5442003) (2.3509585,-0.55132824) (2.3457225,-0.55841089) (2.3404865,-0.56544994) (2.3352505,-0.57244698) (2.3300146,-0.57940351) (2.3247786,-0.58632097) (2.3195426,-0.59320069) (2.3143066,-0.60004397) (2.3090706,-0.60685202) (2.3038346,-0.61362599) (2.2985986,-0.62036699) (2.2933626,-0.62707605) (2.2881266,-0.63375418) (2.2828907,-0.64040234) (2.2776547,-0.64702142) (2.2724187,-0.65361229) (2.2671827,-0.66017578) (2.2619467,-0.66671268) (2.2567107,-0.67322376) (2.2514747,-0.67970973) (2.2462387,-0.68617129) (2.2410028,-0.69260911) (2.2357668,-0.69902383) (2.2305308,-0.70541605) (2.2252948,-0.71178638) (2.2200588,-0.71813537) (2.2148228,-0.72446357) (2.2095868,-0.7307715) (2.2043508,-0.73705967) (2.1991149,-0.74332856) (2.1938789,-0.74957863) (2.1886429,-0.75581034) (2.1834069,-0.76202412) (2.1781709,-0.76822039) (2.1729349,-0.77439955) (2.1676989,-0.78056199) (2.1624629,-0.78670808) (2.157227,-0.79283819) (2.151991,-0.79895267) (2.146755,-0.80505186) (2.141519,-0.81113609) (2.136283,-0.81720566) (2.131047,-0.82326089) (2.125811,-0.82930208) (2.120575,-0.83532951) (2.1153391,-0.84134346) (2.1101031,-0.8473442) (2.1048671,-0.85333199) (2.0996311,-0.85930708) (2.0943951,-0.86526971) (2.0891591,-0.87122013) (2.0839231,-0.87715857) (2.0786871,-0.88308525) (2.0734512,-0.88900038) (2.0682152,-0.89490418) (2.0629792,-0.90079685) (2.0577432,-0.90667858) (2.0525072,-0.91254958) (2.0472712,-0.91841003) (2.0420352,-0.92426012) (2.0367992,-0.9301) (2.0315632,-0.93592987) (2.0263273,-0.94174989) (2.0210913,-0.94756022) (2.0158553,-0.95336101) (2.0106193,-0.95915242) (2.0053833,-0.96493461) (2.0001473,-0.97070771) (1.9949113,-0.97647187) (1.9896753,-0.98222723) (1.9844394,-0.98797392) (1.9792034,-0.99371206) (1.9739674,-0.9994418) (1.9687314,-1.0051632) (1.9634954,-1.0108765) (1.9582594,-1.0165818) (1.9530234,-1.022279) (1.9477874,-1.0279685) (1.9425515,-1.0336503) (1.9373155,-1.0393244) (1.9320795,-1.0449911) (1.9268435,-1.0506503) (1.9216075,-1.0563022) (1.9163715,-1.0619469) (1.9111355,-1.0675845) (1.9058995,-1.0732151) (1.9006636,-1.0788387) (1.8954276,-1.0844555) (1.8901916,-1.0900655) (1.8849556,-1.0956688) (1.8797196,-1.1012656) (1.8744836,-1.1068558) (1.8692476,-1.1124395) (1.8640116,-1.1180169) (1.8587757,-1.1235879) (1.8535397,-1.1291527) (1.8483037,-1.1347114) (1.8430677,-1.140264) (1.8378317,-1.1458105) (1.8325957,-1.1513511) (1.8273597,-1.1568857) (1.8221237,-1.1624145) (1.8168878,-1.1679376) (1.8116518,-1.1734549) (1.8064158,-1.1789665) (1.8011798,-1.1844725) (1.7959438,-1.189973) (1.7907078,-1.195468) (1.7854718,-1.2009575) (1.7802358,-1.2064415) (1.7749998,-1.2119203) (1.7697639,-1.2173937) (1.7645279,-1.2228618) (1.7592919,-1.2283248) (1.7540559,-1.2337825) (1.7488199,-1.2392352) (1.7435839,-1.2446827) (1.7383479,-1.2501252) (1.7331119,-1.2555627) (1.727876,-1.2609952) (1.72264,-1.2664227) (1.717404,-1.2718454) (1.712168,-1.2772632) (1.706932,-1.2826762) (1.701696,-1.2880843) (1.69646,-1.2934877) (1.691224,-1.2988864) (1.6859881,-1.3042804) (1.6807521,-1.3096697) (1.6755161,-1.3150544) (1.6702801,-1.3204345) (1.6650441,-1.3258099) (1.6598081,-1.3311808) (1.6545721,-1.3365472) (1.6493361,-1.3419091) (1.6441002,-1.3472665) (1.6388642,-1.3526194) (1.6336282,-1.3579679) (1.6283922,-1.363312) (1.6231562,-1.3686517) (1.6179202,-1.373987) (1.6126842,-1.3793179) (1.6074482,-1.3846446) (1.6022123,-1.3899669) (1.5969763,-1.3952849) (1.5917403,-1.4005986) (1.5865043,-1.405908) (1.5812683,-1.4112132) (1.5760323,-1.4165142) (1.5707963,-1.421811) (1.5655603,-1.4271035) (1.5603244,-1.4323918) (1.5550884,-1.437676) (1.5498524,-1.442956) (1.5446164,-1.4482318) (1.5393804,-1.4535035) (1.5341444,-1.4587711) (1.5289084,-1.4640345) (1.5236724,-1.4692938) (1.5184364,-1.4745489) (1.5132005,-1.4798) (1.5079645,-1.485047) (1.5027285,-1.4902899) (1.4974925,-1.4955287) (1.4922565,-1.5007634) (1.4870205,-1.505994) (1.4817845,-1.5112206) (1.4765485,-1.5164431) (1.4713126,-1.5216615) (1.4660766,-1.5268759) (1.4608406,-1.5320862) (1.4556046,-1.5372924) (1.4503686,-1.5424946) (1.4451326,-1.5476927) (1.4398966,-1.5528867) (1.4346606,-1.5580767) (1.4294247,-1.5632627) (1.4241887,-1.5684445) (1.4189527,-1.5736223) (1.4137167,-1.5787961) (1.4084807,-1.5839657) (1.4032447,-1.5891313) (1.3980087,-1.5942928) (1.3927727,-1.5994502) (1.3875368,-1.6046035) (1.3823008,-1.6097527) (1.3770648,-1.6148978) (1.3718288,-1.6200388) (1.3665928,-1.6251756) (1.3613568,-1.6303083) (1.3561208,-1.6354369) (1.3508848,-1.6405613) (1.3456489,-1.6456816) (1.3404129,-1.6507977) (1.3351769,-1.6559096) (1.3299409,-1.6610173) (1.3247049,-1.6661208) (1.3194689,-1.6712201) (1.3142329,-1.6763152) (1.3089969,-1.681406) (1.303761,-1.6864925) (1.298525,-1.6915748) (1.293289,-1.6966527) (1.288053,-1.7017264) (1.282817,-1.7067957) (1.277581,-1.7118607) (1.272345,-1.7169213) (1.267109,-1.7219775) (1.261873,-1.7270294) (1.2566371,-1.7320768) (1.2514011,-1.7371198) (1.2461651,-1.7421583) (1.2409291,-1.7471923) (1.2356931,-1.7522218) (1.2304571,-1.7572468) (1.2252211,-1.7622672) (1.2199851,-1.7672831) (1.2147492,-1.7722943) (1.2095132,-1.7773009) (1.2042772,-1.7823029) (1.1990412,-1.7873002) (1.1938052,-1.7922928) (1.1885692,-1.7972806) (1.1833332,-1.8022637) (1.1780972,-1.8072419) (1.1728613,-1.8122154) (1.1676253,-1.817184) (1.1623893,-1.8221477) (1.1571533,-1.8271065) (1.1519173,-1.8320603) (1.1466813,-1.8370091) (1.1414453,-1.8419529) (1.1362093,-1.8468917) (1.1309734,-1.8518254) (1.1257374,-1.8567539) (1.1205014,-1.8616773) (1.1152654,-1.8665954) (1.1100294,-1.8715084) (1.1047934,-1.876416) (1.0995574,-1.8813183) (1.0943214,-1.8862152) (1.0890855,-1.8911068) (1.0838495,-1.8959928) (1.0786135,-1.9008734) (1.0733775,-1.9057484) (1.0681415,-1.9106179) (1.0629055,-1.9154817) (1.0576695,-1.9203398) (1.0524335,-1.9251922) (1.0471976,-1.9300388) (1.0419616,-1.9348796) (1.0367256,-1.9397145) (1.0314896,-1.9445435) (1.0262536,-1.9493664) (1.0210176,-1.9541834) (1.0157816,-1.9589942) (1.0105456,-1.9637989) (1.0053096,-1.9685973) (1.0000737,-1.9733895) (0.99483767,-1.9781754) (0.98960169,-1.9829548) (0.9843657,-1.9877278) (0.97912971,-1.9924943) (0.97389372,-1.9972542) (0.96865773,-2.0020075) (0.96342175,-2.0067541) (0.95818576,-2.0114938) (0.95294977,-2.0162268) (0.94771378,-2.0209528) (0.9424778,-2.0256718) (0.93724181,-2.0303837) (0.93200582,-2.0350885) (0.92676983,-2.0397861) (0.92153385,-2.0444765) (0.91629786,-2.0491594) (0.91106187,-2.0538349) (0.90582588,-2.0585029) (0.90058989,-2.0631632) (0.89535391,-2.0678159) (0.89011792,-2.0724608) (0.88488193,-2.0770978) (0.87964594,-2.0817269) (0.87440996,-2.0863479) (0.86917397,-2.0909608) (0.86393798,-2.0955655) (0.85870199,-2.1001618) (0.853466,-2.1047498) (0.84823002,-2.1093292) (0.84299403,-2.1139) (0.83775804,-2.1184621) (0.83252205,-2.1230153) (0.82728607,-2.1275597) (0.82205008,-2.132095) (0.81681409,-2.1366212) (0.8115781,-2.1411381) (0.80634211,-2.1456457) (0.80110613,-2.1501438) (0.79587014,-2.1546323) (0.79063415,-2.1591111) (0.78539816,-2.16358) (0.78016218,-2.1680391) (0.77492619,-2.172488) (0.7696902,-2.1769267) (0.76445421,-2.1813552) (0.75921822,-2.1857731) (0.75398224,-2.1901805) (0.74874625,-2.1945771) (0.74351026,-2.1989629) (0.73827427,-2.2033376) (0.73303829,-2.2077012) (0.7278023,-2.2120535) (0.72256631,-2.2163944) (0.71733032,-2.2207237) (0.71209433,-2.2250412) (0.70685835,-2.2293469) (0.70162236,-2.2336405) (0.69638637,-2.2379218) (0.69115038,-2.2421908) (0.6859144,-2.2464473) (0.68067841,-2.250691) (0.67544242,-2.2549219) (0.67020643,-2.2591397) (0.66497045,-2.2633443) (0.65973446,-2.2675354) (0.65449847,-2.271713) (0.64926248,-2.2758768) (0.64402649,-2.2800266) (0.63879051,-2.2841622) (0.63355452,-2.2882835) (0.62831853,-2.2923903) (0.62308254,-2.2964823) (0.61784656,-2.3005593) (0.61261057,-2.3046211) (0.60737458,-2.3086676) (0.60213859,-2.3126985) (0.5969026,-2.3167135) (0.59166662,-2.3207126) (0.58643063,-2.3246953) (0.58119464,-2.3286616) (0.57595865,-2.3326111) (0.57072267,-2.3365437) (0.56548668,-2.3404591) (0.56025069,-2.3443571) (0.5550147,-2.3482373) (0.54977871,-2.3520996) (0.54454273,-2.3559437) (0.53930674,-2.3597694) (0.53407075,-2.3635763) (0.52883476,-2.3673643) (0.52359878,-2.371133) (0.51836279,-2.3748821) (0.5131268,-2.3786115) (0.50789081,-2.3823207) (0.50265482,-2.3860096) (0.49741884,-2.3896778) (0.49218285,-2.393325) (0.48694686,-2.396951) (0.48171087,-2.4005554) (0.47647489,-2.4041379) (0.4712389,-2.4076983) (0.46600291,-2.4112361) (0.46076692,-2.4147511) (0.45553093,-2.4182429) (0.45029495,-2.4217113) (0.44505896,-2.4251559) (0.43982297,-2.4285763) (0.43458698,-2.4319722) (0.429351,-2.4353432) (0.42411501,-2.4386891) (0.41887902,-2.4420095) (0.41364303,-2.4453039) (0.40840704,-2.4485721) (0.40317106,-2.4518137) (0.39793507,-2.4550282) (0.39269908,-2.4582154) (0.38746309,-2.4613749) (0.38222711,-2.4645062) (0.37699112,-2.467609) (0.37175513,-2.4706829) (0.36651914,-2.4737275) (0.36128316,-2.4767425) (0.35604717,-2.4797273) (0.35081118,-2.4826816) (0.34557519,-2.4856051) (0.3403392,-2.4884972) (0.33510322,-2.4913576) (0.32986723,-2.4941859) (0.32463124,-2.4969817) (0.31939525,-2.4997445) (0.31415927,-2.5024739) (0.30892328,-2.5051694) (0.30368729,-2.5078308) (0.2984513,-2.5104575) (0.29321531,-2.5130491) (0.28797933,-2.5156051) (0.28274334,-2.5181252) (0.27750735,-2.5206089) (0.27227136,-2.5230558) (0.26703538,-2.5254655) (0.26179939,-2.5278375) (0.2565634,-2.5301713) (0.25132741,-2.5324666) (0.24609142,-2.5347229) (0.24085544,-2.5369398) (0.23561945,-2.5391168) (0.23038346,-2.5412535) (0.22514747,-2.5433496) (0.21991149,-2.5454045) (0.2146755,-2.5474179) (0.20943951,-2.5493893) (0.20420352,-2.5513183) (0.19896753,-2.5532045) (0.19373155,-2.5550475) (0.18849556,-2.5568469) (0.18325957,-2.5586022) (0.17802358,-2.5603132) (0.1727876,-2.5619794) (0.16755161,-2.5636003) (0.16231562,-2.5651758) (0.15707963,-2.5667052) (0.15184364,-2.5681884) (0.14660766,-2.569625) (0.14137167,-2.5710145) (0.13613568,-2.5723567) (0.13089969,-2.5736512) (0.12566371,-2.5748977) (0.12042772,-2.5760959) (0.11519173,-2.5772455) (0.10995574,-2.5783461) (0.10471976,-2.5793976) (0.099483767,-2.5803996) (0.09424778,-2.5813518) (0.089011792,-2.582254) (0.083775804,-2.5831061) (0.078539816,-2.5839076) (0.073303829,-2.5846585) (0.068067841,-2.5853585) (0.062831853,-2.5860074) (0.057595865,-2.5866051) (0.052359878,-2.5871514) (0.04712389,-2.5876461) (0.041887902,-2.5880891) (0.036651914,-2.5884802) (0.031415927,-2.5888195) (0.026179939,-2.5891067) (0.020943951,-2.5893418) (0.015707963,-2.5895247) (0.010471976,-2.5896554) (0.0052359878,-2.5897339) (-4.4408921e-16,-2.58976) (-0.0052359878,-2.5897339) (-0.010471976,-2.5896554) (-0.015707963,-2.5895247) (-0.020943951,-2.5893418) (-0.026179939,-2.5891067) (-0.031415927,-2.5888195) (-0.036651914,-2.5884802) (-0.041887902,-2.5880891) (-0.04712389,-2.5876461) (-0.052359878,-2.5871514) (-0.057595865,-2.5866051) (-0.062831853,-2.5860074) (-0.068067841,-2.5853585) (-0.073303829,-2.5846585) (-0.078539816,-2.5839076) (-0.083775804,-2.5831061) (-0.089011792,-2.582254) (-0.09424778,-2.5813518) (-0.099483767,-2.5803996) (-0.10471976,-2.5793976) (-0.10995574,-2.5783461) (-0.11519173,-2.5772455) (-0.12042772,-2.5760959) (-0.12566371,-2.5748977) (-0.13089969,-2.5736512) (-0.13613568,-2.5723567) (-0.14137167,-2.5710145) (-0.14660766,-2.569625) (-0.15184364,-2.5681884) (-0.15707963,-2.5667052) (-0.16231562,-2.5651758) (-0.16755161,-2.5636003) (-0.1727876,-2.5619794) (-0.17802358,-2.5603132) (-0.18325957,-2.5586022) (-0.18849556,-2.5568469) (-0.19373155,-2.5550475) (-0.19896753,-2.5532045) (-0.20420352,-2.5513183) (-0.20943951,-2.5493893) (-0.2146755,-2.5474179) (-0.21991149,-2.5454045) (-0.22514747,-2.5433496) (-0.23038346,-2.5412535) (-0.23561945,-2.5391168) (-0.24085544,-2.5369398) (-0.24609142,-2.5347229) (-0.25132741,-2.5324666) (-0.2565634,-2.5301713) (-0.26179939,-2.5278375) (-0.26703538,-2.5254655) (-0.27227136,-2.5230558) (-0.27750735,-2.5206089) (-0.28274334,-2.5181252) (-0.28797933,-2.5156051) (-0.29321531,-2.5130491) (-0.2984513,-2.5104575) (-0.30368729,-2.5078308) (-0.30892328,-2.5051694) (-0.31415927,-2.5024739) (-0.31939525,-2.4997445) (-0.32463124,-2.4969817) (-0.32986723,-2.4941859) (-0.33510322,-2.4913576) (-0.3403392,-2.4884972) (-0.34557519,-2.4856051) (-0.35081118,-2.4826816) (-0.35604717,-2.4797273) (-0.36128316,-2.4767425) (-0.36651914,-2.4737275) (-0.37175513,-2.4706829) (-0.37699112,-2.467609) (-0.38222711,-2.4645062) (-0.38746309,-2.4613749) (-0.39269908,-2.4582154) (-0.39793507,-2.4550282) (-0.40317106,-2.4518137) (-0.40840704,-2.4485721) (-0.41364303,-2.4453039) (-0.41887902,-2.4420095) (-0.42411501,-2.4386891) (-0.429351,-2.4353432) (-0.43458698,-2.4319722) (-0.43982297,-2.4285763) (-0.44505896,-2.4251559) (-0.45029495,-2.4217113) (-0.45553093,-2.4182429) (-0.46076692,-2.4147511) (-0.46600291,-2.4112361) (-0.4712389,-2.4076983) (-0.47647489,-2.4041379) (-0.48171087,-2.4005554) (-0.48694686,-2.396951) (-0.49218285,-2.393325) (-0.49741884,-2.3896778) (-0.50265482,-2.3860096) (-0.50789081,-2.3823207) (-0.5131268,-2.3786115) (-0.51836279,-2.3748821) (-0.52359878,-2.371133) (-0.52883476,-2.3673643) (-0.53407075,-2.3635763) (-0.53930674,-2.3597694) (-0.54454273,-2.3559437) (-0.54977871,-2.3520996) (-0.5550147,-2.3482373) (-0.56025069,-2.3443571) (-0.56548668,-2.3404591) (-0.57072267,-2.3365437) (-0.57595865,-2.3326111) (-0.58119464,-2.3286616) (-0.58643063,-2.3246953) (-0.59166662,-2.3207126) (-0.5969026,-2.3167135) (-0.60213859,-2.3126985) (-0.60737458,-2.3086676) (-0.61261057,-2.3046211) (-0.61784656,-2.3005593) (-0.62308254,-2.2964823) (-0.62831853,-2.2923903) (-0.63355452,-2.2882835) (-0.63879051,-2.2841622) (-0.64402649,-2.2800266) (-0.64926248,-2.2758768) (-0.65449847,-2.271713) (-0.65973446,-2.2675354) (-0.66497045,-2.2633443) (-0.67020643,-2.2591397) (-0.67544242,-2.2549219) (-0.68067841,-2.250691) (-0.6859144,-2.2464473) (-0.69115038,-2.2421908) (-0.69638637,-2.2379218) (-0.70162236,-2.2336405) (-0.70685835,-2.2293469) (-0.71209433,-2.2250412) (-0.71733032,-2.2207237) (-0.72256631,-2.2163944) (-0.7278023,-2.2120535) (-0.73303829,-2.2077012) (-0.73827427,-2.2033376) (-0.74351026,-2.1989629) (-0.74874625,-2.1945771) (-0.75398224,-2.1901805) (-0.75921822,-2.1857731) (-0.76445421,-2.1813552) (-0.7696902,-2.1769267) (-0.77492619,-2.172488) (-0.78016218,-2.1680391) (-0.78539816,-2.16358) (-0.79063415,-2.1591111) (-0.79587014,-2.1546323) (-0.80110613,-2.1501438) (-0.80634211,-2.1456457) (-0.8115781,-2.1411381) (-0.81681409,-2.1366212) (-0.82205008,-2.132095) (-0.82728607,-2.1275597) (-0.83252205,-2.1230153) (-0.83775804,-2.1184621) (-0.84299403,-2.1139) (-0.84823002,-2.1093292) (-0.853466,-2.1047498) (-0.85870199,-2.1001618) (-0.86393798,-2.0955655) (-0.86917397,-2.0909608) (-0.87440996,-2.0863479) (-0.87964594,-2.0817269) (-0.88488193,-2.0770978) (-0.89011792,-2.0724608) (-0.89535391,-2.0678159) (-0.90058989,-2.0631632) (-0.90582588,-2.0585029) (-0.91106187,-2.0538349) (-0.91629786,-2.0491594) (-0.92153385,-2.0444765) (-0.92676983,-2.0397861) (-0.93200582,-2.0350885) (-0.93724181,-2.0303837) (-0.9424778,-2.0256718) (-0.94771378,-2.0209528) (-0.95294977,-2.0162268) (-0.95818576,-2.0114938) (-0.96342175,-2.0067541) (-0.96865773,-2.0020075) (-0.97389372,-1.9972542) (-0.97912971,-1.9924943) (-0.9843657,-1.9877278) (-0.98960169,-1.9829548) (-0.99483767,-1.9781754) (-1.0000737,-1.9733895) (-1.0053096,-1.9685973) (-1.0105456,-1.9637989) (-1.0157816,-1.9589942) (-1.0210176,-1.9541834) (-1.0262536,-1.9493664) (-1.0314896,-1.9445435) (-1.0367256,-1.9397145) (-1.0419616,-1.9348796) (-1.0471976,-1.9300388) (-1.0524335,-1.9251922) (-1.0576695,-1.9203398) (-1.0629055,-1.9154817) (-1.0681415,-1.9106179) (-1.0733775,-1.9057484) (-1.0786135,-1.9008734) (-1.0838495,-1.8959928) (-1.0890855,-1.8911068) (-1.0943214,-1.8862152) (-1.0995574,-1.8813183) (-1.1047934,-1.876416) (-1.1100294,-1.8715084) (-1.1152654,-1.8665954) (-1.1205014,-1.8616773) (-1.1257374,-1.8567539) (-1.1309734,-1.8518254) (-1.1362093,-1.8468917) (-1.1414453,-1.8419529) (-1.1466813,-1.8370091) (-1.1519173,-1.8320603) (-1.1571533,-1.8271065) (-1.1623893,-1.8221477) (-1.1676253,-1.817184) (-1.1728613,-1.8122154) (-1.1780972,-1.8072419) (-1.1833332,-1.8022637) (-1.1885692,-1.7972806) (-1.1938052,-1.7922928) (-1.1990412,-1.7873002) (-1.2042772,-1.7823029) (-1.2095132,-1.7773009) (-1.2147492,-1.7722943) (-1.2199851,-1.7672831) (-1.2252211,-1.7622672) (-1.2304571,-1.7572468) (-1.2356931,-1.7522218) (-1.2409291,-1.7471923) (-1.2461651,-1.7421583) (-1.2514011,-1.7371198) (-1.2566371,-1.7320768) (-1.261873,-1.7270294) (-1.267109,-1.7219775) (-1.272345,-1.7169213) (-1.277581,-1.7118607) (-1.282817,-1.7067957) (-1.288053,-1.7017264) (-1.293289,-1.6966527) (-1.298525,-1.6915748) (-1.303761,-1.6864925) (-1.3089969,-1.681406) (-1.3142329,-1.6763152) (-1.3194689,-1.6712201) (-1.3247049,-1.6661208) (-1.3299409,-1.6610173) (-1.3351769,-1.6559096) (-1.3404129,-1.6507977) (-1.3456489,-1.6456816) (-1.3508848,-1.6405613) (-1.3561208,-1.6354369) (-1.3613568,-1.6303083) (-1.3665928,-1.6251756) (-1.3718288,-1.6200388) (-1.3770648,-1.6148978) (-1.3823008,-1.6097527) (-1.3875368,-1.6046035) (-1.3927727,-1.5994502) (-1.3980087,-1.5942928) (-1.4032447,-1.5891313) (-1.4084807,-1.5839657) (-1.4137167,-1.5787961) (-1.4189527,-1.5736223) (-1.4241887,-1.5684445) (-1.4294247,-1.5632627) (-1.4346606,-1.5580767) (-1.4398966,-1.5528867) (-1.4451326,-1.5476927) (-1.4503686,-1.5424946) (-1.4556046,-1.5372924) (-1.4608406,-1.5320862) (-1.4660766,-1.5268759) (-1.4713126,-1.5216615) (-1.4765485,-1.5164431) (-1.4817845,-1.5112206) (-1.4870205,-1.505994) (-1.4922565,-1.5007634) (-1.4974925,-1.4955287) (-1.5027285,-1.4902899) (-1.5079645,-1.485047) (-1.5132005,-1.4798) (-1.5184364,-1.4745489) (-1.5236724,-1.4692938) (-1.5289084,-1.4640345) (-1.5341444,-1.4587711) (-1.5393804,-1.4535035) (-1.5446164,-1.4482318) (-1.5498524,-1.442956) (-1.5550884,-1.437676) (-1.5603244,-1.4323918) (-1.5655603,-1.4271035) (-1.5707963,-1.421811) (-1.5760323,-1.4165142) (-1.5812683,-1.4112132) (-1.5865043,-1.405908) (-1.5917403,-1.4005986) (-1.5969763,-1.3952849) (-1.6022123,-1.3899669) (-1.6074482,-1.3846446) (-1.6126842,-1.3793179) (-1.6179202,-1.373987) (-1.6231562,-1.3686517) (-1.6283922,-1.363312) (-1.6336282,-1.3579679) (-1.6388642,-1.3526194) (-1.6441002,-1.3472665) (-1.6493361,-1.3419091) (-1.6545721,-1.3365472) (-1.6598081,-1.3311808) (-1.6650441,-1.3258099) (-1.6702801,-1.3204345) (-1.6755161,-1.3150544) (-1.6807521,-1.3096697) (-1.6859881,-1.3042804) (-1.691224,-1.2988864) (-1.69646,-1.2934877) (-1.701696,-1.2880843) (-1.706932,-1.2826762) (-1.712168,-1.2772632) (-1.717404,-1.2718454) (-1.72264,-1.2664227) (-1.727876,-1.2609952) (-1.7331119,-1.2555627) (-1.7383479,-1.2501252) (-1.7435839,-1.2446827) (-1.7488199,-1.2392352) (-1.7540559,-1.2337825) (-1.7592919,-1.2283248) (-1.7645279,-1.2228618) (-1.7697639,-1.2173937) (-1.7749998,-1.2119203) (-1.7802358,-1.2064415) (-1.7854718,-1.2009575) (-1.7907078,-1.195468) (-1.7959438,-1.189973) (-1.8011798,-1.1844725) (-1.8064158,-1.1789665) (-1.8116518,-1.1734549) (-1.8168878,-1.1679376) (-1.8221237,-1.1624145) (-1.8273597,-1.1568857) (-1.8325957,-1.1513511) (-1.8378317,-1.1458105) (-1.8430677,-1.140264) (-1.8483037,-1.1347114) (-1.8535397,-1.1291527) (-1.8587757,-1.1235879) (-1.8640116,-1.1180169) (-1.8692476,-1.1124395) (-1.8744836,-1.1068558) (-1.8797196,-1.1012656) (-1.8849556,-1.0956688) (-1.8901916,-1.0900655) (-1.8954276,-1.0844555) (-1.9006636,-1.0788387) (-1.9058995,-1.0732151) (-1.9111355,-1.0675845) (-1.9163715,-1.0619469) (-1.9216075,-1.0563022) (-1.9268435,-1.0506503) (-1.9320795,-1.0449911) (-1.9373155,-1.0393244) (-1.9425515,-1.0336503) (-1.9477874,-1.0279685) (-1.9530234,-1.022279) (-1.9582594,-1.0165818) (-1.9634954,-1.0108765) (-1.9687314,-1.0051632) (-1.9739674,-0.9994418) (-1.9792034,-0.99371206) (-1.9844394,-0.98797392) (-1.9896753,-0.98222723) (-1.9949113,-0.97647187) (-2.0001473,-0.97070771) (-2.0053833,-0.96493461) (-2.0106193,-0.95915242) (-2.0158553,-0.95336101) (-2.0210913,-0.94756022) (-2.0263273,-0.94174989) (-2.0315632,-0.93592987) (-2.0367992,-0.9301) (-2.0420352,-0.92426012) (-2.0472712,-0.91841003) (-2.0525072,-0.91254958) (-2.0577432,-0.90667858) (-2.0629792,-0.90079685) (-2.0682152,-0.89490418) (-2.0734512,-0.88900038) (-2.0786871,-0.88308525) (-2.0839231,-0.87715857) (-2.0891591,-0.87122013) (-2.0943951,-0.86526971) (-2.0996311,-0.85930708) (-2.1048671,-0.85333199) (-2.1101031,-0.8473442) (-2.1153391,-0.84134346) (-2.120575,-0.83532951) (-2.125811,-0.82930208) (-2.131047,-0.82326089) (-2.136283,-0.81720566) (-2.141519,-0.81113609) (-2.146755,-0.80505186) (-2.151991,-0.79895267) (-2.157227,-0.79283819) (-2.1624629,-0.78670808) (-2.1676989,-0.78056199) (-2.1729349,-0.77439955) (-2.1781709,-0.76822039) (-2.1834069,-0.76202412) (-2.1886429,-0.75581034) (-2.1938789,-0.74957863) (-2.1991149,-0.74332856) (-2.2043508,-0.73705967) (-2.2095868,-0.7307715) (-2.2148228,-0.72446357) (-2.2200588,-0.71813537) (-2.2252948,-0.71178638) (-2.2305308,-0.70541605) (-2.2357668,-0.69902383) (-2.2410028,-0.69260911) (-2.2462387,-0.68617129) (-2.2514747,-0.67970973) (-2.2567107,-0.67322376) (-2.2619467,-0.66671268) (-2.2671827,-0.66017578) (-2.2724187,-0.65361229) (-2.2776547,-0.64702142) (-2.2828907,-0.64040234) (-2.2881266,-0.63375418) (-2.2933626,-0.62707605) (-2.2985986,-0.62036699) (-2.3038346,-0.61362599) (-2.3090706,-0.60685202) (-2.3143066,-0.60004397) (-2.3195426,-0.59320069) (-2.3247786,-0.58632097) (-2.3300146,-0.57940351) (-2.3352505,-0.57244698) (-2.3404865,-0.56544994) (-2.3457225,-0.55841089) (-2.3509585,-0.55132824) (-2.3561945,-0.5442003) (-2.3614305,-0.53702529) (-2.3666665,-0.52980132) (-2.3719025,-0.52252636) (-2.3771384,-0.51519828) (-2.3823744,-0.5078148) (-2.3876104,-0.50037348) (-2.3928464,-0.49287173) (-2.3980824,-0.48530674) (-2.4033184,-0.47767555) (-2.4085544,-0.46997493) (-2.4137904,-0.46220145) (-2.4190263,-0.45435138) (-2.4242623,-0.44642069) (-2.4294983,-0.43840503) (-2.4347343,-0.43029968) (-2.4399703,-0.42209948) (-2.4452063,-0.41379881) (-2.4504423,-0.40539153) (-2.4556783,-0.39687087) (-2.4609142,-0.38822939) (-2.4661502,-0.37945885) (-2.4713862,-0.37055012) (-2.4766222,-0.36149299) (-2.4818582,-0.35227605) (-2.4870942,-0.34288642) (-2.4923302,-0.33330953) (-2.4975662,-0.32352878) (-2.5028021,-0.31352511) (-2.5080381,-0.30327648) (-2.5132741,-0.29275719) (-2.5185101,-0.28193697) (-2.5237461,-0.27077978) (-2.5289821,-0.25924214) (-2.5342181,-0.24727081) (-2.5394541,-0.23479951) (-2.54469,-0.22174392) (-2.549926,-0.20799408) (-2.555162,-0.193402) (-2.560398,-0.17776042) (-2.565634,-0.16076333) (-2.57087,-0.1419246) (-2.576106,-0.12038277) (-2.581342,-0.094303316) (-2.586578,-0.057843896)};
\addlegendentry{$p=0.15$}
\addplot[thick,color=blue] coordinates {(-2.2619467,0.060750487) (-2.2567107,0.10856931) (-2.2514747,0.14117281) (-2.2462387,0.16768416) (-2.2410028,0.19066173) (-2.2357668,0.21126161) (-2.2305308,0.23012311) (-2.2252948,0.24764367) (-2.2200588,0.26409022) (-2.2148228,0.27965222) (-2.2095868,0.29446986) (-2.2043508,0.30865028) (-2.1991149,0.32227754) (-2.1938789,0.33541898) (-2.1886429,0.34812956) (-2.1834069,0.36045478) (-2.1781709,0.37243284) (-2.1729349,0.38409615) (-2.1676989,0.39547249) (-2.1624629,0.40658588) (-2.157227,0.41745726) (-2.151991,0.42810501) (-2.146755,0.43854534) (-2.141519,0.44879268) (-2.136283,0.4588599) (-2.131047,0.46875853) (-2.125811,0.47849898) (-2.120575,0.48809067) (-2.1153391,0.49754213) (-2.1101031,0.50686115) (-2.1048671,0.51605486) (-2.0996311,0.52512977) (-2.0943951,0.53409189) (-2.0891591,0.54294674) (-2.0839231,0.55169943) (-2.0786871,0.5603547) (-2.0734512,0.56891693) (-2.0682152,0.57739021) (-2.0629792,0.58577835) (-2.0577432,0.59408489) (-2.0525072,0.60231316) (-2.0472712,0.61046626) (-2.0420352,0.61854713) (-2.0367992,0.62655849) (-2.0315632,0.63450293) (-2.0263273,0.64238288) (-2.0210913,0.65020063) (-2.0158553,0.65795833) (-2.0106193,0.66565805) (-2.0053833,0.67330171) (-2.0001473,0.68089115) (-1.9949113,0.68842811) (-1.9896753,0.69591425) (-1.9844394,0.70335112) (-1.9792034,0.71074024) (-1.9739674,0.71808301) (-1.9687314,0.72538079) (-1.9634954,0.73263488) (-1.9582594,0.7398465) (-1.9530234,0.74701684) (-1.9477874,0.75414701) (-1.9425515,0.76123811) (-1.9373155,0.76829114) (-1.9320795,0.77530711) (-1.9268435,0.78228695) (-1.9216075,0.78923157) (-1.9163715,0.79614185) (-1.9111355,0.8030186) (-1.9058995,0.80986264) (-1.9006636,0.81667473) (-1.8954276,0.82345562) (-1.8901916,0.83020601) (-1.8849556,0.83692659) (-1.8797196,0.84361802) (-1.8744836,0.85028093) (-1.8692476,0.85691594) (-1.8640116,0.86352363) (-1.8587757,0.87010458) (-1.8535397,0.87665932) (-1.8483037,0.8831884) (-1.8430677,0.88969232) (-1.8378317,0.89617157) (-1.8325957,0.90262664) (-1.8273597,0.90905797) (-1.8221237,0.91546603) (-1.8168878,0.92185123) (-1.8116518,0.928214) (-1.8064158,0.93455473) (-1.8011798,0.94087382) (-1.7959438,0.94717165) (-1.7907078,0.95344858) (-1.7854718,0.95970496) (-1.7802358,0.96594114) (-1.7749998,0.97215746) (-1.7697639,0.97835422) (-1.7645279,0.98453175) (-1.7592919,0.99069035) (-1.7540559,0.99683031) (-1.7488199,1.0029519) (-1.7435839,1.0090554) (-1.7383479,1.0151412) (-1.7331119,1.0212093) (-1.727876,1.0272602) (-1.72264,1.033294) (-1.717404,1.039311) (-1.712168,1.0453115) (-1.706932,1.0512956) (-1.701696,1.0572635) (-1.69646,1.0632155) (-1.691224,1.0691518) (-1.6859881,1.0750726) (-1.6807521,1.080978) (-1.6755161,1.0868683) (-1.6702801,1.0927437) (-1.6650441,1.0986043) (-1.6598081,1.1044502) (-1.6545721,1.1102818) (-1.6493361,1.1160991) (-1.6441002,1.1219022) (-1.6388642,1.1276914) (-1.6336282,1.1334669) (-1.6283922,1.1392286) (-1.6231562,1.1449768) (-1.6179202,1.1507117) (-1.6126842,1.1564333) (-1.6074482,1.1621418) (-1.6022123,1.1678374) (-1.5969763,1.17352) (-1.5917403,1.1791899) (-1.5865043,1.1848472) (-1.5812683,1.190492) (-1.5760323,1.1961244) (-1.5707963,1.2017445) (-1.5655603,1.2073524) (-1.5603244,1.2129482) (-1.5550884,1.218532) (-1.5498524,1.2241039) (-1.5446164,1.2296641) (-1.5393804,1.2352125) (-1.5341444,1.2407493) (-1.5289084,1.2462745) (-1.5236724,1.2517883) (-1.5184364,1.2572907) (-1.5132005,1.2627818) (-1.5079645,1.2682616) (-1.5027285,1.2737303) (-1.4974925,1.279188) (-1.4922565,1.2846346) (-1.4870205,1.2900702) (-1.4817845,1.295495) (-1.4765485,1.3009089) (-1.4713126,1.306312) (-1.4660766,1.3117045) (-1.4608406,1.3170863) (-1.4556046,1.3224575) (-1.4503686,1.3278181) (-1.4451326,1.3331682) (-1.4398966,1.3385079) (-1.4346606,1.3438372) (-1.4294247,1.3491561) (-1.4241887,1.3544647) (-1.4189527,1.359763) (-1.4137167,1.365051) (-1.4084807,1.3703289) (-1.4032447,1.3755966) (-1.3980087,1.3808541) (-1.3927727,1.3861015) (-1.3875368,1.3913389) (-1.3823008,1.3965662) (-1.3770648,1.4017835) (-1.3718288,1.4069908) (-1.3665928,1.4121882) (-1.3613568,1.4173756) (-1.3561208,1.4225531) (-1.3508848,1.4277207) (-1.3456489,1.4328784) (-1.3404129,1.4380262) (-1.3351769,1.4431642) (-1.3299409,1.4482923) (-1.3247049,1.4534107) (-1.3194689,1.4585192) (-1.3142329,1.4636179) (-1.3089969,1.4687069) (-1.303761,1.473786) (-1.298525,1.4788554) (-1.293289,1.483915) (-1.288053,1.4889649) (-1.282817,1.494005) (-1.277581,1.4990353) (-1.272345,1.5040558) (-1.267109,1.5090667) (-1.261873,1.5140677) (-1.2566371,1.519059) (-1.2514011,1.5240405) (-1.2461651,1.5290122) (-1.2409291,1.5339742) (-1.2356931,1.5389263) (-1.2304571,1.5438687) (-1.2252211,1.5488012) (-1.2199851,1.5537239) (-1.2147492,1.5586368) (-1.2095132,1.5635398) (-1.2042772,1.568433) (-1.1990412,1.5733163) (-1.1938052,1.5781897) (-1.1885692,1.5830531) (-1.1833332,1.5879067) (-1.1780972,1.5927503) (-1.1728613,1.5975838) (-1.1676253,1.6024074) (-1.1623893,1.607221) (-1.1571533,1.6120245) (-1.1519173,1.6168179) (-1.1466813,1.6216013) (-1.1414453,1.6263745) (-1.1362093,1.6311375) (-1.1309734,1.6358903) (-1.1257374,1.6406329) (-1.1205014,1.6453652) (-1.1152654,1.6500873) (-1.1100294,1.654799) (-1.1047934,1.6595003) (-1.0995574,1.6641913) (-1.0943214,1.6688718) (-1.0890855,1.6735418) (-1.0838495,1.6782013) (-1.0786135,1.6828502) (-1.0733775,1.6874885) (-1.0681415,1.6921161) (-1.0629055,1.6967331) (-1.0576695,1.7013393) (-1.0524335,1.7059347) (-1.0471976,1.7105192) (-1.0419616,1.7150929) (-1.0367256,1.7196556) (-1.0314896,1.7242074) (-1.0262536,1.728748) (-1.0210176,1.7332776) (-1.0157816,1.737796) (-1.0105456,1.7423032) (-1.0053096,1.7467991) (-1.0000737,1.7512837) (-0.99483767,1.7557569) (-0.98960169,1.7602186) (-0.9843657,1.7646688) (-0.97912971,1.7691075) (-0.97389372,1.7735344) (-0.96865773,1.7779497) (-0.96342175,1.7823532) (-0.95818576,1.7867448) (-0.95294977,1.7911245) (-0.94771378,1.7954923) (-0.9424778,1.7998479) (-0.93724181,1.8041914) (-0.93200582,1.8085228) (-0.92676983,1.8128418) (-0.92153385,1.8171485) (-0.91629786,1.8214427) (-0.91106187,1.8257244) (-0.90582588,1.8299935) (-0.90058989,1.8342499) (-0.89535391,1.8384936) (-0.89011792,1.8427244) (-0.88488193,1.8469422) (-0.87964594,1.8511471) (-0.87440996,1.8553388) (-0.86917397,1.8595174) (-0.86393798,1.8636826) (-0.85870199,1.8678345) (-0.853466,1.8719729) (-0.84823002,1.8760977) (-0.84299403,1.8802088) (-0.83775804,1.8843062) (-0.83252205,1.8883898) (-0.82728607,1.8924593) (-0.82205008,1.8965148) (-0.81681409,1.9005562) (-0.8115781,1.9045833) (-0.80634211,1.908596) (-0.80110613,1.9125942) (-0.79587014,1.9165779) (-0.79063415,1.9205468) (-0.78539816,1.924501) (-0.78016218,1.9284402) (-0.77492619,1.9323645) (-0.7696902,1.9362735) (-0.76445421,1.9401674) (-0.75921822,1.9440458) (-0.75398224,1.9479088) (-0.74874625,1.9517561) (-0.74351026,1.9555878) (-0.73827427,1.9594035) (-0.73303829,1.9632033) (-0.7278023,1.966987) (-0.72256631,1.9707545) (-0.71733032,1.9745056) (-0.71209433,1.9782402) (-0.70685835,1.9819583) (-0.70162236,1.9856596) (-0.69638637,1.989344) (-0.69115038,1.9930114) (-0.6859144,1.9966616) (-0.68067841,2.0002946) (-0.67544242,2.0039102) (-0.67020643,2.0075082) (-0.66497045,2.0110885) (-0.65973446,2.0146509) (-0.65449847,2.0181954) (-0.64926248,2.0217218) (-0.64402649,2.0252298) (-0.63879051,2.0287195) (-0.63355452,2.0321906) (-0.62831853,2.035643) (-0.62308254,2.0390765) (-0.61784656,2.042491) (-0.61261057,2.0458863) (-0.60737458,2.0492623) (-0.60213859,2.0526188) (-0.5969026,2.0559557) (-0.59166662,2.0592727) (-0.58643063,2.0625698) (-0.58119464,2.0658468) (-0.57595865,2.0691036) (-0.57072267,2.0723398) (-0.56548668,2.0755555) (-0.56025069,2.0787504) (-0.5550147,2.0819244) (-0.54977871,2.0850772) (-0.54454273,2.0882088) (-0.53930674,2.091319) (-0.53407075,2.0944075) (-0.52883476,2.0974743) (-0.52359878,2.1005191) (-0.51836279,2.1035418) (-0.5131268,2.1065422) (-0.50789081,2.1095201) (-0.50265482,2.1124754) (-0.49741884,2.1154078) (-0.49218285,2.1183173) (-0.48694686,2.1212036) (-0.48171087,2.1240665) (-0.47647489,2.1269059) (-0.4712389,2.1297215) (-0.46600291,2.1325133) (-0.46076692,2.135281) (-0.45553093,2.1380244) (-0.45029495,2.1407434) (-0.44505896,2.1434378) (-0.43982297,2.1461074) (-0.43458698,2.1487519) (-0.429351,2.1513714) (-0.42411501,2.1539654) (-0.41887902,2.1565339) (-0.41364303,2.1590767) (-0.40840704,2.1615936) (-0.40317106,2.1640844) (-0.39793507,2.166549) (-0.39269908,2.168987) (-0.38746309,2.1713984) (-0.38222711,2.173783) (-0.37699112,2.1761406) (-0.37175513,2.178471) (-0.36651914,2.180774) (-0.36128316,2.1830494) (-0.35604717,2.1852971) (-0.35081118,2.1875169) (-0.34557519,2.1897086) (-0.3403392,2.1918719) (-0.33510322,2.1940068) (-0.32986723,2.196113) (-0.32463124,2.1981905) (-0.31939525,2.2002389) (-0.31415927,2.2022581) (-0.30892328,2.2042479) (-0.30368729,2.2062082) (-0.2984513,2.2081388) (-0.29321531,2.2100395) (-0.28797933,2.2119102) (-0.28274334,2.2137506) (-0.27750735,2.2155606) (-0.27227136,2.21734) (-0.26703538,2.2190888) (-0.26179939,2.2208066) (-0.2565634,2.2224933) (-0.25132741,2.2241489) (-0.24609142,2.225773) (-0.24085544,2.2273656) (-0.23561945,2.2289265) (-0.23038346,2.2304556) (-0.22514747,2.2319526) (-0.21991149,2.2334176) (-0.2146755,2.2348502) (-0.20943951,2.2362503) (-0.20420352,2.2376179) (-0.19896753,2.2389528) (-0.19373155,2.2402548) (-0.18849556,2.2415237) (-0.18325957,2.2427596) (-0.17802358,2.2439622) (-0.1727876,2.2451314) (-0.16755161,2.2462671) (-0.16231562,2.2473691) (-0.15707963,2.2484374) (-0.15184364,2.2494719) (-0.14660766,2.2504723) (-0.14137167,2.2514387) (-0.13613568,2.2523709) (-0.13089969,2.2532688) (-0.12566371,2.2541322) (-0.12042772,2.2549612) (-0.11519173,2.2557556) (-0.10995574,2.2565154) (-0.10471976,2.2572403) (-0.099483767,2.2579305) (-0.09424778,2.2585857) (-0.089011792,2.2592059) (-0.083775804,2.2597911) (-0.078539816,2.2603412) (-0.073303829,2.260856) (-0.068067841,2.2613356) (-0.062831853,2.2617799) (-0.057595865,2.2621889) (-0.052359878,2.2625624) (-0.04712389,2.2629005) (-0.041887902,2.2632031) (-0.036651914,2.2634702) (-0.031415927,2.2637018) (-0.026179939,2.2638977) (-0.020943951,2.2640581) (-0.015707963,2.2641828) (-0.010471976,2.264272) (-0.0052359878,2.2643254) (-4.4408921e-16,2.2643433) (0.0052359878,2.2643254) (0.010471976,2.264272) (0.015707963,2.2641828) (0.020943951,2.2640581) (0.026179939,2.2638977) (0.031415927,2.2637018) (0.036651914,2.2634702) (0.041887902,2.2632031) (0.04712389,2.2629005) (0.052359878,2.2625624) (0.057595865,2.2621889) (0.062831853,2.2617799) (0.068067841,2.2613356) (0.073303829,2.260856) (0.078539816,2.2603412) (0.083775804,2.2597911) (0.089011792,2.2592059) (0.09424778,2.2585857) (0.099483767,2.2579305) (0.10471976,2.2572403) (0.10995574,2.2565154) (0.11519173,2.2557556) (0.12042772,2.2549612) (0.12566371,2.2541322) (0.13089969,2.2532688) (0.13613568,2.2523709) (0.14137167,2.2514387) (0.14660766,2.2504723) (0.15184364,2.2494719) (0.15707963,2.2484374) (0.16231562,2.2473691) (0.16755161,2.2462671) (0.1727876,2.2451314) (0.17802358,2.2439622) (0.18325957,2.2427596) (0.18849556,2.2415237) (0.19373155,2.2402548) (0.19896753,2.2389528) (0.20420352,2.2376179) (0.20943951,2.2362503) (0.2146755,2.2348502) (0.21991149,2.2334176) (0.22514747,2.2319526) (0.23038346,2.2304556) (0.23561945,2.2289265) (0.24085544,2.2273656) (0.24609142,2.225773) (0.25132741,2.2241489) (0.2565634,2.2224933) (0.26179939,2.2208066) (0.26703538,2.2190888) (0.27227136,2.21734) (0.27750735,2.2155606) (0.28274334,2.2137506) (0.28797933,2.2119102) (0.29321531,2.2100395) (0.2984513,2.2081388) (0.30368729,2.2062082) (0.30892328,2.2042479) (0.31415927,2.2022581) (0.31939525,2.2002389) (0.32463124,2.1981905) (0.32986723,2.196113) (0.33510322,2.1940068) (0.3403392,2.1918719) (0.34557519,2.1897086) (0.35081118,2.1875169) (0.35604717,2.1852971) (0.36128316,2.1830494) (0.36651914,2.180774) (0.37175513,2.178471) (0.37699112,2.1761406) (0.38222711,2.173783) (0.38746309,2.1713984) (0.39269908,2.168987) (0.39793507,2.166549) (0.40317106,2.1640844) (0.40840704,2.1615936) (0.41364303,2.1590767) (0.41887902,2.1565339) (0.42411501,2.1539654) (0.429351,2.1513714) (0.43458698,2.1487519) (0.43982297,2.1461074) (0.44505896,2.1434378) (0.45029495,2.1407434) (0.45553093,2.1380244) (0.46076692,2.135281) (0.46600291,2.1325133) (0.4712389,2.1297215) (0.47647489,2.1269059) (0.48171087,2.1240665) (0.48694686,2.1212036) (0.49218285,2.1183173) (0.49741884,2.1154078) (0.50265482,2.1124754) (0.50789081,2.1095201) (0.5131268,2.1065422) (0.51836279,2.1035418) (0.52359878,2.1005191) (0.52883476,2.0974743) (0.53407075,2.0944075) (0.53930674,2.091319) (0.54454273,2.0882088) (0.54977871,2.0850772) (0.5550147,2.0819244) (0.56025069,2.0787504) (0.56548668,2.0755555) (0.57072267,2.0723398) (0.57595865,2.0691036) (0.58119464,2.0658468) (0.58643063,2.0625698) (0.59166662,2.0592727) (0.5969026,2.0559557) (0.60213859,2.0526188) (0.60737458,2.0492623) (0.61261057,2.0458863) (0.61784656,2.042491) (0.62308254,2.0390765) (0.62831853,2.035643) (0.63355452,2.0321906) (0.63879051,2.0287195) (0.64402649,2.0252298) (0.64926248,2.0217218) (0.65449847,2.0181954) (0.65973446,2.0146509) (0.66497045,2.0110885) (0.67020643,2.0075082) (0.67544242,2.0039102) (0.68067841,2.0002946) (0.6859144,1.9966616) (0.69115038,1.9930114) (0.69638637,1.989344) (0.70162236,1.9856596) (0.70685835,1.9819583) (0.71209433,1.9782402) (0.71733032,1.9745056) (0.72256631,1.9707545) (0.7278023,1.966987) (0.73303829,1.9632033) (0.73827427,1.9594035) (0.74351026,1.9555878) (0.74874625,1.9517561) (0.75398224,1.9479088) (0.75921822,1.9440458) (0.76445421,1.9401674) (0.7696902,1.9362735) (0.77492619,1.9323645) (0.78016218,1.9284402) (0.78539816,1.924501) (0.79063415,1.9205468) (0.79587014,1.9165779) (0.80110613,1.9125942) (0.80634211,1.908596) (0.8115781,1.9045833) (0.81681409,1.9005562) (0.82205008,1.8965148) (0.82728607,1.8924593) (0.83252205,1.8883898) (0.83775804,1.8843062) (0.84299403,1.8802088) (0.84823002,1.8760977) (0.853466,1.8719729) (0.85870199,1.8678345) (0.86393798,1.8636826) (0.86917397,1.8595174) (0.87440996,1.8553388) (0.87964594,1.8511471) (0.88488193,1.8469422) (0.89011792,1.8427244) (0.89535391,1.8384936) (0.90058989,1.8342499) (0.90582588,1.8299935) (0.91106187,1.8257244) (0.91629786,1.8214427) (0.92153385,1.8171485) (0.92676983,1.8128418) (0.93200582,1.8085228) (0.93724181,1.8041914) (0.9424778,1.7998479) (0.94771378,1.7954923) (0.95294977,1.7911245) (0.95818576,1.7867448) (0.96342175,1.7823532) (0.96865773,1.7779497) (0.97389372,1.7735344) (0.97912971,1.7691075) (0.9843657,1.7646688) (0.98960169,1.7602186) (0.99483767,1.7557569) (1.0000737,1.7512837) (1.0053096,1.7467991) (1.0105456,1.7423032) (1.0157816,1.737796) (1.0210176,1.7332776) (1.0262536,1.728748) (1.0314896,1.7242074) (1.0367256,1.7196556) (1.0419616,1.7150929) (1.0471976,1.7105192) (1.0524335,1.7059347) (1.0576695,1.7013393) (1.0629055,1.6967331) (1.0681415,1.6921161) (1.0733775,1.6874885) (1.0786135,1.6828502) (1.0838495,1.6782013) (1.0890855,1.6735418) (1.0943214,1.6688718) (1.0995574,1.6641913) (1.1047934,1.6595003) (1.1100294,1.654799) (1.1152654,1.6500873) (1.1205014,1.6453652) (1.1257374,1.6406329) (1.1309734,1.6358903) (1.1362093,1.6311375) (1.1414453,1.6263745) (1.1466813,1.6216013) (1.1519173,1.6168179) (1.1571533,1.6120245) (1.1623893,1.607221) (1.1676253,1.6024074) (1.1728613,1.5975838) (1.1780972,1.5927503) (1.1833332,1.5879067) (1.1885692,1.5830531) (1.1938052,1.5781897) (1.1990412,1.5733163) (1.2042772,1.568433) (1.2095132,1.5635398) (1.2147492,1.5586368) (1.2199851,1.5537239) (1.2252211,1.5488012) (1.2304571,1.5438687) (1.2356931,1.5389263) (1.2409291,1.5339742) (1.2461651,1.5290122) (1.2514011,1.5240405) (1.2566371,1.519059) (1.261873,1.5140677) (1.267109,1.5090667) (1.272345,1.5040558) (1.277581,1.4990353) (1.282817,1.494005) (1.288053,1.4889649) (1.293289,1.483915) (1.298525,1.4788554) (1.303761,1.473786) (1.3089969,1.4687069) (1.3142329,1.4636179) (1.3194689,1.4585192) (1.3247049,1.4534107) (1.3299409,1.4482923) (1.3351769,1.4431642) (1.3404129,1.4380262) (1.3456489,1.4328784) (1.3508848,1.4277207) (1.3561208,1.4225531) (1.3613568,1.4173756) (1.3665928,1.4121882) (1.3718288,1.4069908) (1.3770648,1.4017835) (1.3823008,1.3965662) (1.3875368,1.3913389) (1.3927727,1.3861015) (1.3980087,1.3808541) (1.4032447,1.3755966) (1.4084807,1.3703289) (1.4137167,1.365051) (1.4189527,1.359763) (1.4241887,1.3544647) (1.4294247,1.3491561) (1.4346606,1.3438372) (1.4398966,1.3385079) (1.4451326,1.3331682) (1.4503686,1.3278181) (1.4556046,1.3224575) (1.4608406,1.3170863) (1.4660766,1.3117045) (1.4713126,1.306312) (1.4765485,1.3009089) (1.4817845,1.295495) (1.4870205,1.2900702) (1.4922565,1.2846346) (1.4974925,1.279188) (1.5027285,1.2737303) (1.5079645,1.2682616) (1.5132005,1.2627818) (1.5184364,1.2572907) (1.5236724,1.2517883) (1.5289084,1.2462745) (1.5341444,1.2407493) (1.5393804,1.2352125) (1.5446164,1.2296641) (1.5498524,1.2241039) (1.5550884,1.218532) (1.5603244,1.2129482) (1.5655603,1.2073524) (1.5707963,1.2017445) (1.5760323,1.1961244) (1.5812683,1.190492) (1.5865043,1.1848472) (1.5917403,1.1791899) (1.5969763,1.17352) (1.6022123,1.1678374) (1.6074482,1.1621418) (1.6126842,1.1564333) (1.6179202,1.1507117) (1.6231562,1.1449768) (1.6283922,1.1392286) (1.6336282,1.1334669) (1.6388642,1.1276914) (1.6441002,1.1219022) (1.6493361,1.1160991) (1.6545721,1.1102818) (1.6598081,1.1044502) (1.6650441,1.0986043) (1.6702801,1.0927437) (1.6755161,1.0868683) (1.6807521,1.080978) (1.6859881,1.0750726) (1.691224,1.0691518) (1.69646,1.0632155) (1.701696,1.0572635) (1.706932,1.0512956) (1.712168,1.0453115) (1.717404,1.039311) (1.72264,1.033294) (1.727876,1.0272602) (1.7331119,1.0212093) (1.7383479,1.0151412) (1.7435839,1.0090554) (1.7488199,1.0029519) (1.7540559,0.99683031) (1.7592919,0.99069035) (1.7645279,0.98453175) (1.7697639,0.97835422) (1.7749998,0.97215746) (1.7802358,0.96594114) (1.7854718,0.95970496) (1.7907078,0.95344858) (1.7959438,0.94717165) (1.8011798,0.94087382) (1.8064158,0.93455473) (1.8116518,0.928214) (1.8168878,0.92185123) (1.8221237,0.91546603) (1.8273597,0.90905797) (1.8325957,0.90262664) (1.8378317,0.89617157) (1.8430677,0.88969232) (1.8483037,0.8831884) (1.8535397,0.87665932) (1.8587757,0.87010458) (1.8640116,0.86352363) (1.8692476,0.85691594) (1.8744836,0.85028093) (1.8797196,0.84361802) (1.8849556,0.83692659) (1.8901916,0.83020601) (1.8954276,0.82345562) (1.9006636,0.81667473) (1.9058995,0.80986264) (1.9111355,0.8030186) (1.9163715,0.79614185) (1.9216075,0.78923157) (1.9268435,0.78228695) (1.9320795,0.77530711) (1.9373155,0.76829114) (1.9425515,0.76123811) (1.9477874,0.75414701) (1.9530234,0.74701684) (1.9582594,0.7398465) (1.9634954,0.73263488) (1.9687314,0.72538079) (1.9739674,0.71808301) (1.9792034,0.71074024) (1.9844394,0.70335112) (1.9896753,0.69591425) (1.9949113,0.68842811) (2.0001473,0.68089115) (2.0053833,0.67330171) (2.0106193,0.66565805) (2.0158553,0.65795833) (2.0210913,0.65020063) (2.0263273,0.64238288) (2.0315632,0.63450293) (2.0367992,0.62655849) (2.0420352,0.61854713) (2.0472712,0.61046626) (2.0525072,0.60231316) (2.0577432,0.59408489) (2.0629792,0.58577835) (2.0682152,0.57739021) (2.0734512,0.56891693) (2.0786871,0.5603547) (2.0839231,0.55169943) (2.0891591,0.54294674) (2.0943951,0.53409189) (2.0996311,0.52512977) (2.1048671,0.51605486) (2.1101031,0.50686115) (2.1153391,0.49754213) (2.120575,0.48809067) (2.125811,0.47849898) (2.131047,0.46875853) (2.136283,0.4588599) (2.141519,0.44879268) (2.146755,0.43854534) (2.151991,0.42810501) (2.157227,0.41745726) (2.1624629,0.40658588) (2.1676989,0.39547249) (2.1729349,0.38409615) (2.1781709,0.37243284) (2.1834069,0.36045478) (2.1886429,0.34812956) (2.1938789,0.33541898) (2.1991149,0.32227754) (2.2043508,0.30865028) (2.2095868,0.29446986) (2.2148228,0.27965222) (2.2200588,0.26409022) (2.2252948,0.24764367) (2.2305308,0.23012311) (2.2357668,0.21126161) (2.2410028,0.19066173) (2.2462387,0.16768416) (2.2514747,0.14117281) (2.2567107,0.10856931) (2.2619467,0.060750487) (2.2619467,-0.060750487) (2.2567107,-0.10856931) (2.2514747,-0.14117281) (2.2462387,-0.16768416) (2.2410028,-0.19066173) (2.2357668,-0.21126161) (2.2305308,-0.23012311) (2.2252948,-0.24764367) (2.2200588,-0.26409022) (2.2148228,-0.27965222) (2.2095868,-0.29446986) (2.2043508,-0.30865028) (2.1991149,-0.32227754) (2.1938789,-0.33541898) (2.1886429,-0.34812956) (2.1834069,-0.36045478) (2.1781709,-0.37243284) (2.1729349,-0.38409615) (2.1676989,-0.39547249) (2.1624629,-0.40658588) (2.157227,-0.41745726) (2.151991,-0.42810501) (2.146755,-0.43854534) (2.141519,-0.44879268) (2.136283,-0.4588599) (2.131047,-0.46875853) (2.125811,-0.47849898) (2.120575,-0.48809067) (2.1153391,-0.49754213) (2.1101031,-0.50686115) (2.1048671,-0.51605486) (2.0996311,-0.52512977) (2.0943951,-0.53409189) (2.0891591,-0.54294674) (2.0839231,-0.55169943) (2.0786871,-0.5603547) (2.0734512,-0.56891693) (2.0682152,-0.57739021) (2.0629792,-0.58577835) (2.0577432,-0.59408489) (2.0525072,-0.60231316) (2.0472712,-0.61046626) (2.0420352,-0.61854713) (2.0367992,-0.62655849) (2.0315632,-0.63450293) (2.0263273,-0.64238288) (2.0210913,-0.65020063) (2.0158553,-0.65795833) (2.0106193,-0.66565805) (2.0053833,-0.67330171) (2.0001473,-0.68089115) (1.9949113,-0.68842811) (1.9896753,-0.69591425) (1.9844394,-0.70335112) (1.9792034,-0.71074024) (1.9739674,-0.71808301) (1.9687314,-0.72538079) (1.9634954,-0.73263488) (1.9582594,-0.7398465) (1.9530234,-0.74701684) (1.9477874,-0.75414701) (1.9425515,-0.76123811) (1.9373155,-0.76829114) (1.9320795,-0.77530711) (1.9268435,-0.78228695) (1.9216075,-0.78923157) (1.9163715,-0.79614185) (1.9111355,-0.8030186) (1.9058995,-0.80986264) (1.9006636,-0.81667473) (1.8954276,-0.82345562) (1.8901916,-0.83020601) (1.8849556,-0.83692659) (1.8797196,-0.84361802) (1.8744836,-0.85028093) (1.8692476,-0.85691594) (1.8640116,-0.86352363) (1.8587757,-0.87010458) (1.8535397,-0.87665932) (1.8483037,-0.8831884) (1.8430677,-0.88969232) (1.8378317,-0.89617157) (1.8325957,-0.90262664) (1.8273597,-0.90905797) (1.8221237,-0.91546603) (1.8168878,-0.92185123) (1.8116518,-0.928214) (1.8064158,-0.93455473) (1.8011798,-0.94087382) (1.7959438,-0.94717165) (1.7907078,-0.95344858) (1.7854718,-0.95970496) (1.7802358,-0.96594114) (1.7749998,-0.97215746) (1.7697639,-0.97835422) (1.7645279,-0.98453175) (1.7592919,-0.99069035) (1.7540559,-0.99683031) (1.7488199,-1.0029519) (1.7435839,-1.0090554) (1.7383479,-1.0151412) (1.7331119,-1.0212093) (1.727876,-1.0272602) (1.72264,-1.033294) (1.717404,-1.039311) (1.712168,-1.0453115) (1.706932,-1.0512956) (1.701696,-1.0572635) (1.69646,-1.0632155) (1.691224,-1.0691518) (1.6859881,-1.0750726) (1.6807521,-1.080978) (1.6755161,-1.0868683) (1.6702801,-1.0927437) (1.6650441,-1.0986043) (1.6598081,-1.1044502) (1.6545721,-1.1102818) (1.6493361,-1.1160991) (1.6441002,-1.1219022) (1.6388642,-1.1276914) (1.6336282,-1.1334669) (1.6283922,-1.1392286) (1.6231562,-1.1449768) (1.6179202,-1.1507117) (1.6126842,-1.1564333) (1.6074482,-1.1621418) (1.6022123,-1.1678374) (1.5969763,-1.17352) (1.5917403,-1.1791899) (1.5865043,-1.1848472) (1.5812683,-1.190492) (1.5760323,-1.1961244) (1.5707963,-1.2017445) (1.5655603,-1.2073524) (1.5603244,-1.2129482) (1.5550884,-1.218532) (1.5498524,-1.2241039) (1.5446164,-1.2296641) (1.5393804,-1.2352125) (1.5341444,-1.2407493) (1.5289084,-1.2462745) (1.5236724,-1.2517883) (1.5184364,-1.2572907) (1.5132005,-1.2627818) (1.5079645,-1.2682616) (1.5027285,-1.2737303) (1.4974925,-1.279188) (1.4922565,-1.2846346) (1.4870205,-1.2900702) (1.4817845,-1.295495) (1.4765485,-1.3009089) (1.4713126,-1.306312) (1.4660766,-1.3117045) (1.4608406,-1.3170863) (1.4556046,-1.3224575) (1.4503686,-1.3278181) (1.4451326,-1.3331682) (1.4398966,-1.3385079) (1.4346606,-1.3438372) (1.4294247,-1.3491561) (1.4241887,-1.3544647) (1.4189527,-1.359763) (1.4137167,-1.365051) (1.4084807,-1.3703289) (1.4032447,-1.3755966) (1.3980087,-1.3808541) (1.3927727,-1.3861015) (1.3875368,-1.3913389) (1.3823008,-1.3965662) (1.3770648,-1.4017835) (1.3718288,-1.4069908) (1.3665928,-1.4121882) (1.3613568,-1.4173756) (1.3561208,-1.4225531) (1.3508848,-1.4277207) (1.3456489,-1.4328784) (1.3404129,-1.4380262) (1.3351769,-1.4431642) (1.3299409,-1.4482923) (1.3247049,-1.4534107) (1.3194689,-1.4585192) (1.3142329,-1.4636179) (1.3089969,-1.4687069) (1.303761,-1.473786) (1.298525,-1.4788554) (1.293289,-1.483915) (1.288053,-1.4889649) (1.282817,-1.494005) (1.277581,-1.4990353) (1.272345,-1.5040558) (1.267109,-1.5090667) (1.261873,-1.5140677) (1.2566371,-1.519059) (1.2514011,-1.5240405) (1.2461651,-1.5290122) (1.2409291,-1.5339742) (1.2356931,-1.5389263) (1.2304571,-1.5438687) (1.2252211,-1.5488012) (1.2199851,-1.5537239) (1.2147492,-1.5586368) (1.2095132,-1.5635398) (1.2042772,-1.568433) (1.1990412,-1.5733163) (1.1938052,-1.5781897) (1.1885692,-1.5830531) (1.1833332,-1.5879067) (1.1780972,-1.5927503) (1.1728613,-1.5975838) (1.1676253,-1.6024074) (1.1623893,-1.607221) (1.1571533,-1.6120245) (1.1519173,-1.6168179) (1.1466813,-1.6216013) (1.1414453,-1.6263745) (1.1362093,-1.6311375) (1.1309734,-1.6358903) (1.1257374,-1.6406329) (1.1205014,-1.6453652) (1.1152654,-1.6500873) (1.1100294,-1.654799) (1.1047934,-1.6595003) (1.0995574,-1.6641913) (1.0943214,-1.6688718) (1.0890855,-1.6735418) (1.0838495,-1.6782013) (1.0786135,-1.6828502) (1.0733775,-1.6874885) (1.0681415,-1.6921161) (1.0629055,-1.6967331) (1.0576695,-1.7013393) (1.0524335,-1.7059347) (1.0471976,-1.7105192) (1.0419616,-1.7150929) (1.0367256,-1.7196556) (1.0314896,-1.7242074) (1.0262536,-1.728748) (1.0210176,-1.7332776) (1.0157816,-1.737796) (1.0105456,-1.7423032) (1.0053096,-1.7467991) (1.0000737,-1.7512837) (0.99483767,-1.7557569) (0.98960169,-1.7602186) (0.9843657,-1.7646688) (0.97912971,-1.7691075) (0.97389372,-1.7735344) (0.96865773,-1.7779497) (0.96342175,-1.7823532) (0.95818576,-1.7867448) (0.95294977,-1.7911245) (0.94771378,-1.7954923) (0.9424778,-1.7998479) (0.93724181,-1.8041914) (0.93200582,-1.8085228) (0.92676983,-1.8128418) (0.92153385,-1.8171485) (0.91629786,-1.8214427) (0.91106187,-1.8257244) (0.90582588,-1.8299935) (0.90058989,-1.8342499) (0.89535391,-1.8384936) (0.89011792,-1.8427244) (0.88488193,-1.8469422) (0.87964594,-1.8511471) (0.87440996,-1.8553388) (0.86917397,-1.8595174) (0.86393798,-1.8636826) (0.85870199,-1.8678345) (0.853466,-1.8719729) (0.84823002,-1.8760977) (0.84299403,-1.8802088) (0.83775804,-1.8843062) (0.83252205,-1.8883898) (0.82728607,-1.8924593) (0.82205008,-1.8965148) (0.81681409,-1.9005562) (0.8115781,-1.9045833) (0.80634211,-1.908596) (0.80110613,-1.9125942) (0.79587014,-1.9165779) (0.79063415,-1.9205468) (0.78539816,-1.924501) (0.78016218,-1.9284402) (0.77492619,-1.9323645) (0.7696902,-1.9362735) (0.76445421,-1.9401674) (0.75921822,-1.9440458) (0.75398224,-1.9479088) (0.74874625,-1.9517561) (0.74351026,-1.9555878) (0.73827427,-1.9594035) (0.73303829,-1.9632033) (0.7278023,-1.966987) (0.72256631,-1.9707545) (0.71733032,-1.9745056) (0.71209433,-1.9782402) (0.70685835,-1.9819583) (0.70162236,-1.9856596) (0.69638637,-1.989344) (0.69115038,-1.9930114) (0.6859144,-1.9966616) (0.68067841,-2.0002946) (0.67544242,-2.0039102) (0.67020643,-2.0075082) (0.66497045,-2.0110885) (0.65973446,-2.0146509) (0.65449847,-2.0181954) (0.64926248,-2.0217218) (0.64402649,-2.0252298) (0.63879051,-2.0287195) (0.63355452,-2.0321906) (0.62831853,-2.035643) (0.62308254,-2.0390765) (0.61784656,-2.042491) (0.61261057,-2.0458863) (0.60737458,-2.0492623) (0.60213859,-2.0526188) (0.5969026,-2.0559557) (0.59166662,-2.0592727) (0.58643063,-2.0625698) (0.58119464,-2.0658468) (0.57595865,-2.0691036) (0.57072267,-2.0723398) (0.56548668,-2.0755555) (0.56025069,-2.0787504) (0.5550147,-2.0819244) (0.54977871,-2.0850772) (0.54454273,-2.0882088) (0.53930674,-2.091319) (0.53407075,-2.0944075) (0.52883476,-2.0974743) (0.52359878,-2.1005191) (0.51836279,-2.1035418) (0.5131268,-2.1065422) (0.50789081,-2.1095201) (0.50265482,-2.1124754) (0.49741884,-2.1154078) (0.49218285,-2.1183173) (0.48694686,-2.1212036) (0.48171087,-2.1240665) (0.47647489,-2.1269059) (0.4712389,-2.1297215) (0.46600291,-2.1325133) (0.46076692,-2.135281) (0.45553093,-2.1380244) (0.45029495,-2.1407434) (0.44505896,-2.1434378) (0.43982297,-2.1461074) (0.43458698,-2.1487519) (0.429351,-2.1513714) (0.42411501,-2.1539654) (0.41887902,-2.1565339) (0.41364303,-2.1590767) (0.40840704,-2.1615936) (0.40317106,-2.1640844) (0.39793507,-2.166549) (0.39269908,-2.168987) (0.38746309,-2.1713984) (0.38222711,-2.173783) (0.37699112,-2.1761406) (0.37175513,-2.178471) (0.36651914,-2.180774) (0.36128316,-2.1830494) (0.35604717,-2.1852971) (0.35081118,-2.1875169) (0.34557519,-2.1897086) (0.3403392,-2.1918719) (0.33510322,-2.1940068) (0.32986723,-2.196113) (0.32463124,-2.1981905) (0.31939525,-2.2002389) (0.31415927,-2.2022581) (0.30892328,-2.2042479) (0.30368729,-2.2062082) (0.2984513,-2.2081388) (0.29321531,-2.2100395) (0.28797933,-2.2119102) (0.28274334,-2.2137506) (0.27750735,-2.2155606) (0.27227136,-2.21734) (0.26703538,-2.2190888) (0.26179939,-2.2208066) (0.2565634,-2.2224933) (0.25132741,-2.2241489) (0.24609142,-2.225773) (0.24085544,-2.2273656) (0.23561945,-2.2289265) (0.23038346,-2.2304556) (0.22514747,-2.2319526) (0.21991149,-2.2334176) (0.2146755,-2.2348502) (0.20943951,-2.2362503) (0.20420352,-2.2376179) (0.19896753,-2.2389528) (0.19373155,-2.2402548) (0.18849556,-2.2415237) (0.18325957,-2.2427596) (0.17802358,-2.2439622) (0.1727876,-2.2451314) (0.16755161,-2.2462671) (0.16231562,-2.2473691) (0.15707963,-2.2484374) (0.15184364,-2.2494719) (0.14660766,-2.2504723) (0.14137167,-2.2514387) (0.13613568,-2.2523709) (0.13089969,-2.2532688) (0.12566371,-2.2541322) (0.12042772,-2.2549612) (0.11519173,-2.2557556) (0.10995574,-2.2565154) (0.10471976,-2.2572403) (0.099483767,-2.2579305) (0.09424778,-2.2585857) (0.089011792,-2.2592059) (0.083775804,-2.2597911) (0.078539816,-2.2603412) (0.073303829,-2.260856) (0.068067841,-2.2613356) (0.062831853,-2.2617799) (0.057595865,-2.2621889) (0.052359878,-2.2625624) (0.04712389,-2.2629005) (0.041887902,-2.2632031) (0.036651914,-2.2634702) (0.031415927,-2.2637018) (0.026179939,-2.2638977) (0.020943951,-2.2640581) (0.015707963,-2.2641828) (0.010471976,-2.264272) (0.0052359878,-2.2643254) (-4.4408921e-16,-2.2643433) (-0.0052359878,-2.2643254) (-0.010471976,-2.264272) (-0.015707963,-2.2641828) (-0.020943951,-2.2640581) (-0.026179939,-2.2638977) (-0.031415927,-2.2637018) (-0.036651914,-2.2634702) (-0.041887902,-2.2632031) (-0.04712389,-2.2629005) (-0.052359878,-2.2625624) (-0.057595865,-2.2621889) (-0.062831853,-2.2617799) (-0.068067841,-2.2613356) (-0.073303829,-2.260856) (-0.078539816,-2.2603412) (-0.083775804,-2.2597911) (-0.089011792,-2.2592059) (-0.09424778,-2.2585857) (-0.099483767,-2.2579305) (-0.10471976,-2.2572403) (-0.10995574,-2.2565154) (-0.11519173,-2.2557556) (-0.12042772,-2.2549612) (-0.12566371,-2.2541322) (-0.13089969,-2.2532688) (-0.13613568,-2.2523709) (-0.14137167,-2.2514387) (-0.14660766,-2.2504723) (-0.15184364,-2.2494719) (-0.15707963,-2.2484374) (-0.16231562,-2.2473691) (-0.16755161,-2.2462671) (-0.1727876,-2.2451314) (-0.17802358,-2.2439622) (-0.18325957,-2.2427596) (-0.18849556,-2.2415237) (-0.19373155,-2.2402548) (-0.19896753,-2.2389528) (-0.20420352,-2.2376179) (-0.20943951,-2.2362503) (-0.2146755,-2.2348502) (-0.21991149,-2.2334176) (-0.22514747,-2.2319526) (-0.23038346,-2.2304556) (-0.23561945,-2.2289265) (-0.24085544,-2.2273656) (-0.24609142,-2.225773) (-0.25132741,-2.2241489) (-0.2565634,-2.2224933) (-0.26179939,-2.2208066) (-0.26703538,-2.2190888) (-0.27227136,-2.21734) (-0.27750735,-2.2155606) (-0.28274334,-2.2137506) (-0.28797933,-2.2119102) (-0.29321531,-2.2100395) (-0.2984513,-2.2081388) (-0.30368729,-2.2062082) (-0.30892328,-2.2042479) (-0.31415927,-2.2022581) (-0.31939525,-2.2002389) (-0.32463124,-2.1981905) (-0.32986723,-2.196113) (-0.33510322,-2.1940068) (-0.3403392,-2.1918719) (-0.34557519,-2.1897086) (-0.35081118,-2.1875169) (-0.35604717,-2.1852971) (-0.36128316,-2.1830494) (-0.36651914,-2.180774) (-0.37175513,-2.178471) (-0.37699112,-2.1761406) (-0.38222711,-2.173783) (-0.38746309,-2.1713984) (-0.39269908,-2.168987) (-0.39793507,-2.166549) (-0.40317106,-2.1640844) (-0.40840704,-2.1615936) (-0.41364303,-2.1590767) (-0.41887902,-2.1565339) (-0.42411501,-2.1539654) (-0.429351,-2.1513714) (-0.43458698,-2.1487519) (-0.43982297,-2.1461074) (-0.44505896,-2.1434378) (-0.45029495,-2.1407434) (-0.45553093,-2.1380244) (-0.46076692,-2.135281) (-0.46600291,-2.1325133) (-0.4712389,-2.1297215) (-0.47647489,-2.1269059) (-0.48171087,-2.1240665) (-0.48694686,-2.1212036) (-0.49218285,-2.1183173) (-0.49741884,-2.1154078) (-0.50265482,-2.1124754) (-0.50789081,-2.1095201) (-0.5131268,-2.1065422) (-0.51836279,-2.1035418) (-0.52359878,-2.1005191) (-0.52883476,-2.0974743) (-0.53407075,-2.0944075) (-0.53930674,-2.091319) (-0.54454273,-2.0882088) (-0.54977871,-2.0850772) (-0.5550147,-2.0819244) (-0.56025069,-2.0787504) (-0.56548668,-2.0755555) (-0.57072267,-2.0723398) (-0.57595865,-2.0691036) (-0.58119464,-2.0658468) (-0.58643063,-2.0625698) (-0.59166662,-2.0592727) (-0.5969026,-2.0559557) (-0.60213859,-2.0526188) (-0.60737458,-2.0492623) (-0.61261057,-2.0458863) (-0.61784656,-2.042491) (-0.62308254,-2.0390765) (-0.62831853,-2.035643) (-0.63355452,-2.0321906) (-0.63879051,-2.0287195) (-0.64402649,-2.0252298) (-0.64926248,-2.0217218) (-0.65449847,-2.0181954) (-0.65973446,-2.0146509) (-0.66497045,-2.0110885) (-0.67020643,-2.0075082) (-0.67544242,-2.0039102) (-0.68067841,-2.0002946) (-0.6859144,-1.9966616) (-0.69115038,-1.9930114) (-0.69638637,-1.989344) (-0.70162236,-1.9856596) (-0.70685835,-1.9819583) (-0.71209433,-1.9782402) (-0.71733032,-1.9745056) (-0.72256631,-1.9707545) (-0.7278023,-1.966987) (-0.73303829,-1.9632033) (-0.73827427,-1.9594035) (-0.74351026,-1.9555878) (-0.74874625,-1.9517561) (-0.75398224,-1.9479088) (-0.75921822,-1.9440458) (-0.76445421,-1.9401674) (-0.7696902,-1.9362735) (-0.77492619,-1.9323645) (-0.78016218,-1.9284402) (-0.78539816,-1.924501) (-0.79063415,-1.9205468) (-0.79587014,-1.9165779) (-0.80110613,-1.9125942) (-0.80634211,-1.908596) (-0.8115781,-1.9045833) (-0.81681409,-1.9005562) (-0.82205008,-1.8965148) (-0.82728607,-1.8924593) (-0.83252205,-1.8883898) (-0.83775804,-1.8843062) (-0.84299403,-1.8802088) (-0.84823002,-1.8760977) (-0.853466,-1.8719729) (-0.85870199,-1.8678345) (-0.86393798,-1.8636826) (-0.86917397,-1.8595174) (-0.87440996,-1.8553388) (-0.87964594,-1.8511471) (-0.88488193,-1.8469422) (-0.89011792,-1.8427244) (-0.89535391,-1.8384936) (-0.90058989,-1.8342499) (-0.90582588,-1.8299935) (-0.91106187,-1.8257244) (-0.91629786,-1.8214427) (-0.92153385,-1.8171485) (-0.92676983,-1.8128418) (-0.93200582,-1.8085228) (-0.93724181,-1.8041914) (-0.9424778,-1.7998479) (-0.94771378,-1.7954923) (-0.95294977,-1.7911245) (-0.95818576,-1.7867448) (-0.96342175,-1.7823532) (-0.96865773,-1.7779497) (-0.97389372,-1.7735344) (-0.97912971,-1.7691075) (-0.9843657,-1.7646688) (-0.98960169,-1.7602186) (-0.99483767,-1.7557569) (-1.0000737,-1.7512837) (-1.0053096,-1.7467991) (-1.0105456,-1.7423032) (-1.0157816,-1.737796) (-1.0210176,-1.7332776) (-1.0262536,-1.728748) (-1.0314896,-1.7242074) (-1.0367256,-1.7196556) (-1.0419616,-1.7150929) (-1.0471976,-1.7105192) (-1.0524335,-1.7059347) (-1.0576695,-1.7013393) (-1.0629055,-1.6967331) (-1.0681415,-1.6921161) (-1.0733775,-1.6874885) (-1.0786135,-1.6828502) (-1.0838495,-1.6782013) (-1.0890855,-1.6735418) (-1.0943214,-1.6688718) (-1.0995574,-1.6641913) (-1.1047934,-1.6595003) (-1.1100294,-1.654799) (-1.1152654,-1.6500873) (-1.1205014,-1.6453652) (-1.1257374,-1.6406329) (-1.1309734,-1.6358903) (-1.1362093,-1.6311375) (-1.1414453,-1.6263745) (-1.1466813,-1.6216013) (-1.1519173,-1.6168179) (-1.1571533,-1.6120245) (-1.1623893,-1.607221) (-1.1676253,-1.6024074) (-1.1728613,-1.5975838) (-1.1780972,-1.5927503) (-1.1833332,-1.5879067) (-1.1885692,-1.5830531) (-1.1938052,-1.5781897) (-1.1990412,-1.5733163) (-1.2042772,-1.568433) (-1.2095132,-1.5635398) (-1.2147492,-1.5586368) (-1.2199851,-1.5537239) (-1.2252211,-1.5488012) (-1.2304571,-1.5438687) (-1.2356931,-1.5389263) (-1.2409291,-1.5339742) (-1.2461651,-1.5290122) (-1.2514011,-1.5240405) (-1.2566371,-1.519059) (-1.261873,-1.5140677) (-1.267109,-1.5090667) (-1.272345,-1.5040558) (-1.277581,-1.4990353) (-1.282817,-1.494005) (-1.288053,-1.4889649) (-1.293289,-1.483915) (-1.298525,-1.4788554) (-1.303761,-1.473786) (-1.3089969,-1.4687069) (-1.3142329,-1.4636179) (-1.3194689,-1.4585192) (-1.3247049,-1.4534107) (-1.3299409,-1.4482923) (-1.3351769,-1.4431642) (-1.3404129,-1.4380262) (-1.3456489,-1.4328784) (-1.3508848,-1.4277207) (-1.3561208,-1.4225531) (-1.3613568,-1.4173756) (-1.3665928,-1.4121882) (-1.3718288,-1.4069908) (-1.3770648,-1.4017835) (-1.3823008,-1.3965662) (-1.3875368,-1.3913389) (-1.3927727,-1.3861015) (-1.3980087,-1.3808541) (-1.4032447,-1.3755966) (-1.4084807,-1.3703289) (-1.4137167,-1.365051) (-1.4189527,-1.359763) (-1.4241887,-1.3544647) (-1.4294247,-1.3491561) (-1.4346606,-1.3438372) (-1.4398966,-1.3385079) (-1.4451326,-1.3331682) (-1.4503686,-1.3278181) (-1.4556046,-1.3224575) (-1.4608406,-1.3170863) (-1.4660766,-1.3117045) (-1.4713126,-1.306312) (-1.4765485,-1.3009089) (-1.4817845,-1.295495) (-1.4870205,-1.2900702) (-1.4922565,-1.2846346) (-1.4974925,-1.279188) (-1.5027285,-1.2737303) (-1.5079645,-1.2682616) (-1.5132005,-1.2627818) (-1.5184364,-1.2572907) (-1.5236724,-1.2517883) (-1.5289084,-1.2462745) (-1.5341444,-1.2407493) (-1.5393804,-1.2352125) (-1.5446164,-1.2296641) (-1.5498524,-1.2241039) (-1.5550884,-1.218532) (-1.5603244,-1.2129482) (-1.5655603,-1.2073524) (-1.5707963,-1.2017445) (-1.5760323,-1.1961244) (-1.5812683,-1.190492) (-1.5865043,-1.1848472) (-1.5917403,-1.1791899) (-1.5969763,-1.17352) (-1.6022123,-1.1678374) (-1.6074482,-1.1621418) (-1.6126842,-1.1564333) (-1.6179202,-1.1507117) (-1.6231562,-1.1449768) (-1.6283922,-1.1392286) (-1.6336282,-1.1334669) (-1.6388642,-1.1276914) (-1.6441002,-1.1219022) (-1.6493361,-1.1160991) (-1.6545721,-1.1102818) (-1.6598081,-1.1044502) (-1.6650441,-1.0986043) (-1.6702801,-1.0927437) (-1.6755161,-1.0868683) (-1.6807521,-1.080978) (-1.6859881,-1.0750726) (-1.691224,-1.0691518) (-1.69646,-1.0632155) (-1.701696,-1.0572635) (-1.706932,-1.0512956) (-1.712168,-1.0453115) (-1.717404,-1.039311) (-1.72264,-1.033294) (-1.727876,-1.0272602) (-1.7331119,-1.0212093) (-1.7383479,-1.0151412) (-1.7435839,-1.0090554) (-1.7488199,-1.0029519) (-1.7540559,-0.99683031) (-1.7592919,-0.99069035) (-1.7645279,-0.98453175) (-1.7697639,-0.97835422) (-1.7749998,-0.97215746) (-1.7802358,-0.96594114) (-1.7854718,-0.95970496) (-1.7907078,-0.95344858) (-1.7959438,-0.94717165) (-1.8011798,-0.94087382) (-1.8064158,-0.93455473) (-1.8116518,-0.928214) (-1.8168878,-0.92185123) (-1.8221237,-0.91546603) (-1.8273597,-0.90905797) (-1.8325957,-0.90262664) (-1.8378317,-0.89617157) (-1.8430677,-0.88969232) (-1.8483037,-0.8831884) (-1.8535397,-0.87665932) (-1.8587757,-0.87010458) (-1.8640116,-0.86352363) (-1.8692476,-0.85691594) (-1.8744836,-0.85028093) (-1.8797196,-0.84361802) (-1.8849556,-0.83692659) (-1.8901916,-0.83020601) (-1.8954276,-0.82345562) (-1.9006636,-0.81667473) (-1.9058995,-0.80986264) (-1.9111355,-0.8030186) (-1.9163715,-0.79614185) (-1.9216075,-0.78923157) (-1.9268435,-0.78228695) (-1.9320795,-0.77530711) (-1.9373155,-0.76829114) (-1.9425515,-0.76123811) (-1.9477874,-0.75414701) (-1.9530234,-0.74701684) (-1.9582594,-0.7398465) (-1.9634954,-0.73263488) (-1.9687314,-0.72538079) (-1.9739674,-0.71808301) (-1.9792034,-0.71074024) (-1.9844394,-0.70335112) (-1.9896753,-0.69591425) (-1.9949113,-0.68842811) (-2.0001473,-0.68089115) (-2.0053833,-0.67330171) (-2.0106193,-0.66565805) (-2.0158553,-0.65795833) (-2.0210913,-0.65020063) (-2.0263273,-0.64238288) (-2.0315632,-0.63450293) (-2.0367992,-0.62655849) (-2.0420352,-0.61854713) (-2.0472712,-0.61046626) (-2.0525072,-0.60231316) (-2.0577432,-0.59408489) (-2.0629792,-0.58577835) (-2.0682152,-0.57739021) (-2.0734512,-0.56891693) (-2.0786871,-0.5603547) (-2.0839231,-0.55169943) (-2.0891591,-0.54294674) (-2.0943951,-0.53409189) (-2.0996311,-0.52512977) (-2.1048671,-0.51605486) (-2.1101031,-0.50686115) (-2.1153391,-0.49754213) (-2.120575,-0.48809067) (-2.125811,-0.47849898) (-2.131047,-0.46875853) (-2.136283,-0.4588599) (-2.141519,-0.44879268) (-2.146755,-0.43854534) (-2.151991,-0.42810501) (-2.157227,-0.41745726) (-2.1624629,-0.40658588) (-2.1676989,-0.39547249) (-2.1729349,-0.38409615) (-2.1781709,-0.37243284) (-2.1834069,-0.36045478) (-2.1886429,-0.34812956) (-2.1938789,-0.33541898) (-2.1991149,-0.32227754) (-2.2043508,-0.30865028) (-2.2095868,-0.29446986) (-2.2148228,-0.27965222) (-2.2200588,-0.26409022) (-2.2252948,-0.24764367) (-2.2305308,-0.23012311) (-2.2357668,-0.21126161) (-2.2410028,-0.19066173) (-2.2462387,-0.16768416) (-2.2514747,-0.14117281) (-2.2567107,-0.10856931) (-2.2619467,-0.060750487)};
\addlegendentry{$p=0.3$}
\end{axis}\end{tikzpicture}\end{center}

